%% Copernicus Publications Manuscript Preparation Template for LaTeX Submissions
%% ---------------------------------
%% This template should be used for copernicus.cls
%% The class file and some style files are bundled in the Copernicus Latex Package, which can be downloaded from the different journal webpages.
%% For further assistance please contact Copernicus Publications at: production@copernicus.org
%% https://publications.copernicus.org/for_authors/manuscript_preparation.html


%% Please use the following documentclass and journal abbreviations for preprints and final revised papers.

%% 2-column papers and preprints
\documentclass[amt, manuscript]{copernicus}



%% Journal abbreviations (please use the same for preprints and final revised papers)


% Advances in Geosciences (adgeo)
% Advances in Radio Science (ars)
% Advances in Science and Research (asr)
% Advances in Statistical Climatology, Meteorology and Oceanography (ascmo)
% Aerosol Research (ar)
% Annales Geophysicae (angeo)
% Archives Animal Breeding (aab)
% Atmospheric Chemistry and Physics (acp)
% Atmospheric Measurement Techniques (amt)
% Biogeosciences (bg)
% Climate of the Past (cp)
% DEUQUA Special Publications (deuquasp)
% Earth Surface Dynamics (esurf)
% Earth System Dynamics (esd)
% Earth System Science Data (essd)
% E&G Quaternary Science Journal (egqsj)
% EGUsphere (egusphere) | This is only for EGUsphere preprints submitted without relation to an EGU journal.
% European Journal of Mineralogy (ejm)
% Fossil Record (fr)
% Geochronology (gchron)
% Geographica Helvetica (gh)
% Geoscience Communication (gc)
% Geoscientific Instrumentation, Methods and Data Systems (gi)
% Geoscientific Model Development (gmd)
% History of Geo- and Space Sciences (hgss)
% Hydrology and Earth System Sciences (hess)
% Journal of Bone and Joint Infection (jbji)
% Journal of Micropalaeontology (jm)
% Journal of Sensors and Sensor Systems (jsss)
% Magnetic Resonance (mr)
% Mechanical Sciences (ms)
% Natural Hazards and Earth System Sciences (nhess)
% Nonlinear Processes in Geophysics (npg)
% Ocean Science (os)
% Polarforschung - Journal of the German Society for Polar Research (polf)
% Primate Biology (pb)
% Proceedings of the International Association of Hydrological Sciences (piahs)
% Safety of Nuclear Waste Disposal (sand)
% Scientific Drilling (sd)
% SOIL (soil)
% Solid Earth (se)
% State of the Planet (sp)
% The Cryosphere (tc)
% Weather and Climate Dynamics (wcd)
% Web Ecology (we)
% Wind Energy Science (wes)


%% \usepackage commands included in the copernicus.cls:
%\usepackage[german, english]{babel}
%\usepackage{tabularx}
%\usepackage{cancel}
%\usepackage{multirow}
%\usepackage{supertabular}
%\usepackage{algorithmic}
%\usepackage{algorithm}
%\usepackage{amsthm}
%\usepackage{float}
%\usepackage{subfig}
%\usepackage{rotating}

\graphicspath{{Figures/}}

\usepackage{placeins}
\usepackage{longtable}
\usepackage{tabulary}
\usepackage{booktabs,array,multirow}
\usepackage{makecell}

%\usepackage[capitalize,noabbrev]{cleveref}


\begin{document}

\title{Universality of ice particle shapes, densities, and orientations as revealed by the 2D-S in situ optical array probe}


% \Author[affil]{given_name}{surname}

\Author[1,2][edwin.dunnavan@noaa.gov]{Edwin L.}{Dunnavan} %% correspondence author
\Author[3,4]{Joseph A.}{Finlon}
\Author[5]{David J.}{Delene}
\Author[5]{Christian}{Nairy}
\Author[6]{Aaron}{Bansemer}
\Author[7]{Craig}{Hunter}
%\Author[1,5]{Greg M.}{McFarquhar}

\affil[1]{Cooperative Institute for Severe and High-Impact Weather Research and Operations (CIWRO), Norman, Oklahoma}
\affil[2]{NOAA/OAR National Severe Storms Laboratory, Norman, Oklahoma}
\affil[3]{Earth System Science Interdisciplinary Center, University of Maryland, College Park, Maryland, USA}
\affil[4]{NASA Goddard Space Flight Center, Greenbelt, Maryland, USA}
%\affil[5]{School of Meteorology, University of Oklahoma, Norman, OK, USA}
\affil[5]{Department of Atmospheric Sciences, University of North Dakota, Grand Forks, North Dakota, USA}
\affil[6]{NSF National Center for Atmospheric Research, Boulder, Colorado, USA}
\affil[7]{NASA Langley Research Center, Hampton, VA, USA}

%% The [] brackets identify the author with the corresponding affiliation. 1, 2, 3, etc. should be inserted.

%% If an author is deceased, please add \deceased[$Deceased date if applicable$]{$Author number$} (e.g. \deceased[13 November 2015]{2}) at the end of the affiliations. The author number depends on the placement of the author in the author list, e.g. the third author has number 3.


%% If authors contributed equally, please add \equalcontrib{$Author numbers$} (e.g. \equalcontrib{1,3}) at the end of the affiliations. The author number depends on the placement of the author in the author list, e.g. the third author has number 3.


\runningtitle{Universality of ice particle shapes and orientations as revealed by the 2D-S in situ optical array probe}

\runningauthor{Edwin L. Dunnavan}





\received{}
\pubdiscuss{} %% only important for two-stage journals
\revised{}
\accepted{}
\published{}

%% These dates will be inserted by Copernicus Publications during the typesetting process.


\firstpage{1}

\maketitle


%Joint distributions of observed horizontally- and vertically-viewed ellipse aspect ratios, orientations, sizes, and area ratios suggest a universality of 3D ice particle shapes, orientations, and densities for irregular and columnar particle habits.

%While the 2D projection metrics are nearly identical between H and V channels, 2D orientation discrepancies between channels and flights prove that most particle orientations within in situ volumes are neither random nor horizontally aligned and depend crucially upon habit and the precise 2D-S configuration itself.

%Monte Carlo simulations using theoretical triaxial ellipsoids of multivariate shapes, densities, and orientations and simulations using explicit 3D STL model geometries can match ellipse parameter joint distributions for irregular and columnar crystals.

\begin{abstract}
Previous studies have used in situ aircraft optical array probe measurements to infer atmospheric ice particle shapes and densities from their 2D projections despite highly uncertain probe-induced 3D orientations. This study reconciles many different ice particle 3D shapes, densities, and in situ orientations from 2D orthogonal-view images by comparing explicit ellipsoid and STL simulations to nearly 14 million 2D-S ellipse-fitted particle projections collected from 26 flights of the NASA IMPACTS field campaign. Irregular particles below $-18^{\circ}$C from each flight demonstrate a robust universality of ellipse parameters with median distributional deviations of approximately 3\% for aspect ratios and orientations and 6\% for area ratio. Monte Carlo simulations of triaxial ellipsoids with the same beta-distributed shapes and densities achieve more than $90\%$ similarity to each observed irregular ellipse joint distribution from eight individual flights. Ellipsoid and STL simulations of columns and needles elucidate the preferred orientations from different campaign years and slight changes in the 2D-S aircraft placement can change preferential canting by $20^{\circ}$. A suite of habit-specific STL simulations suggests that, when the 2D-S is properly oriented, correction factors needed to convert 2D projection properties to 3D counterparts are nearly identical in both horizontal and vertical channels where the maximum median difference among all tests is only 1.1\% and where full cumulative distributions deviate by at most 5.08\%. These results indicate that knowledge of these characteristic 2D ellipse fit joint distributions can be used to probabilistically supplement 2D-S ice particle retrievals of 3D morphology and habit identification algorithms.
\end{abstract}

%Monte Carlo simulations using theoretical  triaxial ellipsoids and explicit 3D STL model geometries are used to propose plausible shapes, densities, and orientations needed to match ellipse parameter joint distributions for irregular and columnar crystals.

%The observed 2D-S ellipse distribution symmetry between the horizontal and vertical channels for irregular particles, in particular, can be explained by representing irregular particles as triaxial ellipsoids with size-independent aspect ratios and densities rather than the power-law spheroidal representations commonly used in numerical weather prediction models.

%Simulations suggest that particle maximum dimensions are preferentially oriented by approximately $25^{\circ}$ from horizontal. Columns from the 2D-S ellipse fits can be visually distinguished by their power-law like behavior with increasing size along with their $70^{\circ}$ preferential orientation as viewed from the 2D-S vertical channel. 

%\copyrightstatement{TEXT} %% This section is optional and can be used for copyright transfers.

%Ellipsoid simulations representing irregular crystals for eight separate flights achieve more than 90\% distributional similarity with each observed joint ellipse parameter set.


\introduction  %% \introduction[modified heading if necessary]
Ice particle shapes, densities, and orientations are still poorly understood and represented in both local and global numerical weather prediction (NWP) models. These deficiencies result from the vast landscape of various ice particle habit forms and an inability or difficulty in estimating 3D structures from the 2D measurements afforded by available imagery in different environments. There has however been recent progress in identifying and characterizing ice particles. Recent studies have applied machine learning techniques such as convolutional neural networks \citep{Xiao2019,Przybylo2022b,Jaffeux2022,Zhang2023} and multinomial logistic regression \citep{Praz2018} to identify ice particle habits from in situ derived cloud particle imagers. Many of these studies claim up to $97\%$ accuracy in correctly identifying ice particle habits. There are also recent machine learning studies that have made advancements in characterizing snow morphological and microphysical properties from ground-based imagers \citep{Leinonen2021}. 

Despite this recent progress in habit \textit{classification}, there are still large uncertainties regarding the actual \textit{quantification} of certain habit type shapes and densities and current NWP models still use overly simplistic geometric assumptions. In particular, irregular ice particles -- unaggregated pristine crystals without a clear, singular hexagonal structure\footnote{There is some debate regarding the use and appropriateness of the word ``irregular'' in describing such particles. For example, \citet{Stoelinga2007} decries the typical usage of the word ``irregular'' as it suggests that such forms aren't already classified in habit diagrams when such ``regular'' forms are quite well known. Further, \citet{Stoelinga2007} correctly states that probe uncertainties and limited resolutions can frequently lead to misclassifications of such ``irregular'' crystals, thus contributing to confusion regarding what ``irregular'' is supposed to represent. Herein we use the term irregular to simply mean unaggregated, non-pristine crystals; that is, ice crystals that are not plates, columns, dendrites, or needles. Such a definition is consistent with many habit classification conventions and is consistent with the studies of \citet{Korolev2000}, \citet{Korolev2003}, and others.} -- by far exhibit the largest variations and uncertainties in morphology and microphysical properties. Irregular particles exhibit many different exotic habit forms such as rosettes, bullet rosettes, side planes, and many more \citep{Bailey2009} which appear at nearly all subfreezing temperature levels. Each form potentially encapsulates different microphysical and radiative properties even within the same subclass. Thus, it is difficult, if not impossible, for models to utilize theoretical or analytical methods to calculate radiative properties, growth rates, or other microphysical properties for all types of irregular particles even within the same cloud. Instead, approximate or purely empirical methods must be employed. These methods are particularly important when considering that the vast majority of stratiform ice particles are of these irregular forms and are not pristine such as plates, columns, dendrites, or needles \citep{Korolev1999,Korolev2000,Korolev2003}.
In situ aircraft studies are necessary to understand the shapes of such irregular particles since snow particles will almost always be aggregated when such particles are observed from ground instruments such as the Multi-Angle Snowflake Camera \citep[MASC;][]{Garrett2015}.

%% IMPORTANT NOTE
% From abstract of Macke et al. 1998
% "From a broadband solar radiative transfer
%point of view it appears more important to use the correct particle shapes than to average over the correct size
%distributions."

There are currently two paradigms that attempt to bridge the gap between the observed variety of ice particles and the calculations of their properties. The first paradigm involves the detailed, explicit numerical calculations or laboratory experiments of particle properties using synthetic 3D structures. These 3D structures are often completely idealistic and represent perfect hexagonal (e.g., plate and column) or otherwise simple shapes. Such an approach is not only used in numerical experiments \citep[e.g.,][]{Um2009,Schrom2018b,Schrom2019} and in the development of particle ``databases'' \citep[e.g.,][]{Liu2008,Lu2016,Schrom2023}  but also in laboratory studies of, for example, particle fallspeed \citep[e.g.,][]{Westbrook2017,McCorquodale2020a,McCorquodale2020b}. Performing explicit calculations with such imitations provides confidence in the results of such calculations due to the exact nature of the problem, especially for the calculation of radiative properties as these quantities are physically based. In the case of scattering and absorption calculations, it is common for models to use lookup tables generated from these explicit calculations and then integrate the properties across a size distribution. While this paradigm is explicit and physical, its tradeoff lies in the uncertainty regarding the universality of representing a myriad of particles with only a limited number of available synthetic or idealistic representations. Despite this variability and uncertainty, several studies have found that idealized shapes such as Koch fractals \citep{Macke1996,McFarquhar1999} can be used to calculate the mean radiative properties of ice crystals.

The second paradigm potentially overcomes the first paradigm's limitation by replacing explicit single-particle calculations with integrated properties stemming from an assumed analytical geometry and underlying distributions governing such geometry. Thus, fundamental radiative and microphysical properties can be calculated through the integration of each particle's geometry along with their particle size distribution. Oblate and prolate spheroids are frequently used to broadly represent the shapes of various ice particles. Geometric quantities, such as the capacitance needed for growth rate calculations \citep{Pruppacher1997,Lambbook} and radiative quantities such as scattering amplitudes \citep{Ryzhkov2011} can therefore be determined analytically and only intrinsically depend upon the aspect ratios and sizes of the spheroids. An assumed spheroid ``effective'' density is then used to complete the geometry of any particular particle by specifying the mass of each particle. While the second paradigm is widely used, there is often disagreement as to what spheroid aspect ratios should be used for various habit types. For example, \citet{Hogan2012} used oblate spheroids to represent snow aggregates in their radar scattering calculations. They found that aggregates could be sufficiently represented as oblate spheroids with aspect ratios of 0.6 to simultaneously match observations of radar differential reflectivity and dual-wavelength ratios compared to in situ aircraft particle size distribution measurements. This value of 0.6 happens to be consistent with many in situ aircraft and ground-based studies that image ice particles such as \citet{Korolev2000,Korolev2003,Garrett2015}. Thus, it is natural to think that aggregate geometry is \textit{best} represented as a 0.6 aspect ratio oblate spheroid. However, more recent ground-based measurements utilizing artificial intelligence \citep{Jiang2019,Dunnavan2019b} and numerical studies \citep{Dunnavan2019b,Przybylo2022,Schrom2023} have elucidated the best-fit ellipsoids for aggregates. These studies consistently show that fully-developed aggregates are typically not oblate but rather triaxial. \citet{Dunnavan2019b} and \citet{Jiang2019} in particular demonstrate that a bivariate beta distribution of ellipsoid aspect ratios is capable of describing the entire distribution of ellipsoid aspect ratios for most stages of aggregate development. Thus, it is clear that one singular spheroidal aspect ratio might not be sufficient to describe multiple physical properties. This is especially true as radiative studies such as \citet{Macke1998} have found that the diversity of ice particle shapes, rather than explicit size distribution representations, dominates the variability in broadband solar radiative transfer calculations for cirrus clouds. 

%Recent studies such as \citet{Dunnavan2022} also suggest that irregular particles also show such a dispersion of ellipsoidal shapes. Furthermore, \citet{Dunnavan2022} showed that current aspect ratio polarimetric retrieval algorithms for irregular particles for one flight leg of a winter storm were not seemingly consistent with in situ data.

% As a result, there is a tension between treating ice particle shapes in models with generality (i.e., as reduced-density spheroids or with empirical power-law relations) or with fedelity (i.e., using a limited number of explicit model geometries).

In situ aircraft cloud probe instruments are well suited for studying the morphology of ice particles as they can image millions of atmospheric particles spanning large spatial distances. The Two Dimensional Stereo (2D-S) probe \citep{Lawson2006} is particularly well suited for this purpose for two reasons. First, optical array probes (OAPs) such as the 2D-S are capable of measuring many more particles than high-resolution probes such as the Cloud Particle Imager \citep[CPI;][]{Lawson2001} or the Particle Habit Imaging and Polar Scattering \citep[PHIPS;][]{Abdelmonem2016} instrument. Thus the statistics that can be derived from the 2D-S are naturally more robust than these higher-resolution probes. Secondly, the 2D-S is unique because of its two overlapping orthogonal channels -- named the horizontal, or ``H'', channel and the vertical, or ``V'', channel -- rather than one singular channel. Typically, the ``H'' channel is oriented such that images are projected onto a 2D plane (X-Z) that is parallel to the direction of aircraft heading and height whereas the ``V'' channel has images projected onto a 2D plane (X-Y) that it is parallel to the direction of the heading and width of the aircraft. The 2D stereo feature lends itself nicely to the study of ice particle shapes as it potentially allows for the simultaneous viewing of the same particle from two orthogonal directions.

%(see Fig.~\ref{fig:P3} for the 3D coordinate system used herein)

%During IMPACTS, the two available HVPS instruments HVPS-A (horizontally-installed) and HVPS-B (vertically-installed) were separate instruments that were displaced from each other on the port side of the P-3 aircraft.

% Discuss Alexei Korolev's series of papers. Mention Stolenga's paper which discusses the applicability of "irregular" habits. Discuss the aspect ratio dependence on size and how it differs between irregular and primary habits (Korolev and Isaac 2003; Harrington et al. 2013)

While the 2D-S was first introduced by \citet{Lawson2006} 20 years ago, the intended use of the dual channel and ``stereo'' operations of the instrument has still not been fully utilized in in situ aircraft studies. This is mainly because stereo sets of particle images are only available if particles simultaneously appear within both the H and V channel field of view. This overlap region is significantly smaller than the full sampling volume of either the H and V channels and thus any particular flight produces very few actual stereo (i.e., matched) images. Additionally, the 3D shape and orientation reconstruction problem itself requires the use of machine learning algorithms to optimize 3D features appropriately \citep[e.g.,][]{Jiang2019,Leinonen2021}. In fact, for assumed oblate or prolate spheroids, the information regarding the true 3D shapes can be completely obscured even for small changes in orientation and this uncertainty maximizes for the most nonspherical particles \citep{Dunnavan2019a}. Therefore, while shape estimation could be done with the 2D-S, the 2D-S so far is primarily used for the construction of size distributions and their associated moments as well as for habit identification. Often, the vertical channel is used as it is assumed that particles, despite any potential local turbulence due to the aircraft, will have their major axes aligned along the horizontal. Thus, it is assumed that the vertical channel will capture both the particle maximum dimension and will better capture the particle habit (e.g., a better view of a plate basal facet). However, very few studies exist that attempt to understand the 3D orientations of particles observed by optical array probes. Early attempts at such an understanding are provided by W. D. King in a series of papers on the matter \citep{King1985a,King1985b,King1985c,King1985d}. \citet{King1985d} in particular developed a theoretical horseshoe vortex model of airflow which was used to describe the predicted equilibrium angles of columns for probes installed under wing tips and near fuselages. This model predicted that the orientation of columnar particles was not horizontal but rather preferentially angled at around $20^{\circ}$ off horizontal when probes were installed horizontally. This expected equilibrium angle from horizontally-installed probes depends crucially on the location of the probe away from the wing both vertically and horizontally. When probes are installed vertically, as is most common for optical array probes, the predicted column equilibrium angle is more stable for probes installed under wing tips and the \citet{King1985d} model predicted that columns would be preferentially oriented at around $60^{\circ}$ for these vertically-installed probes. \citet{King1985d}, \citet{MacPherson1988}, and more recently \citet{Obrien2016} found good agreement between this simple model and the observed $60^{\circ}$ angle from vertically-installed OAPs. Thus, at least for columnar particles, this preferential orientation contradicts the assertion that particles are randomly or horizontally oriented within OAP volumes. It is currently unclear how these probe-induced orientations would affect estimations of the 3D crystal shapes.

This study seeks to resolve the 2D projection issue by utilizing a large number of 2D-S particle images from 26 cases of NASA's Investigation of Microphysics and Precipitation for Atlantic Coast-Threatening Snowstorms (IMPACTS) field campaign.\footnote{While two separate, orthogonal High-Volume Precipitation Spectrometer (HVPS-3) instruments were available during IMPACTS, we instead focus on only the 2D-S as the 2D-S has better pixel resolution and shares overlapping sampling volumes between the H and V channels.} The stereo capabilities of the 2D-S are reimagined not in terms of matched H and V images but rather in terms of the joint 2D distributions of ellipse fit properties (i.e., ellipse major axis sizes, aspect ratios, orientations, and area ratios) from each channel. In essence, this amounts to marginalizing the complete \textit{matched} 4D observed and simulated joint distributions of, for example, ellipse aspect ratios and orientations, from each channel into 2 separate bivariate distributions -- one for the H channel and the other for the V channel. While this necessarily severs the particle-by-particle information between each channel, it does preserve the large statistics provided by each flight dataset. In addition to presenting and analyzing these 2D ellipse distributions, we also provide Monte Carlo simulations using both a multivariate, analytical ellipsoid distribution model as well as detailed STL models fitted with corresponding minimum volume enclosing triaxial ellipsoids. The purpose of these simulations is not to optimize the 3D shapes and orientations using the 2D projections, but rather to explore what 3D particle characteristics can \textit{plausibly} explain the joint distributions of the H and V ellipse parameters. In this way, we use the 3D simulations as a tool for testing current assumptions regarding 3D ice particle shapes, densities, and in situ orientations and to propose plausible explanations for different habit types based upon the H and V channel ellipse statistics. As a result, this study aims to provide insight and guidance regarding the true morphology and orientations of 3D ice particles with their 2D H and V channel projections and how the 2D ellipse analogs map to their 3D ellipsoid and true morphological counterparts.

% As such, we aim to understand the 3D shapes and orientations on a purely \textit{statistical} level through the juxtaposition of observed and simulated H and V bivariate distributions. 

%In addition, we also produce these types of bivariate distributions including ellipse area ratios as well. 

%The ellipse area ratio can be used to potentially provide information regarding the density of each particle.

% We have developed Matlab programs to simulate similar theoretical distributions with the use of triaxial ellipsoids to potentially explain the in situ-derived observations. Readers can better understand the ellipsoid projection problem and analytical model used in this study by installing Matlab apps located in S1. Alternatively, readers can directly download these apps by searching for ``SingleEllipsoidProjection'' and ``MultipleEllipsoidProjections'' on the ``Get More Apps'' tab within the Matlab console.

% NOTE: check to see if we can use low area ratio + high aspect ratio to determine hexagonal plates.


\section{Methods}

\subsection{Aircraft coordinate system}

Figure~\ref{fig:P3} provides a schematic that shows the 2D-S location on the P-3 aircraft during the NASA IMPACTS campaign and the coordinate system used throughout. The NASA IMPACTS campaign was a large field campaign that operated primarily in the Northeast as well as several Midwestern portions of United States in 2020, 2022, and 2023 \citep{McMurdie2019,McMurdie2022}. During the IMPACTS campaign, the 2D-S probe was mounted under the P-3 aircraft's starboard-side wing tip. The coordinate system used here has the z axis pointed upward whereas the x axis is pointed along flight direction and the y axis pointed toward port side. The 2D-S (highlighted by a rectangular box in Fig.~\ref{fig:P3}a) shares this same (x,y,z) coordinate basis where the vertical and horizontal photodiode arrays (i.e., receivers) were installed on the bottom and starboard side, respectively\footnote{Our analysis shown in Sect.~\ref{sect:obs} actually suggests that these photodiode arrays were mistakenly installed during the 2020 IMPACTS campaign such that there was effectively a 90$^{\circ}$ counter-clockwise rotation of the 2D-S compared to the intended orientation. This likely resulted in the H photodiode array being oriented on bottom and the V photodiode array oriented on the port side during the 2020 flights. We present compelling evidence for this assertion in Appendix B.}. Figure~\ref{fig:P3}b shows a close-up schematic of the 2D-S basis. We use this same basis for ellipsoid and STL model simulations of ice particles in Sect.~\ref{sect:simulations} (see Fig.~\ref{fig:P3}c). Figure~\ref{fig:P3}d shows how the 2D-S diode elements are aligned with respect to the aircraft/ellipse coordinate system. The 2D-S data itself is stored in a different coordinate system such that its x axis is along the -z (-y) H (V) channel aircraft axis and the y axis is along the aircraft x axis. Therefore, we perform the following reflection transformation to align the projected ellipse angle $\theta_{\mathrm{proj}} \in [-90^{\circ},+90^{\circ}]$ where $-90^{\circ} \equiv +90^{\circ}$ such that positive (negative) $\theta_{\mathrm{proj}}$ angles represent counter-clockwise (clockwise) deviations from the x-axis (i.e., aircraft heading direction) in the aircraft/ellipse basis   
\begin{equation}
\theta_{\mathrm{proj}} = 
\begin{cases}
90^{\circ} - \theta_{\mathrm{proj},\mathrm{diode}}, & \theta_{\mathrm{proj},\mathrm{diode}} \geq 0^{\circ}\\
-(90^{\circ} - |\theta_{\mathrm{proj},\mathrm{diode}}|), & \theta_{\mathrm{proj},\mathrm{diode}} < 0^{\circ}.
\end{cases}
\end{equation}


%$-90^{\circ} < \theta_{\mathrm{proj}} < +90^{\circ}$

% The 2D-S coordinate system is such that  the z-axis points down.

% H and V channel 2D-S images are both from 20230113 between 054243 UTC and 054244 UTC.

%(from \url{https://airbornescience.nasa.gov/content/3D_Models_Gallery})

\FloatBarrier
\begin{figure}[ht]
    \centering
    \includegraphics[width=0.9\textwidth]{Figures/P3_WFF_IMPACTS_with_column_ex_rotated_newcoords_MODIFIED_BUFFER_STRIPS_REORG.pdf}
    \caption{3D model of the NASA P-3 aircraft \citep{NASAP3Model} (a) with the 2D-S cloud probe instrument (b). The 2D-S laser transmitters are shown as ``V'' and ``H'' in panel b and the photodiode receivers are on the opposite arms. Panel c shows the corresponding x-y-z coordinate system and a simulation example with a representative STL hexagonal column and a best-fit ellipsoid with corresponding ellipse fits from the ellipsoid (green lines). Panel d shows how the H and V diode coordinates (inset blue ``X'' and ``Y'' basis) correspond to the aircraft and projected ellipse orientations where the H and V images themselves are fit with ellipses (red lines). Note that ellipse angles are defined based on the aircraft/ellipse basis.}
    \label{fig:P3}
\end{figure}
\FloatBarrier

\subsection{In situ data processing}\label{sect:2DS_processing}
We use a modified version of the University of Illinois/Oklahoma Optical array probe Processing Software \citep[UIOOPS;][]{McFarquhar2018} to process 2D-S data from 26 flights of the NASA IMPACTS campaign. UIOOPS is a Matlab-based library that decompresses the binary image base files and calculates particle-by-particle statistics for each distinct particle within 2D-S image buffer strips. Ellipse fits are computed by UIOOPS using the Computational Geometry Algorithms Library \citep[CGAL;][]{CGAL2026} C++ library. An inter-arrival threshold of 10 $\mu$s, a minimum number of 800 shadowed 2D-S pixels, a minimum ellipse major axis length of 0.35 mm, a minimum ellipse aspect ratio of 0.1, and a minimum ellipse area ratio of 0.2 are also used in order to minimize the number of shattered particles and to ensure that each particle is properly resolved.

%UIOOPS creates a Matlab executable file that calls the CGAL ``Min\textunderscore ellipse\textunderscore2'' class to calculate ellipse fits of minimum enclosure for each particle image.

% We have also introduced several data quality checks in the UIOOPS code that are not currently incorporated in the publically available version. In particular, it is quite common for the public UIOOPS code to fit ellipses around images where there are stray shadowed pixels far away from the particle or to fit ellipses around multiple particle images. To correct these instances, we first filter this noise by using the Matlab image processing function ``bwareaopen'' and by specifying that at least 5 pixels must be connected to the main image. We then use the Matlab image processing function ``bwconncomp'' to separately fit any image which contains more than one particle. Particle images are only used if the particles do not touch the edges of the 2D-S buffer strip; that is, the particle must be entirely inside the 2D-S sampling volume. 

The \citet{Schima2024} multi-probe habit classification algorithm was used to classify 2D-S particle habits as well as the irregular and linear STL Monte Carlo simulations. The algorithm is based on the \citet{Holroyd1987} algorithm which uses a decision tree to classify particle images as irregular, linear, dendrite, hexagonal, graupel, aggregate, oriented, or sphere. Classifying particles from the 2D-S is complicated by the limited resolution of the 2D-S and ambiguity related to inferring 3D morphology from 2D binary images. This ambiguity can be compounded by optical diffraction which can be significant for smaller particles. Therefore, this 2D-S classification algorithm is used to generally segregate various habits in the following analysis when appropriate.

%Therefore, for a given set of observations, there are in total 6 parameters governing the distribution of ice particle shapes and orientations $\mathbf{x}=(\overline{\varphi}_{ba},\overline{\varphi}_{cb},\varsigma, \overline{\phi}_{p}, \overline{\theta}_{p},\kappa)$. 

\subsection{Analytical shape and orientation models}\label{methods:math}

Ice particles are modeled as tri-axial ellipsoids with semi-axis lengths: $a$, $b$, and $c$, where $a\geq b \geq c$. Two aspect ratios can therefore characterize a particular ellipsoid shape. As shown in \citet{Jiang2019} and \citet{Dunnavan2019b}, observations of ice particle aggregates and Monte Carlo simulated aggregates indicate that beta distributions can accurately characterize the entire distribution of possible shapes with 
\begin{subequations}\label{eq:beta_asp_dists}
    \begin{equation}\label{n_ba}
        p(\varphi_{ba}|\alpha_{ba},\beta_{ba}) = \frac{1}{B(\alpha_{ba},\beta_{ba})} \varphi_{ba}^{\alpha_{ba}-1} (1-\varphi_{ba})^{\beta_{ba}-1}, \quad 0 < \varphi_{ba} \leq 1
    \end{equation}
    \begin{equation}\label{n_cb}
    p(\varphi_{cb}|\alpha_{ba},\beta_{ba},\beta_{cb}) = \frac{1}{B(\alpha_{ba}+\beta_{ba},\beta_{cb})} \varphi_{cb}^{\alpha_{ba}+\beta_{ba}-1} (1-\varphi_{cb})^{\beta_{cb}-1}, \quad 0 < \varphi_{cb} \leq 1
    \end{equation}    
    \begin{equation}\label{n_ca}
    p(\varphi_{ca}|\alpha_{ba},\beta_{ba},\beta_{cb}) = \frac{1}{B(\alpha_{ba},\beta_{ba}+\beta_{cb})} \varphi_{ca}^{\alpha_{ba}-1} (1-\varphi_{ca})^{\beta_{ba}+\beta_{cb}-1},  \quad 0 < \varphi_{ca} \leq \varphi_{ba} 
    \end{equation}
\end{subequations}
where $\varphi_{ba} \equiv b/a$, $\varphi_{cb} \equiv c/b$, $\varphi_{ca} \equiv c/a$ and $B(\alpha,\beta)$ is the beta function. Throughout this study we assume that these beta distributions can be used to represent the variety of observed particles. Note that for oblate spheroids, $\varphi_{ba} = 1.0$ and $\varphi_{cb} = \varphi_{ca}$ whereas for prolate spheroids, $\varphi_{cb}=1.0$ and $\varphi_{ba} = \varphi_{ca}$. For the independent random variates $\Phi_{ba}$ and $\Phi_{cb}$ which follow probability distributions Eqs.~\ref{n_ba} and~\ref{n_cb}, it can be shown from the Mellin convolution theorem that $\Phi_{ca} = \Phi_{ba} \Phi_{cb}$ and that $\Phi_{ca}$ follows Eq.~\ref{n_ca}. Therefore, while $p(\varphi_{ba},\varphi_{cb}) = p(\varphi_{ba}) p(\varphi_{cb})$,  $\Phi_{ba}$ and $\Phi_{ca}$ are jointly correlated such that they follow the bivariate probability distribution

\begin{equation}
    p(\varphi_{ba},\varphi_{ca}|\alpha_{ba},\beta_{ba},\beta_{cb}) = \frac{1}{B(\alpha_{ba},\beta_{ba}) B(\alpha_{ba}+\beta_{ba},\beta_{cb})} \varphi_{ca}^{\alpha_{ba}+\beta_{ba}-1} (\varphi_{ba}-\varphi_{ca})^{\beta_{cb}-1} \varphi_{ba}^{-\beta_{ba}-\beta_{cb}} (1-\varphi_{ba})^{\beta_{ba}-1}.
\end{equation}

%\FloatBarrier
%\begin{figure}[ht]
%    \centering
%    \includegraphics[width=\textwidth]%{Figures/MASC_model_bivariate_dist_pics.png}
%    \caption{Caption}
%    \label{fig:bivbeta_example_agg}
%\end{figure}
%\FloatBarrier

%\FloatBarrier
%\begin{figure}[ht]
%    \centering
 %   \includegraphics[width=\textwidth]{Figures/Ellipsoid_bivbeta_hailstones.pdf}
%    \caption{Caption}
%    \label{fig:bivbeta_examples}
%\end{figure}
%\FloatBarrier

It is useful to parameterize $\alpha_{ba}$, $\beta_{ba}$, and $\beta_{cb}$ in terms of mean quantities: $\overline{\varphi}_{ba}$ and $\overline{\varphi}_{cb}$, where each mean aspect ratio can be independently specified (exclusively) between 0 and 1.0. However, because the beta distributions are only valid for $\alpha >0$ and $\beta > 0$, specifying these mean aspect ratios will necessarily lead to a valid range of $\alpha_{ba}$ between some minimum value $\alpha_{ba,\min}$ to some maximum value $\alpha_{ba,\max}$. $\alpha_{ba,\min}$ is determined by the chosen $\overline{\varphi}_{ba}$ and $\overline{\varphi}_{cb}$ values whereas $\alpha_{ba,\max}$ is chosen to prevent numerical difficulties as the beta distributions become more and more narrow. Thus, we parameterize $\alpha_{ba}$, $\beta_{ba}$, and $\beta_{cb}$ as
\begin{subequations}\label{eq:ar_beta_params}
    \begin{equation}
        \alpha_{ba} = \varsigma \alpha_{ba,\max} + (1-\varsigma) \alpha_{ba,\min}
    \end{equation}
    \begin{equation}
        \beta_{ba} = \alpha_{ba} \left(\frac{1}{\overline{\varphi}_{ba}}-1 \right)
    \end{equation}
    \begin{equation}
        \beta_{cb} = \frac{\alpha_{ba}}{\overline{\varphi}_{ba}} \left( \frac{1}{\overline{\varphi}_{cb}} - 1\right), 
    \end{equation}
\end{subequations}
where 
\begin{subequations}
    \begin{equation}
        \alpha_{ba,\min} = \max(\varphi_{ba,\min},\varphi_{cb,\min})
    \end{equation}
    \begin{equation}
        \varphi_{ba,\min} = \frac{\overline{\varphi}_{ba}}{1-\overline{\varphi}_{ba}}
    \end{equation}
    \begin{equation}
        \varphi_{cb,\min} = \frac{\overline{\varphi}_{ba} \overline{\varphi}_{cb}}{1-\overline{\varphi}_{cb}},
    \end{equation}
\end{subequations}
and $\alpha_{ba,\max}$ is some upper limit of $\alpha_{ba}$ (herein set to 20) to prevent numerical difficulties. Therefore, the entire distribution of ellipsoidal shapes can be determined by the mean aspect ratios $\overline{\varphi}_{ba}$ and $\overline{\varphi}_{cb}$ and $\varsigma$ which determines the amount of dispersion inherent in the chosen bivariate beta distribution. Larger values of $\varsigma$ indicate a more concentrated (less dispersed) distribution of aspect ratios whereas smaller values of $\varsigma$ indicate a less concentrated (more dispersed) distribution of aspect ratios.

The above approach is used to simulate ellipsoids which generally represent irregular crystals. However, the same approach is modified for hexagonal plates and columns where power-law relations are used to simulate ellipsoid proxies and to appropriately scale the aspect ratios of STL models for  hexagonal plates and columns using the relations
\begin{subequations}\label{eq:power-law}
\begin{equation}\label{eq:power-law_PLATES}
\textrm{Plates:} \qquad \overline{\varphi}_{cb} = 0.0392 D^{-0.602} \qquad (b=a)
\end{equation} 
\begin{equation}\label{eq:power-law_COLUMNS}
 \textrm{Columns:} \qquad \overline{\varphi}_{ba} = 0.2201 D^{-0.5} \qquad (b=c),
\end{equation}
\end{subequations}
where $D\equiv 2 a$ in mm. These plate and column aspect ratio relations are derived from those of \citet{Auer1970} and \citet{Mitchell1994a}, respectively. For the ellipsoid simulations of linear (columnar) particles, the mean aspect ratio in Eq.~\ref{eq:power-law_COLUMNS}, $\overline{\varphi_{cb}}=0.95$, and $\varsigma=0.5$ are used to randomly sample $\varphi_{ba}$ and $\varphi_{cb}$ after the ellipsoid semi-major axis $a$ is sampled. Equation~\ref{eq:power-law} is also used to specify STL hexagonal column (i.e., linear) and plate aspect ratios where $D$ is the caliper maximum dimension.

%These parameterizations take the following form
%\begin{subequations}
%\begin{equation}
%\overline{\varphi}_{ba} = s_{ba} (2 a)^{h_{ba}}
%\end{equation}
%\begin{equation}
%\overline{\varphi}_{cb} = s_{cb} (2 a)^{h_{cb}},
%\end{equation}
%\end{subequations}
%where $s_{ba}$ ($s_{cb}$) and $h_{ba}$ ($h_{cb}$) are the prefactor and exponents of the aspect ratio-size for $\overline{\varphi}_{ba}$ ($\overline{\varphi}_{cb}$). 

%Effective densities for each ellipsoid can, in theory, be recovered through a dimensionality scaling \textit{ansatz} where
%\begin{equation}
%\mathrm{AR} = \rho_{\mathrm{eff}}^{1/3}.
%\end{equation}
%This ansatz essentially represents the average amount of space that a 3D particle would take up relative to a 2D projection (i.e., a projection along the $y$ axis represents 1 spatial dimension compared to 3 spatial dimensions of the particle). Theoretical support for this scaling \textit{ansatz} is provided in \cref{sect:simulations}.

2D projected ellipse area ratios ($\mathrm{AR}$) for idealized ellipsoids are simulated directly by similarly using beta distributions. The 2D ellipse area ratios are bounded according to an assumed minimum value $\mathrm{AR}_{\min}$ herein set to 0.1. Then, the target mean ellipse area ratio $\overline{\mathrm{AR}}$ is normalized to this interval 
\begin{equation}\label{eq:AR_mean}
\widehat{\mathrm{AR}} = \frac{\overline{\mathrm{AR}}-\mathrm{AR}_{\min}}{1-\mathrm{AR}_{\min}}.
\end{equation}
A dispersion parameter $\nu_{\mathrm{AR}}$ acts as a proxy for the distribution variance and is defined using a control parameter $\varrho_{\mathrm{frac}}$
\begin{equation}\label{eq:nu_AR}
\nu_{\mathrm{AR}} = \max\left(\rho_{\mathrm{frac}}^{-1} - 1, 2\right).
\end{equation}
The area ratio beta distribution parameters are then determined from
\begin{subequations}\label{eq:beta_AR}
\begin{equation}
\alpha_{\mathrm{AR}} = \widehat{\mathrm{AR}}\, \nu_{\mathrm{AR}}
\end{equation}
\begin{equation}
\beta_{\mathrm{AR}} = \left(1-\widehat{\mathrm{AR}}\right) \nu_{\mathrm{AR}}.
\end{equation}
\end{subequations}
A normalized random variable is then sampled as $X \sim \mathrm{Beta}\left(\alpha_{\mathrm{AR}},\beta_{\mathrm{AR}} \right)$ which is then mapped to calculate the area ratio using
\begin{equation}\label{eq:AreaRatio}
\mathrm{AR} = \mathrm{AR}_{\min} + X\left( 1 - \mathrm{AR}_{\min} \right).
\end{equation}
While both the H and V channels use the same AR parameterization and parameter values, each projection is separately drawn from the AR beta distribution. If the ellipsoid effective density $\rho_{\mathrm{ellip}} \sim \mathrm{AR}^{3/2}$, then $\rho_{\mathrm{ellip}}$ will follow a generalized beta distribution. 

%Effective ellipsoid densities are used to estimate projection ellipse area ratios. We assume that the 3D ellipsoid effective densities are beta-distributed ($p(\rho) \sim \mathrm{Beta}(\alpha_{\rho},\beta_{\rho})$). In this way, the beta distribution is used to represent Gaussian-like variability in the area ratios similar to what might be expected from 2D-S observations. As in Equation~\ref{eq:ar_beta_params}, we use the method-of-moments approach to parameterize the beta distribution parameters in terms of the mean effective density $\overline{\rho}$ and a fractional variance $\varrho_{\mathrm{frac}}$ which modulates the maximum allowable variance as

%\begin{equation}
%\varrho = \min(\varrho_{\mathrm{frac}} \varrho_{\mathrm{max}},20).
%\end{equation}

%The maximum allowable variance is determined similarly to \cref{eq:ar_beta_params} where
%\begin{equation}
%\varrho_{\mathrm{max}} = \frac{\overline{\rho}^{2} - \overline{\rho}^{3} }{\overline{\rho}+1}
%\end{equation}
%and
%\begin{subequations}
%\begin{equation}
%\alpha_{\rho} = \overline{\rho}\left(\frac{\overline{\rho}(1-\overline{\rho})}{\varrho} -1 \right)
%\end{equation}
%\begin{equation}
%\beta_{\rho} = (1-\overline{\rho})\left(\frac{\overline{\rho}(1-%\overline{\rho})}{\varrho} -1 \right).
%\end{equation}
%\end{subequations}


The geometry for the particle basis is shown in Fig.~\ref{fig:ellip_orient} where the $c$ (blue arrow), $b$ (green arrow), and $a$ (red arrow) axes (defined by $a\geq b \geq c$) of an ice particle ellipsoid (green shell) are first aligned along the coordinate reference axes $x$, $y$, and $z$, respectively (Fig.~\ref{fig:ellip_orient}a). This basis ensures that the first two Euler angle rotations, $\phi_p$ (Fig.~\ref{fig:ellip_orient}c) and $\theta_p$ (Fig.~\ref{fig:ellip_orient}b), specify a vector $\mathbf{z^{\prime}}$ representing the semi-major ellipsoid axis $a$, is rotated to point in the direction of a spot on a circumscribing sphere with radius $a$ dictated by a spherical coordinate system given by $(\phi_p,\theta_p,a)$. The final Euler angle rotation, $\psi_p$ (Fig.~\ref{fig:ellip_orient}d), is then used to complete the arbitrary orientation of the 3D ellipsoid where $\psi_p$ is a rotation about $\mathbf{z^{\prime}}$. In this way, $0\leq \theta_p \leq 180^{\circ}$ and $0 \leq \phi_p (\psi_p) \leq 360^{\circ}$. Figure~\ref{fig:ellip_orient} shows a schematic for this Euler angle convention. It is important to realize that triaxial ellipsoids are symmetric about their origins. Thus, in the reference frame of the ellipsoid (before rotation), the shape is reflected from the negative poles (i.e., $x=-c$, $y=-b$, and $z=-a$) to the positive poles (i.e., $x=+c$, $y=+b$, and $z=+a$). The Euler angle rotations obviously preserve this symmetry about the $\mathbf{x^{\prime}}$, $\mathbf{y^{\prime}}$, and $\mathbf{z^{\prime}}$ principal axes. As a result, the full range of $\phi_p$ and $\psi_p$ can be halved to avoid double counting each unique ellipsoid orientation and its projections (i.e., restricting $\phi_p$ and $\psi_p$ from $0^{\circ}$ to $180^{\circ}$). 

%This is especially important when training the ANN model to preserve a 1-1 mapping between probability distribution parameters and the projection probability distributions.

%The Euler angle convention chosen here thus can be characterized by a point on a sphere with radius $a$, dictated by spherical coordinate angles ($\theta_p$, $\phi_p$) and an additional angle, $\psi_p$, rotated about this vector.

 We use a 3-dimensional von Mises-Fisher (or simply, Fisher) probability distribution to specify and randomly draw the first two Euler angles $\phi_p$ and $\theta_p$. This distribution is expressed as

\begin{equation}\label{eq:fisher}
    p(\mathbf{x}| \mathbf{r},\kappa) = \frac{\kappa}{4 \pi \sinh\kappa} \exp(\kappa\, \mathbf{r} \cdot \mathbf{x}),
\end{equation}
where $\mathbf{x} = (x,y,z)$ represents the Cartesian coordinates of a sphere
\begin{subequations}\label{eq:xyz}
    \begin{equation}
        x = \sin \theta_p \cos \phi_p
    \end{equation}
        \begin{equation}
        y = \sin \theta_p \sin \phi_p
    \end{equation}
        \begin{equation}
        z = \cos \theta_p .
    \end{equation}
\end{subequations}
Here, $\mathbf{r} \equiv \mathbf{z^{\prime}}/||\mathbf{z^{\prime}}||$ is the unit-normalized Cartesian $\mathbf{z^{\prime}}$ vector specified by the mean $\phi_p$ ($\overline{\phi}_{p}$) and mean $\theta_p$ ($\overline{\theta}_{p}$). $\kappa$ is a concentration parameter where low values indicate more dispersion and high values indicate more concentrated values about $\overline{\phi}_p$ and $\overline{\theta}_p$. The Fisher distribution tends toward a random distribution of points on a sphere as $\kappa \to 0$ (herein, $\kappa=0.001$ is used to represent such a random distribution). Figure~\ref{fig:fisher} schematically shows how the Fisher probability distribution changes for various $\kappa$ values using a fixed mean orientation direction specified by ($\overline{\phi}_{p},\overline{\theta}_p$).  

\FloatBarrier
\begin{figure}[ht]
    \centering
\includegraphics[width=0.95\textwidth]{Figures/Ellipsoid_orient_coord_sys_new.png}
    \caption{Euler angle rotation convention used throughout using $\phi_p = 60^{\circ}$, $\theta_p = 70^{\circ}$, and $\psi_p = 160^{\circ}$ as example rotations for an ellipsoid with aspect ratios of $\varphi_{ba} = 0.5$ and $\varphi_{cb}=0.7$. a.) represents the initial orientation, b.) represents the $\theta_{p}$ rotation, c.) represents the $\phi_p$ rotation, d.) represents the $\psi_p$ rotation. Note that computationally the order of b.) and c.) is reversed but is shown in this way to better illustrate each rotation.}
    \label{fig:ellip_orient}
\end{figure}
\FloatBarrier

$\psi_{p}$ can be characterized by another von Mises-Fisher distribution with a dimensionality of 2 (otherwise known as simply the von Mises distribution) 
\begin{equation}\label{eq:p_psi}
    p(\psi_{p}|\overline{\psi}_{p},\kappa_{\psi}) = \frac{1}{2\pi I_{0}(\kappa_{\psi})} e^{\kappa_{\psi} \cos (\psi_{p} - \overline{\psi}_{p})},
\end{equation}
where $I_{0}$ is the first modified Bessel function of order 0. Equation~\ref{eq:p_psi} tends toward a uniform distribution as $\kappa_{\psi} \to 0$. Throughout, we use a lower limit of $\kappa_{\psi} = 0.001$ to represent a random $\psi_{p}$ orientation distribution and $\kappa = \kappa_{\psi} = 0.001$ to represent completely random (i.e., chaotic) orientations.


\FloatBarrier
\begin{figure}[ht]
    \centering
    \includegraphics[width=\textwidth]{Figures/Fisher_orient_ex_ellip_kap_multi.png}
    \caption{Examples of the Fisher distribution (Equation~\ref{eq:fisher}) for $\overline{\phi}_p=45^{\circ}$ and $\overline{\theta}_{p} = 70^{\circ}$ and for three different values of $\kappa$.}
    \label{fig:fisher}
\end{figure}
\FloatBarrier

The combination of Eqs.~\ref{eq:beta_asp_dists},~\ref{eq:fisher},  and~\ref{eq:p_psi} can completely specify an entire multivariate distribution of ellipsoid aspect ratios and orientations for the 8-parameter set $(\overline{\varphi}_{ba},\overline{\varphi}_{cb},
\varsigma,\overline{\phi}_{p},\overline{\theta}_{p},\kappa,\overline{\psi}_{p},\kappa_{\psi})$ whereas area ratios only require specifying the two-parameter set $(\overline{AR},\rho_{\mathrm{frac}})$ in Eqs.(\ref{eq:AR_mean})--(\ref{eq:AreaRatio}). Unless otherwise stated, $\overline{\varphi}_{ba}$ and $\overline{\varphi}_{cb}$ use fixed values whereas plate and linear habit simulations utilize Eq.~\ref{eq:power-law_PLATES} and Eq.~\ref{eq:power-law_COLUMNS}, respectively.
Simulating distributions of projected ellipse aspect ratios, area ratios, and orientation angles therefore only requires drawing random variables from these multivariate distributions, calculating the corresponding H (i.e., x-z plane) and V (i.e., x-y plane) ellipse aspect ratio and orientation angle for each drawn sample, and binning the corresponding ellipse projection aspect ratio, area ratio, and orientation if the ellipses pass the filtering procedure described in Section~\ref{sect:MC_sim}. The mapping from ellipsoid to ellipse projection is analytical and given explicitly in Appendix~\ref{appen:A}.

If we assume that the major axis $a$ is projected in the same direction as each ellipse, then Eq.~\ref{eq:xyz} can be used to estimate the 3D mean orientation angles in Eq.~\ref{eq:fisher} solely from the observed mean 2D ellipse angles from the ellipse fit distributions. These estimates are given by

% x-z and x-y plane 3D retrieval estimates
\begin{subequations}\label{eq:ret_orient}
    \begin{equation}\label{eq:ret_orient_phi}
	\overline{\phi}_{p} \approx \begin{cases} 

	 \overline{\theta}_{\mathrm{proj},v}, & \overline{\theta}_{\mathrm{proj},v} \geq 0^{\circ}, \\
	 \overline{\theta}_{\mathrm{proj},v} + 180^{\circ},  & \overline{\theta}_{\mathrm{proj},v}  < 0^{\circ}
	\end{cases}    
    \end{equation}
    \begin{equation}\label{eq:ret_orient_theta}
        \overline{\theta}_{p} \approx \mathrm{atan2}(\sec \overline{\theta}_{\mathrm{proj},v},\tan \overline{\theta}_{\mathrm{proj},h}),
    \end{equation}
\end{subequations}
where $\mathrm{atan2}$ is the two-argument arctangent function defined by
\begin{equation}
    \mathrm{atan2}(y,x) = \begin{cases}
        \arctan(y/x), & \textrm{if } x>0, \\
        \arctan(y/x)+180^{\circ} & \textrm{if } (x<0) \land (y\geq 0), \\
        \arctan(y/x) - 180^{\circ} & \textrm{if } (x<0) \land (y<0),\\
        90^{\circ} & \textrm{if } (x=0) \land (y>0),\\
        -90^{\circ} & \textrm{if } (x=0) \land (y<0),\\
        \textrm{undefined} & \textrm{if } (x=0) \land (y=0).
    \end{cases}
\end{equation}
The use of Eq.~\ref{eq:ret_orient} allows us to more efficiently determine \textit{plausible} orientation parameters for a given set of ellipsoid parameters from the observations without needing to guess or test multiple orientation parameter sets or otherwise optimize the 8-parameter multivariate ellipsoid shape and orientation and 2-parameter area ratio parameters. As such, this does restrict our analysis to only using a single set of parameters for representing a particular habit class from a particular set of observations. However, results in Sect.~\ref{sect:simulations} illustrate that, at least for irregular and linear habits, the actual orientation distributions \textit{must} be very precise given the observed 2D ellipse orientation distributions and our overall comparisons very closely mimic the joint distribution observations between each combination of the four ellipse parameters. 

%Orientation parameters that are described as ``(fixed)'' refer to the orientation angle itself being fixed rather than drawn from the Fisher distribution.

%Irregular 2020 and bullet rosette 2020 refers to retrieved orientation parameters using \cref{eq:ret_orient} for the 2020 dataset assuming that the 2020 2D-S was oriented according to the schematic in \cref{fig:P3}c.

%\FloatBarrier
%\begin{table}[ht]
%\centering
%\caption{Orientation distribution parameters used for ellipsoid and STL simulations. Irregular 2020 and Columns 2020 refers to retrieved orientation parameters using \cref{eq:ret_orient} for the 2020 dataset assuming that the 2020 2D-S was oriented according to the schematic in \cref{fig:P3}c.}
%\begin{tabular}{c c c c c c c c c c c}
%\toprule
%\toprule
%Name & $\theta_{\mathrm{proj},h}$ & $\theta_{\mathrm{proj},v}$ &$\overline{\phi}_{p}$ & $\overline{\theta}_{p}$ &$\kappa$ & $\overline{\psi}_{p}$ & $\kappa_{\psi}$ & %POP frac. \\
%\midrule
%Planar Flutter & --- & --- & 0$^{\circ}$ & %$90^{\circ}$ & 6 & $0^{\circ}$ & 6 & --- \\
%Irregular 2020 & -40$^{\circ}$ & %45$^{\circ}$ &  45$^{\circ}$ & %$120.68^{\circ}$ &  3.5 & $0^{\circ}$ & %0.001 & ---\\
%Columns 2020 & -70$^{\circ}$ & 80$^{\circ}$ %&  80$^{\circ}$ & $115.506^{\circ}$ &  8 & %$0^{\circ}$ & 0.001 & ---\\
%Irregular 2022-2023 POP 1 & -40$^{\circ}$ & %-10$^{\circ}$ &  170$^{\circ}$ & %$50.43^{\circ}$ &  3.5 & $0^{\circ}$ & 0.001 & 0.75\\
%Irregular 2022-2023 POP 2 & -70$^{\circ}$ & -75$^{\circ}$ &  105$^{\circ}$ & $54.58^{\circ}$ &  10 & $0^{\circ}$ & 0.001 & 0.25\\
%Columns 2022-2023 & -80$^{\circ}$ & %-80$^{\circ}$ &  100$^{\circ}$ & $45.44^{\circ}$ &  8 & $0^{\circ}$ & 0.001 & %---\\
%Irregular 2022-2023 & 45$^{\circ}$ & %-10$^{\circ}$ &  45$^{\circ}$ & $116.341^{\circ}$ &  4 & $0^{\circ}$ & 0.001\\
%Columns 2022-2023 & -70$^{\circ}$ & 80$^{\circ}$ &  80$^{\circ}$ & $115.506^{\circ}$ &  3 & $0^{\circ}$ & 0.001\\
%Irregular Parallel $\overline{\psi}_{p}$ & %-35$^{\circ}$ & 45$^{\circ}$ &  45$^{\circ}$ & $116.341^{\circ}$ &  4 & $0^{\circ}$ & 10 & ---\\
%Irregular Skewed $\overline{\psi}_{p}$ & -35$^{\circ}$ & 45$^{\circ}$ &  45$^{\circ}$ & $116.341^{\circ}$ &  4 & $45^{\circ}$ & 10& ---\\
%Irregular Orthogonal $\overline{\psi}_{p}$ & -35$^{\circ}$ & 45$^{\circ}$ &  45$^{\circ}$ & $116.341^{\circ}$ &  4 & $90^{\circ}$ & 10 & ---\\
%\bottomrule
%\end{tabular}
%\label{table:ellip_orient_tests}
%\end{table}
%\FloatBarrier

%Exponential functions and power-law functions are additionally used to capture the size dependence of ellipsoid aspect ratios and effective densities that appear in the in situ observations. These functions take the form
%\begin{subequations}
%\begin{equation}\label{eq:fscale}
%f_{X} = e^{-a/a_{\mathrm{X}}}
%\end{equation} 
%\begin{equation}\label{eq:power-law}
% c = \alpha_{c} a^{\beta_{c}}.
%\end{equation}
%\end{subequations}

% but modulated according to the scaling factor shown in \cref{eq:fscale}

%This parameterization ensures that $\alpha_{\rho}>1$ such that $\lim_{\rho \to 0} p(\rho) = 0$.

Particle maximum dimensions ($D = 2 a$ for ellipsoids) are governed by gamma distributions which are parameterized as
\begin{equation}
\tilde{n}(D|D_{\mathrm{mean}},\nu) = \frac{1}{\Gamma(\nu)} \frac{\nu^{\nu}}{D_{\mathrm{mean}}} \left( \frac{D}{D_{\mathrm{mean}}}\right)^{\nu-1} e^{- \nu \frac{D}{D_{\mathrm{mean}}}},
\end{equation}
where $\nu$ is the gamma distribution shape parameter,  $D_{\mathrm{mean}}$ is the mean maximum dimension, and $\tilde{n}$ is the number-concentration-normalized number density function: $\tilde{n} \equiv n(D)/N_{t}$. The gamma distribution parameters $\nu$ and $D_{\mathrm{mean}}$ are chosen subjectively so as to best represent the observed filtered ellipse fit distributions.

%The $\nu$ parameters used throughout are generally  much larger than typical values that might be expected. These larger values of $\nu$ are due to the filtering of small particles less than $0.35$ mm as well as the additional quality control checks as described in the next section. 

% (\url{https://www.mathworks.com/matlabcentral/fileexchange/9542-minimum-volume-enclosing-ellipsoid})

\subsection{Synthetic 3D STL ice particles}\label{sect:STL}

In addition to the ellipsoid distributions described in Sect.~\ref{methods:math}, STL models are used to mimic various ice particle morphologies to investigate plausible explanations for the joint probability distributions shown in Sect.~\ref{sect:obs}. Each STL model is fit with a minimum volume enclosing trixial ellipsoid using the Matlab Community File Exchange function ``MinVolEllipse.m'' \citep{Moshtagh2024}. Orientations for each STL particle therefore follow the methodology outlined in Sect.~\ref{methods:math} where the Fisher-distributed Euler angles are used to orient both the ellipsoids and STL models after the STL model is resized and reoriented according to the best-fit ellipsoid parameters. The VOXELISE.m ``Mesh voxelisation'' Matlab Community File Exchange code \citep{Aitkenhead2024} is used to voxelize each STL particle after the particle is rotated using a 3D grid with a gridspacing of $10\mu$m in each spatial direction. This is done for a number of different habit types including: irregular, bullet rosettes, plates, dendrites, linear particles (i.e., columns and needles), and capped columns.

To represent irregular particles, we use the Spherical Harmonics Particle Shape Generator (SHPSG) python package \citep{Zhao_SHPSG} which is described in detailed and used in several papers including \citep{Zhao2015,Zhao2017}. The SHPSG package generates random 3D triangulation meshes using four prescribed parameters: an elongation aspect ratio (EI; analogous to $\varphi_{ba}$), a flatness aspect ratio (FI; analogous to $\varphi_{cb}$), a roughness parameter ($R_{G}$), and a roundness parameter ($R_{M}$). These roughness and roundness parameters are not traditional and are rather related to the spherical harmonic coefficients used to specify each 3D shape. When is 0, then the particles are idealized triaxial ellipsoids with aspect ratios governed by the elongation and flatness aspect ratio parameters. In this case, the resulting ellipsoid is considered a ``first-degree ellipsoid'' (FDE) characterized by $\varphi_{ba}$ and $\varphi_{cb}$. The FDE is not the same as the minimum-enclosing ellipsoids. Rather, FDE in this context represents a characteristic ellipsoid from which the SHPSG generates random perturbations based on the 16 spherical harmonic coefficients that are randomly drawn. Therefore, the FDE and best-fit ellipsoids are different from the same 3D STL particle unless $R_{M}=0$. For the SHPSG tests used in this study, we use $R_{M}=0.5$ and $R_{G}=0$ values to ensure volumes are finite as large $R_{G}$ values can yield unrealistic manifolds. By themselves, the STL models produce area ratios that are, strictly speaking, higher than the IMPACTS observations even when $R_{M}$ and $R_{G}$ are large. This, in part, is due to the inability of the SHPSG code to simulate branches and any potential hollowing effects. To generate more complicated spatial features such as branches, which decreases aspect ratios, we use the Python ``manifold3d'' package to randomly fuse together either three or four SHPSG-generated particles together. In this way, these irregular particles can better mimic the various polycrystaline shapes ubiquitously found in the IMPACTS dataset. 1,000 irregular SHPSG fused particles are randomly generated for $\mathrm{EI}$ and $\mathrm{FI}$ values which are randomly selected between 0.2 and 0.9.

%(\url{https://www.mathworks.com/matlabcentral/fileexchange/27390-mesh-voxelisation})

%For the SHPSG tests used in this study, we use $R_{G}=R_{M}=1.0$ as these parameters visually produce more complicated structures than used in \citet{Zhao2015} and \citet{Zhao2017} to represent sand granules. 10,000 irregular SHPSG particles are randomly generated for $\mathrm{EI}=\mathrm{FI}=0.75$.

Figure~\ref{fig:STL_irreg_diagram} provides two examples of SHPSG generated particles as well as two morphologically similar irregular particle images from the PHIPS imager. In general, SHPSG can generate particles that have similar random morphologies however there are some caveats. Various irregular particles such as side planes are difficult to accurately mimic due to the many gaps present in between branches of varying sizes and shapes. While fusing three to four SHPSG particles helps mimic these branches, SHPSG particles are still smoother than the 3D and 2D observational images but can still resemble the overall structure of many irregular particles. More complicated irregular-like particles such as bullet rosettes are not possible to generate accurately using SHPSG. Instead, bullet rosettes with branches between 3 and 6 are randomly generated by aligning a single bullet STL model at the bullet tip and randomly orienting each bullet ensuring a minimum angle between all bullets of at least 75$^{\circ}$, 70$^{\circ}$, 60$^{\circ}$, 50$^{\circ}$, 45$^{\circ}$, and 40$^{\circ}$ for 3 through 8 branches in that order. Capped columns are similarly randomly generated with column widths between 0.3 mm and 0.6 mm and lengths between 1.0 mm and 2.5 mm. The capped plates for these capped columns have widths that are randomly between 1.5 and 3 times that of each drawn column width along with thicknesses randomly sampled between 0.05 mm and 0.2 mm. Occasionally, the bullet rosette and fused SHPSG STL models will produce very slight gaps in the H and V projections where the individual STL models fuse together. To eliminate these gaps in the projection data, we use the Matlab ``imfill'' and ``imclose'' functions from the Image Processing and Computer Vision Matlab toolkit to fill in these gaps. Example STL models for each habit type are shown in Fig.~\ref{fig:stl_ex} along with similar PHIPS and CPI images.

%Single hexagonal plate and column aspect ratio are determined through the power-law parameters from \citet[][D74]{Davis1974} and \citet[][AV70]{Auer1970}, respectively. 

\begin{figure}[ht]
    \centering
 %   \includegraphics[width=0.7\textwidth]{Figures/IMPACTS_3D_STL_diagram_w_2DS_new.png}
        \includegraphics[width=0.7\textwidth]{Figures/IMPACTS_3D_STL_diagram_w_2DS_VECTORIZED.pdf}
    \caption{(a and b) Comparison between an example irregular particle observed by the PHIPS instrument and similar SHPSG particles along with best-fit ellipsoids and 2D ellipse fits in the H and V planes from the ellipsoid projections (green lines) and from the CGAL ellipse fitting algorithm on the H and V projected shadows (red lines). (c) Comparison between example of IMPACTS 2D-S images with ellipse fits and the SHPSG irregular particle images processed using the SODA-2 software \citep{Bansemer2026} in a synthetically similar buffer strip.}
    \label{fig:STL_irreg_diagram}
\end{figure}
\FloatBarrier

%We utilize various explicit STL models to represent plates, columns, needles, and dendrites. 

% NEW !
%\FloatBarrier
%\begin{table}[ht]
%\centering
%\caption{Ellipsoid and STL distribution parameters used for tests.}
%\begin{tabular}{c c c c c c c c c c}
%\toprule
%\toprule
%Name & $\nu$ & $D_{\mathrm{mean}}$ & $\overline{\varphi}_{ba}$ & $\overline{\varphi}_{cb}$ & $\zeta$ & $\overline{\mathrm{AR}}$ & $\varrho_{\mathrm{frac}}$ & $\alpha_{\varphi}$ & $\beta_{\varphi}$ \\
%\midrule
%Irregular 2020 & 1 & 2.0 mm & 0.75 & 0.75 & 0.6 & 0.59 & 0.02  &  --- & ---\\
%Irregular 2022 POP 1 & 1 & 2.0 mm & 0.78 & 0.78 & 0.6 & 0.59 & 0.02  &  --- & ---\\
%Irregular 2020 & 1 & 2.0 mm & 0.75 & 0.75 & 0.6 & 0.59 & 0.02  &  --- & ---\\
%Linear & 3 & 0.75 mm & --- & 0.95 & 0.5 & 0.65 & 0.03 & 0.2201 & -0.5\\
%Hexagonal Plates & 3 & 0.75 mm & --- & --- & --- & 0.65 & 0.03 &  0.0392 & 0.398 \\
%Dendrites & 4 & 1.0 mm & --- & --- & --- & 0.3 & 0.05 & --- & ---  &  & \\
%Needles & & & & & & & & & & & \\
%\bottomrule
%\end{tabular}
%\label{table:ellip_dist_params}
%\end{table}
%\FloatBarrier
% NEW !

% ORIGINAL!
%\FloatBarrier
%\begin{table}[ht]
%\centering
%\caption{Ellipsoid and STL distribution parameters used for tests.}
%\begin{tabular}{c c c c c c c c c c c c}
%\toprule
%\toprule
%Name & $\nu$ & $D_{\mathrm{mean}}$ & $\overline{\varphi}_{ba}$ & $\overline{\varphi}_{cb}$ & $\zeta$ & $\overline{\rho}$ & $\varrho_{\mathrm{frac}}$ & $a_{\varphi}$ & $a_{\rho}$ & $\alpha_{c}$ & $\beta_{c}$ \\
%\midrule
%Irregular & 8 & 0.6 mm & 0.75 & 0.75 & 0.6 & 0.8 & 0.12 & 2 & 4/3 & --- & ---\\
%%Hexagonal Columns & 4 & 0.5 mm & --- & --- & --- & 0.75 & 0.05 & --- & 4 & 0.1973 & 0.414\\
%%Hexagonal Plates & 4 & 0.5 mm & --- & --- & --- & 0.75 & 0.05 &  --- & --- & 0.0392 & 0.398 \\
%%Dendrites & 4 & 1.0 mm & --- & --- & --- & 0.3 & 0.05 & --- & ---  &  & \\
%%Needles & & & & & & & & & & & \\
%\bottomrule
%\end{tabular}
%\label{table:ellip_dist_params}
%\end{table}
%\FloatBarrier
% ORIGNAL !

% Sphericity values: BR03 = 0.5628, BR04 = 0.5463, BR08 = 0.4764, Dendrite = 0.1378, Needle = 0.5472, Capped Column = 0.4469

%\FloatBarrier
%\begin{table}[ht]
%\centering
%\caption{STL models and orientations used in %\cref{sect:simulations} with size distributions specified in \cref{table:ellip_dist_params}. Particle aspect ratios ($\varphi$) are defined in terms of the ratio of the prism facet length over the basal facet length. In this way, oblate particles like plates and dendrites have $\varphi < 1.0$ whereas prolate particles like columns and needles have $\varphi > 1.0$.}
%\begin{tabular}{c c c c c c}
%\toprule
%\toprule
%Habit Type (\# Const.) & $\varphi$ & $\varphi_{ba}$ $(\varphi_{cb})$ & Ellipsoid Volume Frac. & Orient. & Size Dist. \\
%\midrule
%Irregular (1) &   0.2** -- 0.9** & --- & --- & Irregular 2020 %& Irregular \\
%%Bullet Rosette (3) & 4$^{\ast}$ &   0.9124 (0.2660)& 0.1375 %& Irregular 2020 & Fixed 1 mm \\
%%Bullet Rosette (4) & 4$^{\ast}$  & 0.9148 (0.6892) & 0.1082 %& Irregular 2020 & Fixed 1 mm \\
%%Bullet Rosette (8) &  4$^{\ast}$  & 0.9678 (0.8666) & 0.1671 %& Irregular 2020 & Fixed 1 mm \\
%Hexagonal Plates (1) &  D74 &  --- & --- & Planar Flutter & %Hex. Plates \\
%Hexagonal Columns (1) & MA94 & --- & --- &  Columns 2022-2023 %& Hex. Columns \\
%Dendrite (1) &  ---  & 0.9588 (0.0203) & 0.1283 & Planar %Flutter & Fixed 1 mm\\
%Needle (1) &  10 & 0.0816 (0.866) & 0.4775 & Columns 2020 & %Fixed 1 mm \\
%Capped Column (3) & 0.1$^{\ast}$ caps / 4$^{\ast}$ col. & %0.5573 (0.9556)& 0.217 & Columns 2020 & Fixed 1 mm  \\
%\bottomrule
%\multicolumn{6}{c}{$^{\ast}$ Constituent crystal aspect %ratio}\\
%\multicolumn{6}{c}{$^{\ast\ast}$ SHPSG $E_{i}/F_{i}$ aspect %ratios}\\
%\end{tabular}
%\label{table:sim_tests}
%\end{table}
%\FloatBarrier

%\FloatBarrier
%\begin{table}[ht]
%\centering
%\caption{STL models used in this section.}
%\begin{tabular}{m{0.5cm} c c c c c c}
%\toprule
%\toprule
%STL Num. & Habit Type & \# Constituents & AR & Sphericity & $\varphi_{ba}$ $(\varphi_{cb})$ & Ellip. Frac. \\
%\midrule
%1& Irregular &  1 &  0.75$^{\ast\ast}$/0.75$^{\ast\ast}$ & --- & --- & --- \\
%2& Bullet Rosette & 3 & 4$^{\ast}$ &  0.5628 &   0.9124 (0.2660)& 0.1375 \\
%3& Bullet Rosette & 4 & 4$^{\ast}$ & 0.5463 & 0.9148 (0.6892) & 0.1082  \\
%4& Bullet Rosette & 8 & 4$^{\ast}$  & 0.4764 & 0.9678 (0.8666) & 0.1671 \\
%5& Plates & 1 &  0.01 & 0.1209  & 0.8660 (0.0141) & 0.4775\\
%5 & Plates & 1 & D74 & --- & --- & --- \\
%6& Dendrite & 1 &  --- & 0.1378 & 0.9588 (0.0203) & 0.1283\\
%7& Column & 1 & 0.4 & 0.6964 & 0.2041 (0.866) & 0.4775  \\
%7& Hexagonal Columns & 1 & AV70  & --- & --- & ---  \\
%8& Needle & 1 &  10 & 0.5472 & 0.0816 (0.866) & 0.4775  \\
%9 & Capped Column & 3 & 0.1$^{\ast}$ caps / 4$^{\ast}$ col. & 0.4469 & 0.5573 (0.9556)& 0.217  \\
%\bottomrule
%\multicolumn{7}{c}{$^{\ast}$ Constituent crystal %aspect ratio}\\
%\multicolumn{7}{c}{$^{\ast\ast}$ SHPSG $E_{i}/F_{i}$ aspect ratios}\\
%\end{tabular}
%\label{table:sim_tests}
%\end{table}
%\FloatBarrier



%\FloatBarrier
%\begin{table}[ht]
%\centering
%\caption{STL tests.}
%\begin{tabular}{c c c c c c c}
%\toprule
%\toprule
%Name & STL Num. & Orientation & Size Distribution  \%\
%\midrule
%Irregular 2020 & 1 & Irregular 2020 & Irregular \\
%Columns 2020 & 7 & Columns 2020 & Columns MA94\\
%Plate & 5 & Planar flutter & D74\\
%Needle & 8 & Columns 2020 & Fixed 1 mm\\
%Dendrite & 6 & Planar flutter & Fixed 1 mm \\
%Capped Column & 9 & Columns 2020 & Fixed 1 mm\\
%BR03 & 2 & Irregular 2020 & fixed 1 mm\\
%%BR04 & 3 & Irregular 2020 & fixed 1 mm \\
%BR08 & 4 & Irregular 2020& fixed 1 mm \\
%\bottomrule
%\end{tabular}
%\label{table:sim_test_dist_params}
%\end{table}
%\FloatBarrier

\FloatBarrier
\begin{figure}[ht]
\centering
\includegraphics[width=\textwidth]{STL_IMPACTS_COMP_NEW_PHI.png}
\caption{Example STL models for simulation tests along with similar cloud probe images from IMPACTS for each habit type. Irregular particle images are from the PHIPS imager whereas the others are from the CPI. Example aspect ratios ($\varphi$) values shown for plates, columns, and needles are defined as the hexagonal prism facet maximum dimension divided by the basal facet maximum dimension. ``Linear'' particles are considered to be columns and needles. Three bullet rosette examples with 3 branches (BR03), 4 branches (BR04), and 8 branches (BR08) are shown as well.}
\label{fig:stl_ex}
\end{figure}
\FloatBarrier

\subsection{Monte Carlo simulation methodology}\label{sect:MC_sim}

Monte Carlo simulations are performed using both the analytical ellipsoid formulation and the STL model methodology outlined in Sect.~\ref{methods:math} and Sect.~\ref{sect:STL}, respectively, where sizes are drawn from gamma distributions and orientations are randomly drawn from habit- and case-specific Fisher and von Mises distributions. Ellipsoids additionally have aspect ratios and area ratios drawn from the beta distributions outlined in Sect.~\ref{methods:math}. After each ellipsoid or STL particle is drawn, the 2D ellipse projection geometries are used to filter the particles in these simulations to mimic the 2D-S filtering outlined in Sect.~\ref{sect:2DS_processing}. The probability that the projection is not entirely in the buffer is determined from
\begin{equation}
P_{\mathrm{in}} = \frac{\max\left(0,W_{\mathrm{array}}-W_{\mathrm{cross}}\right)}{W_{\mathrm{array}}}
\end{equation}
where
\begin{equation}\label{W_cross}
W_{\mathrm{cross}} = 2 \sqrt{\left( r_{a} \sin \theta_{\mathrm{proj}} \right)^{2} + (r_{b} \cos \theta_{\mathrm{proj}})^{2}},
\end{equation}
$W_{\mathrm{array}}=1.28$ mm, $r_{a}$ and $r_{b}$ are the semi-major and semi-minor ellipse axes, respectively, and $\theta_{\mathrm{proj}}$ is the ellipse angle as defined in Appendix~\ref{appen:A}. Modified sigmoid functions are used to simulate a smooth transition in the filtering for the size, aspect ratio, area ratio, and pixel number thresholds. These sigmoid functions are defined as
\begin{equation}
P({x|k}) = \frac{1}{1+e^{-k \left( x-x_{\mathrm{thresh}} \right)}},
\end{equation}
where $k=50$ for the projected ellipse maximum dimension $D_{\mathrm{major}} = 2 r_{a}$, $k=50$ for the projected aspect ratio $\varphi_{\mathrm{proj}}\equiv r_{b}/r_{a}$, $k=30$ for the projected area ratio AR, and $k=0.1$ for the number of projected pixels $N_{\mathrm{pixel}}$. The thresholds ($x_{\mathrm{thresh}}$) mimic the 2D-S thresholds in Sect.~\ref{sect:2DS_processing} where $D_{\mathrm{thresh}}=0.35$ mm, $\varphi_{\mathrm{thresh}}=0.1$, $\mathrm{AR}_{\mathrm{thresh}}=0.2$, and $N_{\mathrm{thresh}}=800$. These probabilities are multiplied together for each sampled particle
\begin{equation}
P_{\mathrm{keep}} = \prod_{c \in \{xy, xz\}} \left(P_{\mathrm{in},c} \prod_{x} P_{x,c}\right),
\end{equation}
where $x \in \{D_{\mathrm{major}}, \varphi_{\mathrm{proj}},\mathrm{AR},N_{\mathrm{pixel}}\}$. A corresponding stochastic detection threshold $U$ is also drawn from a uniform distribution [$U \sim \mathrm{Uniform}(0,1)$] such that the particle is rejected if $U > P_{\mathrm{keep}}$. 

An STL particle is also instantly rejected before these probabilities are calculated if the X-Y or X-Z projection of the enclosing ellipsoid meets any of the following three criteria
\begin{enumerate}
\item The projection maximum dimension $2 r_{a}$ is less than $0.5D_{\mathrm{thresh}}$.
\item The cross-track width exceeds the array width plus a 0.5 mm tolerance: $W_{\mathrm{cross}} >W_{\mathrm{array}}+0.5\, \mathrm{mm} = 1.78\, \mathrm{mm}$
\item The maximum theoretical pixel count falls below the threshold: $N_{\mathrm{max}} < N_{\mathrm{thresh}}$,
\end{enumerate} 
where 
\begin{equation}\label{eq:N_max}
N_{\max} = \frac{\pi r_{a} r_{b}}{A_{\mathrm{pixel}}},
\end{equation}
and $A_{\mathrm{pixel}}= 0.0001\, \mathrm{mm}^{2}$ represents the 10 $\mu$m $\times$ 10 $\mu$m 2D-S pixel resolution. This best-fit ellipsoid prefiltering is included to reduce the number of computationally intensive voxel calculations used in the STL projection and corresponding CGAL ellipse fitting.

Each Monte Carlo simulation consists of randomly sampling and performing the rotations and projection calculations from the database of 1,000 fused SHPSG irregular particle or 10,000 particles from the other STL habits shown in Fig.~\ref{fig:stl_ex}. Each STL model is Monte Carlo simulated a total of 1,000,000 times. However, the synthetic 2D-S filtering procedure leads to a roughly 5\%-15\% success frequency depending upon each STL habit. This, along with any \citet{Schima2024} habit restrictions therefore leads to fewer simulated H and V counts.
 
%500,000 times for joint histograms and 200,000 times for univariate histograms. 

% STL orientations




%\FloatBarrier
%\begin{figure}[ht]
%\centering
%\includegraphics[width=\textwidth]%{IMPACTS_2DS_SHPSG_Simulation_PCGL_Irregular_2020_exp.png}
%\caption{As in \cref{fig:EllipSim_Irreg_2020} but using SHPSG-generated irregular particles.}
%\label{fig:SHPSG_irreg_PCGL_2020}
%\end{figure}
%\FloatBarrier



%For this study, we only use UIOOP particle statistics generated by the \citet{Holroyd1987} habit classification algorithm and the ellipse fit data. 




\section{NASA IMPACTS Observations}\label{sect:obs}
%\subsection{NASA IMPACTS Campaign}

Table~\ref{table:cases} provides the number of 2D-S ellipse fits that were performed on particles from various temperature regions for each case. Temperature regions are determined using the convention and conceptual model of \citet{Varcie2023} where the Polycrystalline Growth Layer (PCGL) are temperatures less than -18$^{\circ}$C, the Dendritic Growth Layer (DGL) are temperatures between $-12^{\circ}$C and $-18^{\circ}$C, the Plate Growth Layer (PGL) are temperatures between $-8^{\circ}$C and $-12^{\circ}$C, the Needle Growth Layer (NGL) are temperatures between $-3^{\circ}$C and $-8^{\circ}$C, and the aggregation layer (AGG) are temperatures between $0^{\circ}$C and $-3^{\circ}$C. Overall, there are substantially more particles/images at lower temperatures than higher temperatures where more than half of all data comes from the DGL and PCGL. Table~\ref{table:cases} also demonstrates the heterogeneous sampling from case to case. Some cases such as 20220213 and 20220219 only sampled the lowest temperature regions whereas other cases such as 20200201 primarily sampled higher temperatures. Figure~\ref{fig:temp_height} shows timeseries of 2D-S-derived habit fractions from each of the 26 cases using the \citet{Schima2024} algorithm. Irregular particles overwhelmingly dominate almost all time periods except periods in the aggregation layer (e.g., 20200201, 20200220, and 20230209) where large supercooled liquid droplets and early-onset melting can produce liquid droplets. There are also a number of NGL periods throughout where there is a notable increase in linear classified particles (e.g., 20200125). There are overall very few dendrites or hexagonal particles that the \citet{Schima2024} identifies. This, in part, is likely is due to the limited resolution of the 2D-S and the impact that orientation has on these particles; that is, dendrites and hexagonal particles would be best identified if viewed along the basal facet whereas these particles would likely be identified as linear particles if viewed along the prism facet. Furthermore, various studies such as \citet{Korolev2000}, \citet{Korolev2003}, \citet{Dunnavan2022} demonstrate that midlatitude cloud systems are dominated by irregular particles. A recent study by \citet{Garcia2026} similarly found that dendrites and hexagonal crystals were rare throughout the IMPACTS campaign and that irregular complex polycrystals dominated the DGL.  Graupel particles are identified sporadically throughout various cases with peaks in graupel habit fractions occurring in different layers. Oriented particles were rarely identified. Figure~\ref{fig:CPI_54ex} shows fifty-four example CPI 5-second panels of representative habits throughout the 26 cases shown in Fig.~\ref{fig:temp_height}. These CPI images show that these irregular particles have various forms such as bullet rosettes and side planes and potentially aggregates of such forms. Linear particles show a variety of different kinds of columns such as shown in panels 7, 13, 28 and 36. Needles also appear in panels 2, 33, 39, 46, 47, 48, 49 and 50. Capped columns occasionally appear such as demonstrated in panels 3 and 25 and these linear particles are sometimes rimed as shown in panels 5 and 37. Pristine dendrites, while rare, do appear with different stellar structures as shown in panels 1, 11, 19, 32, 35, and 38. More complicated and irregular-looking dendrites and spatial planar particles also appear in these regions as shown in panels 34 and 52. Hexagonal particles often are not pristine as shown in panels 17 and 30; the \citet{Schima2024} algorithm will often characterize these more irregular-like particles as hexagonals or dendrites if viewed along the basal facet of one of the branches. Two different types of graupel periods are shown in panels 24 and 26. Panel 24 shows heavily rimed, quasi-spherical particles whereas panel 26 shows lighter rimed particles with some rimed columns. Therefore, it would seem that periods with graupel particles might or might not exhibit low aspect ratios due to whether columns are present.

\FloatBarrier
\begin{sidewaystable}[ht]
    \centering
        \caption{List of IMPACTS cases used in this study and the total number of horizontal and vertical ellipse fits used in different temperature ranges. Vertical channel 2D-S ellipse fit counts are in parentheses whereas horizontal channel counts are not. }
%\resizebox{\textheight}{!}{%
  %  \begin{tabular}{>{\centering\arraybackslash}m{1.1cm} >{\centering\arraybackslash}m{2.3cm} >{\centering\arraybackslash}m{2.3cm} >{\centering\arraybackslash}m{2.3cm} >{\centering\arraybackslash}m{2.3cm} >{\centering\arraybackslash}m{2.3cm}}
    
   \begin{tabular}{c c c c c c}
    \toprule 
    \toprule
 %       &   \multicolumn{5}{c}{\large Number of Horizontal (Vertical) Channel Particle Images}\tabularnewline
%        \midrule
          Date    &  AGG&  NGL  & PGL & DGL & PCGL\tabularnewline
          \midrule
        01/18/2020  & 26,148 (26,478)  & 67,419 (73,677)  & 69,422 (69,293) & 76,579 (77,659) &  37,781 (38,229) \tabularnewline
        01/25/2020  & 4,995 (4,954)  & 81,058 (82,210) & 16,852 (17,220)  & 78,143 (78,235)& 214,290 (219,425) \tabularnewline
        02/01/2020  & 31,417 (32,838)  & 47,324 (49,987) & 69,749 (70,990) & 13,445 (13,665) & 1,168 (1,141) \tabularnewline
        02/05/2020  & 23,864 (24,571)  & 38,807 (39,592) & 89,908 (90,736)  & 99,466 (101,214) & 50,630 (52,445) \tabularnewline
        02/07/2020 & 22,768 (22,939) & 79,219 (78,894) & 38,992 (39,552) & 116,599 (116,067)&  81,171 (83,575) \tabularnewline
        02/13/2020 & 17,170 (17,274) & 28,938 (29,917) & 9,919 (10,190) & 92,513 (93,864) & 4,487 (4,465)  \tabularnewline
        02/18/2020 & 10,784 (10,713)  & 48,077 (48,914) & 25,734 (25,542) & 33,272 (34,335) & 96,247 (98,626) \tabularnewline
        02/20/2020 & 11,086 (11,505)  & 27,255 (26,974) & 45,129 (44,364) & 16,957 (16,593) & 267,728 (272,820)  \tabularnewline
        02/08/2022 & 202 (202)  & 51,087 (50,633) & 12,721 (12,175) & 32,344 (31,239) &  4,875 (4,626) \tabularnewline
        02/13/2022 & 0 (0)  & 0 (0) & 0 (0) & 8,718 (8,669) & 34,668 (34,388) \tabularnewline
        02/17/2022 & 2 (5,780)  & 1,694 (8,881) & 16,947 (16,912)  & 98,173 (100,645) &  266,905 (269,369) \tabularnewline
        02/19/2022 & 0 (0)  & 0 (0) & 0 (0)  & 128,292 (149,362) & 84,468 (108,727) \tabularnewline
        02/25/2022 & 8,533 (8,905) & 22,299 (22,693) & 4,278 (4,390)  & 8,129 (8,242) &  62,116 (64,644) \tabularnewline
        01/13/2023 & 7,890 (7,300)  & 54,158 (54,353) & 49,758 (49,872) & 116,599 (109,114) &  70,663 (69,866) \tabularnewline
        01/15/2023 & 2 (5) & 67,023 (70,667) & 204,338 (218,536) & 68,316 (71,066) & 122,288 (130,399) \tabularnewline
        01/19/2023 & 12,669 (12,804) & 29,624 (29,451) & 83,145 (84,904) & 94,291 (96,691) &  20,835 (21,240)  \tabularnewline
        01/23/2023 & 7,472 (7,751) & 55,433 (56,875) & 62,516 (63,393) & 145,406 (150,590) & 15,116 (15,232)  \tabularnewline
        01/25/2023 & 4,158 (4,051)  & 46,692 (49,125) & 13,512 (13,733) & 86,008 (89,358) & 58,486 (59,502) \tabularnewline
        01/29/2023 & 7,415 (7,460)  & 13,195 (13,474) & 9,120 (9,041) & 82,439 (83,558) & 107,838 (110,401)  \tabularnewline
        02/05/2023 & 0 (0) & 5,382 (5,770) & 9,962 (10,049) & 12,161 (12,903) &  113,197 (116,051) \tabularnewline
        02/09/2023 & 9,336 (9,509) & 77,162 (78,426) & 8,096 (8,337)  & 69,896 (72,046) & 33,898 (34,996)  \tabularnewline
        02/12/2023 & 48,876 (49,496) & 83,623 (84,206) & 71,290 (73,249) & 290,067 (298,069) & 65,845 (68,248)  \tabularnewline
        02/14/2023 & 0 (0) & 26,779 (27,178) & 53,465 (54,344) & 259,414 (265,656) & 177,678 (182,187) \tabularnewline
        02/16/2023 & 9,941 (10,382) & 19,298 (20,349) & 47,043 (47,983) & 191,976 (201,284) & 159,196 (169,656)  \tabularnewline
        02/23/2023 & 0 (0) & 4,328 (4,483) & 38,970 (40,147) & 74,273 (76,097) &  2,554 (2,692) \tabularnewline
        02/28/2023 & 6,002 (6,331) & 56,689 (57,739) & 4,849 (5,282) & 313 (328) & 5,824 (5,944)   \tabularnewline
       \midrule 
       \textbf{Total} & \textbf{270,730 (281,248)}  & \textbf{1,032,563 (1,064,468)}  & \textbf{1,055,715 (1,080,234)}  & \textbf{2,293,789 (2,356,549)} & \textbf{2,159,952 (2,238,894)} \tabularnewline
       \midrule 
       $\mathbf{T<\, 0\, ^{\circ}\mathrm{C}}:$& \multicolumn{5}{c}{ \textbf{6,812,749 (7,021,393)}}\tabularnewline
       \bottomrule
 \end{tabular}
    \label{table:cases}
\end{sidewaystable}
\FloatBarrier

%\textbf{270,730}\par \textbf{(281,248)} 

Figure~\ref{fig:2DS_habit_examples} shows examples for how each \citet{Schima2024} habit generally appears as 2D-S images in the horizontal and vertical channels along with the UIOOPS CGAL ellipse fits. Irregular particles are generally small and often around $0.5$ mm in maximum dimensions. Dendrites often represent larger irregular types like side planes or are indeed  dendrites. Hexagonal particles often appear hollow due to optical transmittance if hexagonal prisms are viewed along their basal facet. Some hexagonal particles are more pristine with more intricate patterns whereas some have little to no branching. Graupel particles appear much like irregular particles but generally appear more circular with perimeters that have noticeable quasi-periodic irregularities. Aggregates observed by the 2D-S are often linear chain aggregates as those particles can more easily stay fully within the 2D-S buffer strips if they are oriented horizontally (i.e., with ellipse angles around $0^{\circ}$). Spheres are characterized by high aspect ratios with area ratios close to unity.

\FloatBarrier
\begin{figure}[ht]
    \centering
    \includegraphics[width=\textwidth]{Figures/IMPACTS_2DS_Habit_Fraction_26flights_56ex.pdf}
    \caption{Time series of habit fractions for the 2D-S (vertical channel) during each of the 26 IMPACTS flights. Shaded regions correspond to colors representing each habit growth region based on the temperature ranges outlined in Table~\ref{table:cases} for various 2D-S habits: PCGL (irregular; black), DGL (dendrite; blue), PGL (hexagonal; teal), NDL (linear; orange) and AGG (aggregate; green). Regions where  $T_{c}>0\, ^{\circ}\mathrm{C}$ are shown with the color of the ``sphere''  category (dark red). Temperature regions are smoothed with a 20-minute moving average to better highlight various regions.}
    \label{fig:temp_height}
\end{figure}
\FloatBarrier



\FloatBarrier
\begin{figure}[ht]
   \centering
    \includegraphics[width=0.85\textwidth]{Figures/CPI_54ex_IMPACTS_9x6.png}
    \caption{Cloud Particle Imager (CPI) 5-second image collage from each of the 54 times shown in Fig.~\ref{fig:temp_height}.}
    \label{fig:CPI_54ex}
\end{figure}

\FloatBarrier

\FloatBarrier
\begin{figure}[ht]
   \centering
    \includegraphics[width=\textwidth]{Figures/2DS_example_images_habits.pdf}
    \caption{Example 2D-S images for various habits identified by the \citet{Schima2024} 2D-S habit classification algorithm.}
    \label{fig:2DS_habit_examples}
\end{figure}
\FloatBarrier

Figure~\ref{fig:habit_heatmaps} shows the fraction of each \citet{Schima2024} identified habit from each of the five temperature layers. Most habit classifications are consistent with their respective temperature layers. Most irregular particles were classified in the PCGL and DGL whereas dendrites and linear particles were predominantly identified in the DGL and NGL, respectively. Other habits such as spheres appear mostly at lower temperatures, probably due to the larger number concentration of small particles at those temperatures compared to the lower number concentrations of supercooled liquid droplets in the AGG layer. Aggregates are identified mostly in the DGL, PGL, and NGL rather than the PCGL or AGG layers. The habit fractions are overall very  consistent between both H and V channels.

\FloatBarrier
\begin{figure}[ht]
\centering
\includegraphics[width=\textwidth]{IMPACTS_all_habit_layer_heatmap.pdf}
\caption{Heatmaps of \citet{Schima2024} habit classifications from the horizontal (top) and vertical (bottom) 2D-S channels for all 26 IMPACTS cases within each of the temperature layer shown in Table~\ref{table:cases}. Numbers in each cell represent the total number of images in the layer and \citet{Schima2024} habit category. Colors indicate the fraction of each habit in each layer such that the summation of all habit counts in each row for a column gives the total number of particles for that column (shown below AGG cells).}
\label{fig:habit_heatmaps}
\end{figure}
\FloatBarrier

Most cases individually illustrate a characteristic distribution of ellipse fit parameter values. This can be seen in Fig.~\ref{fig:boxplots} which shows the ellipse aspect ratios, orientations, and area ratios for all H and V channel ellipse data from each individual case. Most cases show median ellipse aspect ratios and area ratios around 0.7. The ellipse orientation angles are extremely similar to one another, but only for the 2020 cases and 2022-2023 cases, separately. PCGL irregular particles, which represent nearly 2 million samples in each channel as shown in Fig.~\ref{fig:habit_heatmaps}, in particular exhibit a remarkably stable set of ellipse parameter distributions among all flights. Two-sample Kolmogorov–Smirnov tests comparing ellipse aspect ratios, orientations, and area ratios of \citet{Schima2024} PCGL irregular particles between the pooled 2020 ensemble and each of the eight individual flights yield median deviations of 0.0319 (0.0275), 0.0236 (0.0269), and 0.0575 (0.0666) for the H (V) channel, respectively. Corresponding median deviations for the eighteen 2022–2023 flights evaluated against their respective campaign ensemble are 0.0347 (0.0407), 0.0327 (0.0409), and 0.0616 (0.0543). This suggests that the morphological characteristics from the variety of observed irregular particles in the PCGL could possibly be accurately represented using only one set of shape and density distributions. Section~\ref{sect:simulations} tests this hypothesis. Some cases do show more prominent deviations such as 20220208, 20220225, and 20230228 where the interquartile range (IQR) of aspect ratios are generally larger and the median aspect ratios are lower. These cases are generally warmer as seen in Table~\ref{table:cases} and Fig.~\ref{fig:temp_height} which suggests the prevalence of low aspect ratio habits such as columns and plates. In particular, the vast majority of 20230228 particles come from the NGL (Table~\ref{table:cases} and Fig.~\ref{fig:temp_height}) which suggests that many of these particles are most likely needles or columns.

\FloatBarrier
\begin{figure}[ht]
\centering
\includegraphics[width=\textwidth]{IMPACTS_boxplots_ellipse_params_all_enlarged.pdf}
\caption{Boxplots of (left) ellipse aspect ratios, (middle) ellipse orientation angles, and (right) ellipse area ratios for all 26 cases and for both H (blue) and V (orange) channels. Whiskers represent $\pm$ 1.5 IQR from the 75th and 25th percentiles, respectively.}
\label{fig:boxplots}
\end{figure}
\FloatBarrier

%\section{In situ ellipse fit results}\label{sect:insitu_results}

To account for the diverging orientation distributions from the 2020 cases to the 2022-2023 cases, Figures~\ref{fig:ellipse_fit_joint2_2020} and~\ref{fig:ellipse_fit_joint2_2022_2023} show joint distributions for ellipse aspect ratios, ellipse major dimensions, ellipse orientation angles, and ellipse area ratios for all cases during the 2020 (Fig.~\ref{fig:ellipse_fit_joint2_2020}) and the combined 2022\footnote{It is worth noting that the 2022 data is generally of lower quality due to occasional optical fogging and pressurization issues after the 2D-S was re-installed in 2021. These instrumental issues were resolved before the 2023 flights.} and 2023 campaigns (Fig.~\ref{fig:ellipse_fit_joint2_2022_2023}). The distributions among all cases show nearly identical 2D ellipse aspect ratios, area ratios, and major dimensions in both the H and V channels. There are also at least four separate regions of particle habits which shows the ellipse characteristics of irregular, linear, aggregate, and spherical particles. The vast majority of irregular particles, which dominate all particles as demonstrated in Fig.~\ref{fig:temp_height}, have ellipse aspect ratios between about 0.6 and 0.9 with a mode around 0.75 in both channels. Similarly, the ellipse area ratios for irregular particles concentrate between 0.65 and 0.7 and major ellipse dimensions are between 0.4 mm and 0.7 mm. The irregular particle orientation angles however show a remarkable difference between 2020 and 2022-2023 campaigns. While the irregular ellipse orientations in 2020 are roughly symmetric with mean orientations at -45$^{\circ}$ and $+45^{\circ}$ from the H and V channels, respectively, the 2022 and 2023 vertical channels instead show a mean orientation angle closer to $0^{\circ}$ with some ellipses exhibiting slightly negative orientations. The stark difference between the 2020 and 2022-2023 distributions and the extraordinarily similar orientation distributions between the 2022 and 2023 are quite surprising and difficult to explain on first glance. However, it appears that these differences can be, at least in part, explained by a mislabeling and rotation of the 2D-S channels either before or after the 2020 deployments as well as a 3 in. displacement and 2.8 degree rotation in the 2D-S location between campaign years. This is described in more detail in Appendix B.

% The difference in the equilibrium orientation angle location between the 2020 and 2022-2023 cases appears to be the result of a mislabeling and rotation of the 2D-S channels that occurred after the 2020 campaign. \cref{sect:discuss_King} provides compelling evidence to support this claim.

Linear particles are characterized in both the H and V channel distributions by a power-law-like ``tail'' of ellipse aspect ratio with increasing ellipse major axis lengths. These tails are apparent for 2020, 2022, and 2023 and appear generally the same in the H and V channel. The area ratios of these linear are similar to that of irregular particles as seen in the area ratio and aspect ratio joint distributions. Like the ellipse aspect ratio and ellipse major dimension joint distributions, the area ratio/aspect ratio joint distributions are also consistent between the H and V channels. The linear particle orientation angles show up most clearly in the 2020 vertical channel and the 2022-2023 horizontal channels where the joint aspect ratio/orientation distributions show notable dips. The orientations associated with these linear particles in the opposite channels are generally more disperse but seem to be generally more concentrated around $\pm 90$ rather than around $0^{\circ}$. It is worth noting that previous studies such as \citet{King1985c,King1985d} and \citet{Obrien2016} have found that columns, when viewed from the vertical channel, have a preferential orientation of about $60^{\circ}$. This preferred orientation theoretically results as an equilibrium due to the net torque exerted on the particles due to the aircraft vortices generated near wing tips \citep{King1985d}. The analytical model proposed by  \citet{King1985d} suggests that the location of the probe under the wing can change this expected equilibrium angle increases with probe placements that are farther in front of the wing and closer to the aircraft. The slightly larger equilibrium orientation of $+70^{\circ}$ shown in Fig.~\ref{fig:ellipse_fit_joint2_2020} and $-70^{\circ}$ in Fig.~\ref{fig:ellipse_fit_joint2_2022_2023} can be explained in terms of the extended pylon used during the IMPACTS campaign. 

The ellipse aspect ratios and area ratios show an apparent slight size dependence that is most obvious for major axis lengths larger than 1.0 mm. 2D-S images indicate that these particles are largely aggregates or large spatial crystals. These larger particles exhibit near $0^{\circ}$ orientation in both the H and V channels. Simulations in Sect.~\ref{sect:simulations} confirm that these orientations are the result of the 2D-S probe processing and limitations as the maximum size along the photodiode dimension (i.e., the z and y axes in the H and V channels, respectively, see Fig.~\ref{fig:P3}b) is only 128 elements (1.28 mm) long. Therefore, any particle with a 2D projected ellipse maximum dimension as large or larger than this dimension will necessarily have to have low orientation angles in order for the particle to stay fully inside the buffer strips shown in Fig.~\ref{fig:2DS_habit_examples}. As a result, the joint distributions for the largest particles are predominantly millimeter-sized linear chain aggregates. 

Spherical particles appear in these joint distributions in small parameter regions. This is most apparent in the joint aspect ratio/area ratio figures where these supercooled liquid droplets and/or raindrops are concentrated in a somewhat circular region centered on an ellipse aspect ratio of about 0.8 with area ratios between 0.9 and 1.0. These droplets also appear distinct in the area ratio/ellipse angle joint distributions where it is clear that these particles are nearly horizontally oriented in both the H and V channels.

\FloatBarrier
\begin{figure}[ht]
\centering
%\includegraphics[width=\textwidth]{Figures/IMPACTS_2DS_2020_6panel.png}
\includegraphics[width=\textwidth]{Figures/IMPACTS_2DS_MultiFlight_2020_All_JointDists_3x2_enlarged.pdf}
\caption{Joint ellipse fit parameter distributions for all eight 2020 cases from Table~\ref{table:cases} involving ellipse aspect ratios and major dimension (top left), ellipse aspect ratio and angle (top right), ellipse area ratio and major dimension (middle left), ellipse area ratio and angle (middle right), ellipse aspect ratio and area ratio (bottom left), and ellipse major dimension and angle (bottom right).  }
\label{fig:ellipse_fit_joint2_2020}
\end{figure}
\FloatBarrier

\FloatBarrier
\begin{figure}[ht]
\centering
%\includegraphics[width=\textwidth]{Figures/IMPACTS_2DS_2022_joint_distributions_all.png}
\includegraphics[width=\textwidth]{Figures/IMPACTS_2DS_MultiFlight_2022_2023_All_JointDists_3x2_enlarged.pdf}
\caption{As in Fig.~\ref{fig:ellipse_fit_joint2_2020} but for all 2022 and 2023 cases (18 in total; 5 from 2022 and 13 from 2023).}
\label{fig:ellipse_fit_joint2_2022_2023}
\end{figure}
\FloatBarrier

%\FloatBarrier
%\begin{figure}[ht]
%\centering
%\includegraphics[width=\textwidth]{Figures/IMPACTS_2DS_2023_joint_distributions_all.png}
%\includegraphics[width=\textwidth]{Figures/IMPACTS_2DS_MultiFlight_2023_All_JointDists_3x2.pdf}
%\caption{As in \cref{fig:ellipse_fit_joint2_2020} but for all 13 2023 cases.}
%\label{fig:ellipse_fit_test}
%\end{figure}
%\FloatBarrier

% 20220225, 20230115, 20230228 show potential rimed particles with angles ~ 0 deg.

%Because there are only 5 cases where the 2D-S was operational during the 2022 campaign, the vast majority of these particles seem to come from only one case: 20220225. This particular case is somewhat of an outlier as most of the ellipse fit orientations were approximately $0^{\circ}$ in both 2D-S channels.

One notable difference between the 2020 and 2022-2023 joint distributions is the presence of low aspect ratio particles with near $0^{\circ}$ orientations in both the horizontal and vertical planes. While these low aspect ratio ellipses with near $0^{\circ}$ orientation occur sporadically from case-to-case, they are most prominent in only a select number of cases including: 20220225, 20230115, and 20230228. CPI images (for example, CPI panels 26 and 31 in Fig.~\ref{fig:CPI_54ex}) and the habit identification of spheres and graupel in the Fig.~\ref{fig:temp_height} timeseries plots suggest that there was riming throughout these case with many of the rimed particles seemingly columnar. Thus, it appears that graupel particles could possibly be identified from the near $0^{\circ}$ ellipse orientations in either 2D-S channel. In general, these low aspect ratio particles corresponding to ``linear'' particles (i.e., columns and needles) as well as rimed columns can be identified from the joint distributions  of ellipse parameters in the five different temperature levels shown in Table~\ref{table:cases}. This temperature dependence is shown in Fig.~\ref{fig:ar_angle_2020_temp} for the 2022-2023 aspect ratio/orientation joint distribution. As temperature increases, the irregular particle population from the PCGL slowly becomes more disperse and lower aspect ratio particles become more prominent. The lower aspect ratios corresponding to unrimed linear particles have counts concentrated at $-70^{\circ}$ in the horizontal channel and a broader distribution of $\pm 90^{\circ}$ orientation angles in the vertical channel. These lower aspect ratio particles appear most prominently in the NGL which suggests that these linear particles are likely columns or needles. The aggregation layer shows the appearance of supercooled liquid droplets with high aspect ratios around 0.8 to 0.9 and orientations close $0^{\circ}$ in both H and V channels. 

%Interestingly enough, the NASA ER-2 airborne triple-frequency radar analysis performed by \citet{Dunnavan2023} for this case (performed hours earlier than the P-3 flight legs) suggested that the snow for this particular case was dominated by needle aggregates, possibly resulting from rime splintering. Despite the temporal and spatial differences between the ER-2 airborne radar measurements and the P-3 measurements, it is reasonable to speculate that this particular case had prevalent riming.

%; the joint distributions shown in  \cref{fig:ellipse_fit_joint2_2020} do indeed suggest that these parameters are apparently universal even for the individual cases themselves rather than a smoothing that results from mixing disparate distributions






%The 2020 cases similarly show these temperature dependencies as well and some cases such as 20220208 and 20230228 have a more pronounced transition between irregular particles and linear particles.

\FloatBarrier
\begin{figure}[ht]
\centering
%\includegraphics[width=0.75\textwidth]{IMPACTS_2DS_AspectRatio_EllipseAngle_2020_temp_levels.png}
\includegraphics[width=0.7\textwidth]{IMPACTS_2DS_AspectRatio_EllipseAngle_2022_2023_All.pdf}
\caption{Joint distributions of ellipse aspect ratios and orientations for the five temperature regions shown in Table~\ref{table:cases} for all eighteen 2022 and 2023 cases. The colorbar for each panel is identical to those shown in Figs.~\ref{fig:ellipse_fit_joint2_2020} and~\ref{fig:ellipse_fit_joint2_2022_2023}. }
\label{fig:ar_angle_2020_temp}
\end{figure}
\FloatBarrier

\section{Simulations}\label{sect:simulations}

Table~\ref{table:ellip_orient_tests} shows the simulation tests used in this section along with their various size, shape, density, and orientation distribution parameters. Note that the simulations are generally split into all eight 2020 cases and all eighteen 2022 and 2023 cases. As shown in Sect.~\ref{sect:obs}, the 2020 and 2022-2023 cases demonstrate separate orientation distributions which are likely due to the 2D-S being installed differently in 2020 than in 2022 and 2023  (Appendix B provides both detailed CFD modeling and observational evidence that supports this conclusion). To investigate plausible morphological and orientation characteristics for the 2020 and 2022-2023 cases, we perform five ellipsoid-based Monte Carlo simulations and a single SHPSG-based STL simulation for irregular particles. We also show two sets of ellipsoid and STL simulations for linear particles for the 2020 and the 2022-2023 flights.  As shown in Table~\ref{table:ellip_orient_tests}, we used the 2020 and 2022-2023 PCGL \citet{Schima2024} identified irregular particle mean orientation angles to estimate $\overline{\phi}_{p}$ and $\overline{\theta}_{p}$ from $\overline{\theta}_{\mathrm{proj},h}$ and $\overline{\theta}_{\mathrm{proj},v}$. Similarly, $\overline{\theta}_{\mathrm{proj},h}$ and $\overline{\theta}_{\mathrm{proj},v}$ are used to estimate the orientation parameters for the ``Linear'' and ``Capped Column'' simulations where we assume that these habits share the same orientation parameters. We assume that the bullet rosettes from the 2020 cases (``Bullet Rosettes 2020'') will assume the same orientations as those inferred for the ``Irregular 2020'' simulations. The ``Plates'' and ``Dendrite'' simulations use random orientations as these habits are sparser in the IMPACTS dataset (Fig.~\ref{fig:habit_heatmaps}) and are more difficult to reliably identify using the \citet{Schima2024} habit classification algorithm. As such, we use random orientations for these two planer crystal habit types as a stress test for investigating their potential H and V channel projection properties. To prevent the generation of unrealistically small dendrites, the dendrite simulation only considers a single 1 mm STL dendrite rather than drawing sizes from a gamma distribution.

%We perform only one Monte Carlo simulation for each habit type (irregular, linear, dendrite, plates, capped columns, and bullet rosettes) because the STL simulations are much more computationally demanding than the ellipsoid simulations.

\FloatBarrier
\begin{table}[ht]
\centering
\caption{Size, shape, density, and orientation distribution parameters used for ellipsoid and/or STL simulations. Random (i.e., chaotic) orientations use $\kappa=\kappa_{\psi} = 0.001$. $D_{\mathrm{mean}}$ is in units of millimeters.}
\resizebox{\textwidth}{!}{%
\begin{tabular}{c c c c c c c c c c c c c c c c c c c c}
\toprule
\toprule
Simulation Name & $\nu$ & $D_{\mathrm{mean}}$ & $\overline{\varphi}_{ba}$ & $\overline{\varphi}_{cb}$ & $\zeta$ & $\overline{\mathrm{AR}}$ & $\varrho_{\mathrm{frac}}$ & $\overline{\theta}_{\mathrm{proj},h}$ & $\overline{\theta}_{\mathrm{proj},v}$ &$\overline{\phi}_{p}$ & $\overline{\theta}_{p}$ &$\kappa$ & $\overline{\psi}_{p}$ & $\kappa_{\psi}$ & Frac. & Ellipsoids/STL \\
\midrule
Irregular 2020 & \multirow{4}{*}{1} & \multirow{4}{*}{2.0} & \multirow{4}{*}{0.75} & \multirow{4}{*}{0.75} & \multirow{4}{*}{0.6} & \multirow{4}{*}{0.59} & \multirow{4}{*}{0.02}  &  \multirow{4}{*}{$-40^{\circ}$} & \multirow{4}{*}{$45^{\circ}$} &  \multirow{4}{*}{$45^{\circ}$} & \multirow{4}{*}{$120.68^{\circ}$} &  \multirow{4}{*}{$4.5$} & $0^{\circ}$ & $0.001$ & \multirow{4}{*}{---} & Both \\
Irregular 2020 Parallel &   &  &  &  &  &  &   &  & &  & & & $0^{\circ}$ & $10$ & & Ellipsoids\\
Irregular 2020 Skewed &   &  &  &  &  &  &   &  & & & &  & $45^{\circ}$ & $10$ & & Ellipsoids\\
Irregular 2020 Orthogonal &   &  &  &  &  &  &   &  & & & &  & $90^{\circ}$  & $10$ & & Ellipsoids\\
\midrule
\multirow{2}{*}{Irregular 2022-2023} & $1$ & $0.5$ & $0.78$ & $0.78$ & $0.6$ & \multirow{2}{*}{$0.59$} & \multirow{2}{*}{$0.02$} & $-40^{\circ}$ & $-10^{\circ}$ &  $170^{\circ}$ & $50.43^{\circ}$ &  $3.5$ & \multirow{2}{*}{$0^{\circ}$} & \multirow{2}{*}{$0.001$} & $0.75$ & \multirow{2}{*}{Ellipsoids}\\
 & $4$ & $0.8$ & $0.65$ & $0.95$ & $1$ & &  & $-70^{\circ}$ & $-75^{\circ}$ &  $105^{\circ}$ & $54.58^{\circ}$ &  $10$ &  &  & $0.25$ & \\
\midrule
Linear 2020 & $3$ & $0.75$ & Eq.~\ref{eq:power-law_COLUMNS} & $0.95$ & $0.5$ & $0.65$ & $0.03$ & $-70^{\circ}$ & $80^{\circ}$ &  $80^{\circ}$ & $115.51^{\circ}$ &  $8$ & $0^{\circ}$ & $0.001$ & --- & Both\\
Linear 2022-2023 & $3$ & $0.75$ & Eq.~\ref{eq:power-law_COLUMNS} & $0.95$ & $0.5$ & $0.65$ & $0.03$ & $-80^{\circ}$ & $-80^{\circ}$ &  $100^{\circ}$ & $45.44^{\circ}$ &  $8$ & $0^{\circ}$ & $0.001$ & --- & Both\\
Dendrite &  \multicolumn{2}{c}{Fixed 1mm}  & --- & --- & --- & --- & --- & --- & --- & --- & --- & $0.001$ & --- & $0.001$ & --- & STL\\
Plates & $1$ & $2.0$ & --- & Eq.~\ref{eq:power-law_PLATES} & --- & --- & --- & --- & --- & --- &  --- & $0.001$ & --- & $0.001$ & --- & STL\\
Capped Columns 2020 & $1$ & $2.0$ & --- & --- & --- & --- & --- & $-70^{\circ}$ & $80^{\circ}$ &  $80^{\circ}$ & $115.51^{\circ}$ &  8 & $0^{\circ}$ & $0.001$ & --- & STL\\
Capped Columns 2022-2023 & $1$ & $2.0$ & --- & --- & --- & ---  & --- & $-80^{\circ}$ & $-80^{\circ}$ & $100^{\circ}$& $45.44^{\circ}$ & 8 & $0^{\circ}$ & $0.001$ & --- & STL\\
Bullet Rosettes 2020 & $1$ & $2.0$ & --- & --- & ---  &  --- & --- & $-40^{\circ}$ & $45^{\circ}$ &  $45^{\circ}$ & $120.68^{\circ}$ &  $4.5$ & $0^{\circ}$ & $0.001$ & --- & STL\\
%Plates & & & & & & & & & & & --- & --- & 0$^{\circ}$ & $90^{\circ}$ & 6 & $0^{\circ}$ & 6 & --- \\
%Irregular 2022-2023 POP 2 & & & & & & & & & & & -70$^{\circ}$ & -75$^{\circ}$ &  105$^{\circ}$ & $54.58^{\circ}$ &  10 & $0^{\circ}$ & 0.001 & 0.25\\
%Linear 2020 & & & & & & & & & & & -70$^{\circ}$ & 80$^{\circ}$ &  80$^{\circ}$ & $115.506^{\circ}$ &  8 & $0^{\circ}$ & 0.001 & ---\\
%Linear 2022-2023 & & & & & & & & & & & -80$^{\circ}$ & -80$^{\circ}$ &  100$^{\circ}$ & $45.44^{\circ}$ &  8 & $0^{\circ}$ & 0.001 & ---\\
%Irregular Parallel & & & & & & & & & & $\overline{\psi}_{p}$ & -35$^{\circ}$ & 45$^{\circ}$ &  45$^{\circ}$ & $116.341^{\circ}$ &  4 & $0^{\circ}$ & 10 & ---\\
%Irregular Skewed & & & & & & & & & & & $\overline{\psi}_{p}$ & -35$^{\circ}$ & 45$^{\circ}$ &  45$^{\circ}$ & $116.341^{\circ}$ &  4 & $45^{\circ}$ & 10& ---\\
%Irregular Orthogonal & & & & & & & & & & $\overline{\psi}_{p}$ & -35$^{\circ}$ & 45$^{\circ}$ &  45$^{\circ}$ & $116.341^{\circ}$ &  4 & $90^{\circ}$ & 10 & ---\\
\bottomrule
\end{tabular}%
}
\label{table:ellip_orient_tests}
\end{table}
\FloatBarrier

%\subsection{Idealized triaxial ellipsoids}\label{sim:ellipsoids}

%The purpose of this section is to understand the nature of the ellipsoid projection problem for the proposed analytical distributions given in \cref{methods:math}. We also answer the question as to what ellipsoid characteristics can and can not explain the joint distribution characteristics between H and V channels as shown in \cref{fig:ellipse_fit_joint2_2020,fig:ellipse_fit_joint2_2022,fig:ellipse_fit_joint2_2022,fig:ellipse_fit_test} for the \citet{Schima2024} habit types.
% \cref{table:ellip_dist_params} provides the ellipsoid parameters used to represent various habits.

%%The tests given in the next two sections can be easily replicated using the ``MultipleEllipsoidProjections'' Matlab App provided in S1.

%1.) What triaxial ellipsoid characteristics \textit{could plausibly} explain the observed joint distributions shown in \cref{sect:insitu_results} and 2.) What triaxial ellipsoid characteristics \textit{can not plausibly} explain the observed joint distributions shown in \cref{sect:insitu_results}. As shown in \cref{fig:ellipse_fit_joint2_2020,fig:ellipse_fit_joint2_2022,fig:ellipse_fit_test}, there are certain parameters that simply would not explain the observed distributions. For example, the observed non-uniform ellipse angle distributions rules out the possibility that the particles are randomly oriented as is sometimes assumed for in situ cloud probes. This issue of simulating plausible ellipsoid characteristics is quite challenging considering the vast parameter space of the analytical ellipsoid model described in \cref{methods:math}. As such, we utilize \cref{eq:ret_orient} to determine plausible sets of orientation parameters for different habits. 


%\FloatBarrier
%\begin{table}[ht]
%\centering
%\caption{Ellipsoid orientation tests.}
%\begin{tabular}{c c c c}
%\toprule
%\toprule
%Variable & PCGL Irregular 2020 & NGL Linear 2020 & DGL Dendrite 2020\\
%\midrule
%$\nu$ & $8$ & $4$ & $6$ \\
%$D_{\mathrm{mean}}$ & 0.6 mm& 0.7 mm & 1.0 mm  \\
%$\overline{\varphi}_{ba}$ & 0.8 & --- & --- \\
%$\overline{\varphi}_{cb}$ & 0.8 & --- & --- \\
%$\zeta$ & 0.5 & --- & --- \\
%$\theta_{\mathrm{proj},h}$ & -35$^{\circ}$ & %-70$^{\circ}$ & --- \\
%$\theta_{\mathrm{proj},v}$ & 45$^{\circ}$  & %80$^{\circ}$ & ---\\
%$\overline{\phi}_{p}$ & 45$^{\circ}$ & %80$^{\circ}$ & $0^{\circ}$ \\
%$\overline{\theta}_{p}$ & $116.341^{\circ}$ & %$115.506^{\circ}$ & $90^{\circ}$ \\
%$\kappa$ & 5 & 3 & 3 \\
%$\overline{\psi}_{p}$ & $0^{\circ}$ & $0^{\circ}$ %& $0^{\circ}$  \\
%$\kappa_{\psi}$ & 0.001 & 0.001 &  \\
%$\overline{\rho}$ & 0.8 & 0.75 & \\
%$\varrho_{\mathrm{frac}}$ & 0.12 & 0.05 &  \\
%$a_{\varphi}$ & 2 & --- & \\
%$a_{\mathrm{\rho}}$ & 4/3 & 4 &  \\
%$\alpha_{c}$ & --- & 0.1673 & \\
%$\beta_{c}$ & --- & 0.4 & \\
%\bottomrule
%\end{tabular}
%\label{table:ellip_irreg_params}
%\end{table}
%\FloatBarrier

Figures~\ref{fig:EllipSim_Irreg_2020} and~\ref{fig:SHPSG_Irreg_2020} compare the observed ellipse fit parameters of irregular-identified particles from the 2020 PCGL data with the irregular ellipsoid and STL tests described in Table~\ref{table:ellip_orient_tests}.
 As shown in Fig.~\ref{fig:EllipSim_Irreg_2020}, the triaxial ellipsoid parameters for the ``Irregular 2020'' simulation in Table~\ref{table:ellip_orient_tests} can very accurately simulate the characteristic features for each joint distribution of observed ellipse sizes, aspect ratios, area ratios, and orientations compared to the 2020 PCGL Irregular cases. In particular, the Fisher orientations from both the ellipsoid distributions and STL models do a remarkable job explaining the observed mirrored H and V channel ellipse angle distributions where $\theta_{\mathrm{proj},h} = -40^{\circ}$ and $\theta_{\mathrm{proj},v} = 45^{\circ}$ (see Table~\ref{table:ellip_orient_tests}) that are present in all eight of the 2020 cases (see Fig.~\ref{fig:boxplots}). The bivariate beta aspect ratio distributions and quasi-random SHPSG-generated STL models also perfectly explain the beta-like ellipse aspect ratio distributions due to the quasi-random spatial configurations present in various irregular particles. Furthermore, the pronounced transition of the symmetric orientation angles at smaller major axis lengths to a domination of $0^{\circ}$ values toward $2.0$ mm is extremely accurately simulated due to the synthetic 2D-S filtering technique described in Sect.~\ref{sect:MC_sim}. This leads to the appearance of a size dependence in the aspect ratio/major axis and area ratio/major axis joint distributions despite the underlying 3D ellipsoid distribution and SHPSG-generated models not exhibiting any explicit size dependence. The 2D-S probe filtering also explains the sharp discontinuities for the smallest projections which are caused by the size, shape, area ratio, and pixel number thresholds used for the observations. The main discrepancies between the observed and simulated ellipse fit parameters lies in the prevalence of simulated high area ratio ellipses for moderately-sized particles between 0.5 mm and 1.0 mm. This is particularly apparent for the Irregular 2020 SHPSG Fused Irregular simulation of Fig.~\ref{fig:SHPSG_Irreg_2020} which shows a more disperse set of ellipse area ratios at each major axis size compared to both the observations and the ellipsoid simulations. This discrepancy can be explained by comparing the two STL simulations in Fig.~\ref{fig:SHPSG_Irreg_2020} which shows the tendency for the \citet{Schima2024} algorithm to classify these high area ratio/aspect ratio projections most likely as graupel rather than as irregular. 
 
\FloatBarrier
\begin{figure}[ht]
    \centering
    %\includegraphics[width=\textwidth]{Figures/IMPACTS_2DS_Simulation_STL_Irregular_ellipsoids_6x3.png}
        \includegraphics[width=\textwidth]{Figures/IMPACTS_2020_ellipse_fit_joint_dists_stack_OBS_ELLIPSOIDS_NEW.pdf}
    \caption{Comparison of PCGL irregular particle ellipse parameters (left column) and simulated ellipsoids (right column) using parameters provided in the text for all eight 2020 cases. Log-scaled colormap values and ranges for each panel are consistent with Fig.~\ref{fig:ellipse_fit_joint2_2020}.}
    \label{fig:EllipSim_Irreg_2020}
\end{figure}
\FloatBarrier

\FloatBarrier
\begin{figure}[ht]
    \centering
    \includegraphics[width=\textwidth]{Figures/IMPACTS_2020_irregular_ellipse_fit_joint_dists_stack_OBS_STL_STL_Irreg_FIXED.pdf}
       % \includegraphics[width=\textwidth]{Figures/IMPACTS_2020_ellipse_fit_joint_dists_stack_OBS_SHPSG_fused_NEW.pdf}
    \caption{As in Fig.~\ref{fig:EllipSim_Irreg_2020} but with SHPSG-generated irregular particles filtered by the  \citet{Schima2024} habit identification algorithm identification as ``irregular'' (center column) and without the additional habit identification filtering (right column).}
    \label{fig:SHPSG_Irreg_2020}
\end{figure}
\FloatBarrier
 
 
The joint distributions in Figs.~\ref{fig:EllipSim_Irreg_2020} and~\ref{fig:SHPSG_Irreg_2020} appear to visually capture the correct ellipse parameter joint distribution shapes and magnitudes of the observations but it is not as obvious how well these simulations quantitatively mimic the H and V channel observations. While there are several metrics that could be used here to compare the fidelity of the ellipse parameter joint distribution functions, here we develop a similarity metric (as a percentage) via the Maximum Mean Discrepancy \citep[MMD; see][]{Gretton2012} between the 2020 PCGL Irregular observations and the ellipsoid and STL simulations. \citet{Gretton2012} defines the MMD as a non-parametric integral probability metric that measures the distance between two multivariate distributions after  embedding each distribution into a Reproducing Kernel Hilbert Space. To do this with the discrete joint probability distribution functions, we first standardize the values from each six joint distributions using the maximum and minimum values from the combined pool of observations and simulated ellipse parameter values. We then use a Radial Basis Function (RBF) as the embedding kernel $K(x,y) \equiv \exp\left[-||x-y||^{2}/(2\sigma^{2}) \right]$ with $\sigma=0.5$ and we calculate
\begin{equation}\label{eq:MMD}
\mathrm{MMD}(X, Y) = \sqrt{ \frac{1}{n^2} \sum_{i=1}^n \sum_{j=1}^n k(x_i, x_j) + \frac{1}{m^2} \sum_{i=1}^m \sum_{j=1}^m k(y_i, y_j) - \frac{2}{nm} \sum_{i=1}^n \sum_{j=1}^m k(x_i, y_j) }
\end{equation}
for 10,000 randomly sampled data pairs between the observations and simulations. This subsampling is used to reduce the memory footprint as the number of calculations increases as $\mathcal{O}(N^{2})$. To ensure that the  subsampled population is statistically representative, we bootstrap sample MMD values 30 times and report only the largest MMD value. To simplify interpretation of the MMD results, we develop a similarity metric using the maximum possible value of the MMD: $\mathrm{MMD}_{\max} = \sqrt{2 - 2 \exp\left[- 2/(2\cdot 0.5^{2})\right]} \approx 1.401$. We then define a similarity percentage based on each normalized MMD value
\begin{equation}
S \equiv \left(1 - \frac{\mathrm{MMD}}{\mathrm{MMD}_{\max}}\right) \cdot 100\%,
\end{equation}
where $S = 100\%$ is perfect similarity and $S=0\%$ indicates perfect dissimilarity. These similarity percentages are shown in Fig.~\ref{fig:Irreg_2020_Similarity} for each 2020 PCGL irregular simulation compared against each eight 2020 flight data as well as the 2020 combined data used in Figs.~\ref{fig:EllipSim_Irreg_2020} and~\ref{fig:SHPSG_Irreg_2020}. The ellipsoid simulations by far exhibit the most similarity where $S$ percentages for each ellipse parameter combination in each channel are all greater than $90\%$ for each individual 2020 case and greater than $95\%$ for all 2020 cases combined. The STL simulations generally exhibit lower similarity values than the ellipsoid simulations. Furthermore, as seen in Fig.~\ref{fig:SHPSG_Irreg_2020}, the irregular filtering also impacts the simulated similarity as the filtered irregular particle simulation produces higher $S$ percentages for the area ratio-aspect ratio joint distributions due to the high area ratio hitch caused by the \citet{Schima2024} decision tree logic. Conversely, including all simulated STL PCGL irregular particles leads to generally higher similarity percentages for the ellipse aspect ratio/orientation angle joint distributions because the \citet{Schima2024} algorithm disturbs the bivariate beta-like symmetry by removing counts for higher projected aspect ratio values and for angles closer to $\pm 90^{\circ}$. Overall, even the STL SHPSG simulations achieve at least $84.35\%$ similarity compared to the 2020 PCGL Irregular observations for each individual case.
 
 %The here is a notable transition to near $0^{\circ}$ ellipse angles for particles larger than about 1.5 mm in the observations. As stated before, however, this can be simply explained by the inability for the larger particles to stay fully within the 2D-S buffer strips unless these large particles are oriented horizontally. The second main discrepancy comes from the observed beta-like area ratio distribution about the mean area ratio at each ellipse major dimension. There is a slight decrease in area ratio for increasing sizes that can be somewhat captured by the exponential density modulation factor used for the ellipsoid simulation. However, it is more difficult to simulate the symmetric nature of the observed area ratios with this modulation factor. The STL simulation does a better job simulating this symmetric area ratio distribution, however the area ratio distribution is still somewhat negatively skewed like the ellipsoid simulation.  
 
While the tests in Fig.~\ref{fig:EllipSim_Irreg_2020} can accurately explain over $90\%$ of the joint ellipse distribution characteristics for irregular particles, there is an outstanding question as to how many similar ellipsoid distributions could also yield these irregular particle ellipse parameters. Of note is the suspicious, nearly identical H and V aspect ratio projection distributions for the observed mirrored ellipse orientation distributions. Figure~\ref{fig:EllipSim_Irreg_AR_2020_hists} investigates this possibility by providing several sensitivity tests of various oblate-like and prolate-like distributions which each yields consistent mean projection aspect ratios to the observations  using the ellipsoid orientation parameters from the Irregular 2020 simulation in Table~\ref{table:ellip_orient_tests}. In addition, three other orientation distributions (Parallel $\overline{\psi}_{p}$, Skewed $\overline{\psi}_{p}$, and Orthogonal $\overline{\psi}_{p}$) are used to investigate the effects that various $\psi_{p}$ orientations other than random might have on the projections. Oblate-like distributions (top row) have projections that are the most sensitive to these various orientation sets for $\psi_{p}$. For these ellipsoids, the random $\psi_{p}$ orientations that were used for the ellipsoids in Fig.~\ref{fig:EllipSim_Irreg_2020} produce an unobserved bimodality in the projected H and V ellipse aspect ratios. The parallel and orthogonal $\psi_{p}$ orientations produce asymmetric projected ellipse aspect ratio distributions where the parallel $\overline{\psi}_{p}$ orientations peaks at around 0.6 in the horizontal channel whereas aspect ratios peak near unity in vertical channel. The reverse behavior occurs with the orthogonal $\overline{\psi}_{p}$ simulation where the $a$ and $b$ axes are predominantly orthogonal to the horizontal channel but parallel to the vertical channel. As similarly demonstrated in \citet{Dunnavan2022}, the only way to reasonably yield the observed symmetric, nearly identical H and V aspect ratio distributions for oblate-like shapes is if the $\psi_{p}$ orientation is equally skewed with respect to each channel (i.e., $\psi_{p}=45^{\circ}$). However, in this case, the aspect ratio distribution in the vertical channel becomes much more narrow than in the horizontal channel. Overall, the simplest explanation is that the irregular particles are best represented as triaxial ellipsoids with $\varphi_{ba}=\varphi_{cb}\approx 0.75$ with preferential orientations at $\phi_{p} = 45^{\circ}$ and $\theta_{p} \approx 120^{\circ}$. Prolate-like distributions (bottom row of Fig.~\ref{fig:EllipSim_Irreg_AR_2020_hists}) have projections with the least amount of sensitivity to various $\psi_{p}$ distributions. There would be no sensitivity if each ellipsoid was a prolate spheroid as the $\psi_{p}$ rotation (as shown in Fig.~\ref{fig:ellip_orient}) is about the $a$ axis, which is an axis of symmetry for prolate spheroids. All projection variability when using $\psi_{p}$ distributions therefore comes from slight deviations in the 3D ellipsoids from prolate form due to the use of Eq.~\ref{eq:beta_asp_dists}.  While prolate-like ellipsoid distributions better explain the Irregular IMPACTS 2020 observations due to more symmetric forms in both H and V channels, there are some notable deviations that are most prominent when the projected ellipse aspect ratios are closest to unity. For the horizontal channel, the projection distributions switch from concave up to concave down for each $\psi_{p}$ orientation set. This non-beta behavior is inconsistent with both the IMPACTS PCGL irregular observations and the triaxial ellipsoid distribution tests used in Fig.~\ref{fig:EllipSim_Irreg_2020}. Furthermore, there are slight shifts between the H and V distribution modes which are inconsistent with observations. The $\psi_{p}$ orientation tests for the $\overline{\varphi}_{ba}=\overline{\varphi}_{cb}=0.75$ ellipsoids (middle row) show the same behaviors as both the oblate-like and prolate-like distributions but with a sensitivity in between each spheroid. As such, the simplest explanation for the emergence of the H and V channel joint distributions of ellipse aspect ratio and the orientation is that $\overline{\varphi}_{ba}=\overline{\varphi}_{cb}$ and $\psi_{p}$ is randomly distributed.


%by oblate-like or even prolate-like ellipsoids.


\FloatBarrier
\begin{figure}[ht]
    \centering
    \includegraphics[width=0.95\textwidth]{Figures/Similarity_2020_cases_hist.pdf}
    \caption{Maximum Mean Discrepancy (MMD) Similarity (as a percentage) for the ellipsoid and STL PCGL irregular simulations for all eight individual 2020 flights and from all 2020 flights combined (i.e., Figs.~\ref{fig:EllipSim_Irreg_2020} and~\ref{fig:SHPSG_Irreg_2020}).}
    \label{fig:Irreg_2020_Similarity}
\end{figure}
\FloatBarrier

%\begin{table}[htbp]
%\centering
%\caption{Bootstrapped Maximum Mean Discrepancy \citep[MMD; see][]{Gretton2012} evaluating the divergence between observed and simulated 2D joint distributions ($N=10,000$ per draw, $n=30$ iterations, and a Gaussian RBF kernel bandwidth $\sigma=0.5$) for the joint distributions shown in Figs.~\ref{fig:EllipSim_Irreg_2020} and~\ref{fig:SHPSG_Irreg_2020}. The observed and simulation results are normalized to [0,1] using the shared minimum and maximum values before performing the MMD calculations. MMD values are reported as Mean $\pm$ Standard Deviation.}
%\label{tab:mmd_full_comparison}
%\begin{tabular}{lcccc}
%\toprule
%\toprule
%\textbf{Joint Distribution} & \textbf{Channel} & \textbf{Ellipsoids} & \textbf{STL (Irregular Filtered)} & \textbf{STL (All Particles)} \\
%\midrule
%\multirow{2}{*}{Major Axis vs. Aspect Ratio} 
% & Horizontal & $0.0333 \pm 0.0035$ & $0.1081 \pm 0.0041$ & $0.0489 \pm 0.0031$ \\
% & Vertical   & $0.0271 \pm 0.0041$ & $0.0916 \pm 0.0033$ & $0.0404 \pm 0.0037$ \\
%\addlinespace
%\multirow{2}{*}{Major Axis vs. Area Ratio} 
% & Horizontal & $0.0144 \pm 0.0029$ & $0.0686 \pm 0.0039$ & $0.1187 \pm 0.0043$ \\
% & Vertical   & $0.0267 \pm 0.0032$ & $0.0541 \pm 0.0036$ & $0.0958 \pm 0.0033$ \\
%\addlinespace
%\multirow{2}{*}{Ellipse Angle vs. Aspect Ratio} 
% & Horizontal & $0.0551 \pm 0.0039$ & $0.1005 \pm 0.0043$ & $0.0377 \pm 0.0037$ \\
% & Vertical   & $0.0603 \pm 0.0040$ & $0.0856 \pm 0.0035$ & $0.0363 \pm 0.0046$ \\
%\addlinespace
%\multirow{2}{*}{Ellipse Angle vs. Area Ratio} 
% & Horizontal & $0.0475 \pm 0.0040$ & $0.0408 \pm 0.0035$ & $0.0548 \pm 0.0028$ \\
% & Vertical   & $0.0561 \pm 0.0047$ & $0.0391 \pm 0.0044$ & $0.0488 \pm 0.0043$ \\
%\addlinespace
%\multirow{2}{*}{Ellipse Angle vs. Major Axis} 
% & Horizontal & $0.0471 \pm 0.0060$ & $0.0633 \pm 0.0035$ & $0.1136 \pm 0.0036$ \\
% & Vertical   & $0.0585 \pm 0.0047$ & $0.0504 \pm 0.0046$ & $0.0906 \pm 0.0047$ \\
%\addlinespace
%\multirow{2}{*}{Area Ratio vs. Aspect Ratio} 
% & Horizontal & $0.0339 \pm 0.0035$ & $0.1190 \pm 0.0029$ & $0.1151 \pm 0.0053$ \\
% & Vertical   & $0.0330 \pm 0.0028$ & $0.0962 \pm 0.0047$ & $0.0898 \pm 0.0047$ \\
%\bottomrule
%\end{tabular}
%\end{table}


\FloatBarrier
\begin{figure}[ht]
    \centering
    \includegraphics[width=\textwidth]{Figures/Irregular_2020_ellipsoid_tests_hists_new2.png}
    \caption{Ellipsoid distribution projected aspect ratio tests using parameters in Table~\ref{table:ellip_orient_tests}. Horizontal and vertical projected aspect ratio distributions are shown in the left and middle columns, respectively. Each row consists of a separate set of ellipsoid distribution mean aspect ratios with the right column showing the Monte Carlo joint distribution of sampled $\varphi_{ba}$ and $\varphi_{ca}$ values. IMPACTS 2020 distributions (filled) represent the \citet{Schima2024} identified irregular particles for all eight 2020 cases in the PCGL.}
    \label{fig:EllipSim_Irreg_AR_2020_hists}
\end{figure}
\FloatBarrier

While the multivariate analytical ellipsoid distribution model in Sect.~\ref{methods:math} can very accurately reproduce the full suite of 2020 PCGL irregular ellipse joint distributions, it is more difficult to use this model to reproduce the 2022-2023 ellipse distributions. The corresponding comparison of observations with theoretical ellipsoids is shown in Fig.~\ref{fig:EllipSim_Irreg_2022_2023}. While the 2022-2023 H channel ellipse fit results are quite similar to that of 2020, the vertical channel demonstrates a larger spread of orientations where the modal ellipse angle is approximately $-10^{\circ}$ but with many counts toward more negative angles. For this primary population of irregular particles, we assume $\overline{\phi}_{p} = 170^{\circ}$ and $\overline{\theta}_{p} \approx 50^{\circ}$ which is needed to match the $-40^{\circ}$ H channel mode with the $-10^{\circ}$ V channel mode  (Table~\ref{table:ellip_orient_tests}). Thus, unlike the Irregular 2020 PCGL simulations which exhibited a skewed $\overline{\phi}_{p} = 45^{\circ}$ angle, the primary particles needed for the 2022-2023 irregular projections appear to require the particles' major axes to align almost entirely along the $\hat{x}$ direction (i.e., along air flow/aircraft heading). The observations also show an apparent bimodality in the V channel angles that is most apparent in the ellipse major axis size and orientation angle plot. We hypothesize that this additional mode is the result of a separate population of prolate-like particles with preferred orientations consistent with those of ``Linear'' particles. As shown in Table~\ref{table:ellip_orient_tests}, we use highly oriented, linear-like prolate particles with mean projection orientations of $\theta_{\mathrm{proj},h}=-70^{\circ}$ and $\theta_{\mathrm{proj},v}=-75^{\circ}$ and a population fraction of 0.25 to simulate this additional mode. Although this comparison is not as representative of the 2022-2023 PCGL irregular particles as the 2020 PCGL particles from Fig.~\ref{fig:EllipSim_Irreg_2020}, the dual-population ellipsoid simulation is capable of explaining many of the ellipse joint distribution characteristics. For example, the secondary mode of prolate particles appears with slightly higher area ratios than the primary irregular habit such that there are notches in the observed ellipse joint area ratio/orientation plots for $\mathrm{AR}>0.8$ and $\overline{\theta}_{\mathrm{proj},h}$ and $\overline{\theta}_{\mathrm{proj},v} \approx -75^{\circ}$. These notches align with the notches of the observed area ratio/aspect ratio joint distribution plots for aspect ratios between about 0.4 and 0.8 and $\mathrm{AR}>0.8$. As discussed above, these same notches appear in Figs.~\ref{fig:EllipSim_Irreg_2020} and~\ref{fig:SHPSG_Irreg_2020} and appear to result from the use of the \citet{Schima2024} algorithm to filter only irregular habits. While the theoretical ellipsoid simulations do not exactly replicate these notch locations, they do appear to be very close to the exact same H and V channel orientation modes for the linear habits from 2020 ($\overline{\theta}_{\mathrm{proj},h}=-70^{\circ}$ and $\overline{\theta}_{\mathrm{proj},h}=+80^{\circ}$; Figs.~\ref{fig:ellipse_fit_joint2_2020} and~\ref{fig:EllipSim_Linear_2020}) and from the 2022-2023 cases ($\overline{\theta}_{\mathrm{proj},h}=-80^{\circ}$ and $\overline{\theta}_{\mathrm{proj},h}=-80^{\circ}$; Figs.~\ref{fig:ellipse_fit_joint2_2022_2023} and~\ref{fig:EllipSim_Linear_2022_2023}). This provides strong evidence that even in the PCGL there are prolate- or columnar-like particles which exhibit a much more concentrated set of equilibrium orientations where $\kappa\sim 10$ rather than the lower kappa values for irregular particles (see Table~\ref{table:ellip_orient_tests}). Appendix~\ref{sect:discuss_King} provides explicit CFD simulations that both explain these 2022-2023 linear equilibrium orientations in the H and V channels and why a mistaken $90^{\circ}$ rotation in the 2020 2D-S probe configuration could explain why orientations from flights during that year differ from those derived from the 2022 and 2023 flights.  

\FloatBarrier
\begin{figure}[ht]
    \centering
    %\includegraphics[width=\textwidth]{Figures/IMPACTS_2DS_Simulation_STL_Irregular_ellipsoids_6x3.png}
        \includegraphics[width=\textwidth]{Figures/IMPACTS_2022_2023_ellipse_fit_joint_dists_stack_OBS_ELLIPSOIDS_NEW.pdf}
    \caption{As in Fig.~\ref{fig:EllipSim_Irreg_2020} but for all eighteen 2022 and 2023 cases.}
    \label{fig:EllipSim_Irreg_2022_2023}
\end{figure}
\FloatBarrier

Figures~\ref{fig:EllipSim_Linear_2020} and~\ref{fig:EllipSim_Linear_2022_2023} compare the 2020 and 2022-2023 NGL \citet{Schima2024} identified ``Linear'' ellipse fit parameters with both theoretical triaxial ellipsoid simulations as well as explicit Hexagonal Column STL models. Both the theoretical ellipsoid and STL simulations explain the modal locations of joint ellipse parameters in each plot where aspect ratios and area ratios are approximately 0.3 and 0.7 and the H (V) orientations are $-70^{\circ}$ ($80^{\circ}$) and $-70^{\circ}$ ($-70^{\circ}$) for the 2020 and 2022-2023 cases, respectively. The STL simulations present a much more concentrated set of values and the ellipsoids present a more broad set of ellipse parameter values in each joint distribution. This can be explained by the use of the \citet{Schima2024} algorithm to filter linear habits in both the observations and in the STL simulations. The decision tree logic of the habit identification algorithm does not classify aspect ratios close to unity as linear despite the possibility that such linear particles could be oriented with their major axes oriented toward the H or V channel receiver. Conversely, the larger variability in the observed joint distributions compared to the STL simulations can be explained by misclassifications of other habits as linear particles. For example, planar crystals can appear linear if their prism facets are oriented toward the H or V channel receivers. Furthermore, the STL simulations, unlike the ellipsoid simulations, were similarly filtered using the exact same \citet{Schima2024} algorithms as used in the observations in Fig.~\ref{fig:EllipSim_Linear_2022_2023}. Therefore, the pronounced discontinuities present in the orientation angles at approximately $-80^{\circ}$, $-25^{\circ}$ and $85^{\circ}$ from both the observations and the STL simulations are almost certainly the result of misclassifications of linear (columnar) habits into one or more of the other \citet{Schima2024} habit categories.

The theoretical ellipsoid and STL simulations for linear particles also exhibit the same H/V channel symmetry as in the observations. This curious H/V symmetry results from a combination of relatively precise Fisher orientation parameters where $\kappa = 8$ and mean orientation angles are $\overline{\phi}_{p}=80^{\circ}$ ($\overline{\phi}_{p}=100^{\circ}$) and $\overline{\theta}_{p} \approx 115$ ($\overline{\theta}_{p} \approx 45^{\circ}$) for the 2020 (2022-2203) simulations  (Table~\ref{table:ellip_orient_tests}). In this way, the linear particles for both 2020 and 2022-2023 are preferentially oriented such that their major axes are within $10^{\circ}$ of the $\hat{y}$ direction (see Figs.~\ref{fig:P3} and ~\ref{fig:ellip_orient}). This suggests that these linear particles are indeed orienting themselves such that their broadsides are more-or-less along the $\hat{x}$ air flow direction. As speculated by \citet{King1985d}, aircraft lift also acts to provide preferential canting in the X-Z (H channel) plane. For the 2022-2023 flights, this precise canting of almost exactly $45^{\circ}$ combined with near broadside flow orientation lead to both H and V channels viewing roughly the same amount projection foreshortening. For the 2020 cases, it appears that this degree of canting is instead about $25^{\circ}$ which creates a slight asymmetry in the projected aspect ratios between the H and V channels. This slight asymmetry can actually be seen in  Fig.~\ref{fig:ellipse_fit_joint2_2020} where the power-law-like tail of NGL linear particles in the ellipse major axis/aspect ratio joint distribution plot is more pronounced in the V channel than in the H channel. Area ratios between the simulations and observations are largely consistent with a majority of projections having area ratios between 0.6 and 0.7 and orientation modes around -70$^{\circ}$ with the exception of the \citet{Schima2024} algorithm discontinuities which obscures the true location of this mode.

% The simulated ellipse aspect ratio and ellipse major axis length joint distributions show a similar power-law-like dependence as the observations but with less variability due to the idealized nature of the prolate-like ellipsoids used.

\FloatBarrier
\begin{figure}[ht]
    \centering
  %  \includegraphics[width=\textwidth]{Figures/IMPACTS_2DS_Simulation_STL_Columns_ellipsoids_6x3.png}
    \includegraphics[width=\textwidth]{Figures/IMPACTS_2020_linear_ellipse_fit_joint_dists_stack_OBS_ELLIPSOIDS_STL_NEW.pdf}
    \caption{As in Fig.~\ref{fig:EllipSim_Irreg_2022_2023} but for 2020 ``Linear'' NGL particles using the Hexagonal Columns parameters from Table~\ref{table:ellip_orient_tests}.}
    \label{fig:EllipSim_Linear_2020}
\end{figure}
\FloatBarrier

\FloatBarrier
\begin{figure}[ht]
    \centering
  %  \includegraphics[width=\textwidth]{Figures/IMPACTS_2DS_Simulation_STL_Columns_ellipsoids_6x3.png}
    \includegraphics[width=\textwidth]{Figures/IMPACTS_2022_2023_linear_ellipse_fit_joint_dists_stack_OBS_ELLIPSOIDS_STL_NEW.pdf}
    \caption{As in Fig.~\ref{fig:EllipSim_Irreg_2022_2023} but for 2022-2023 ``Linear'' NGL particles using the Hexagonal Columns parameters from Table~\ref{table:ellip_orient_tests}.}
    \label{fig:EllipSim_Linear_2022_2023}
\end{figure}
\FloatBarrier

%\FloatBarrier
%\begin{figure}[ht]
%    \centering
%    \includegraphics[width=0.85\textwidth]{Figures/IMPACTS_2DS_Simulation_STL_Columns_NGL_linear_comp.png}
%    \caption{As in \cref{fig:EllipSim_Linear_2020} but for the NGL ``Linear'' particles and the Hexagonal Columns AV70 particles identified as ``Linear'' using the \citet{Schima2024} algorithm.}
%    \label{fig:STL_Linear_2020}
%\end{figure}
%\FloatBarrier

%\cref{fig:STL_MISC_tests} shows a suite of simulations using the STL models from \cref{fig:stl_ex}. Needle ellipse aspect and area ratios are generally consistent with the ``Linear'' results from \cref{fig:EllipSim_Linear_2022_2023} but with aspect ratios that peak at values slightly above 0.1. Dendrite and hexagonal plates with planar flutter orientations have low aspect ratios that peak close to the 3D STL model aspect ratios as viewed in the horizontal channel but high, near unity aspect ratio values in the vertical channel. Area ratios in the horizontal are much broader than in the vertical channel. The capped column STL simulations are generally consistent with the irregular distributions shown in \cref{fig:EllipSim_Irreg_AR_2020_hists} where a majority of the values span from approximately 0.6 to 0.8 and 0.4 to 0.8 for aspect ratios and area ratios, respectively. This suggests that capped column particles might be difficult to identify from their ellipse fit parameters compared to irregular particles. Bullet rosettes show a wide range of different ellipse fit parameter distributions depending on the number of bullets. BR03 has the broadest distributions where the 1.0 mm BR03 STL model can have ellipse aspect ratio fits that range from 0.2 to 1.0 and area ratios that range from 0.3 to 0.7. BR04 has a slightly lower range of ellipse aspect ratio values but very narrow values of area ratios whereas BR08 has a narrower range of aspect ratio values but a broader distribution of area ratio. Overall, the bullet rosette simulations illustrate how variable ellipse fit parameters are for only three various manifestations. Much of this variability inevitably stems from how the 3D ellipsoid fit itself does not accurately characterize the intricate bullet rosette morphology.  

%Figure~\ref{fig:STL_MISC_tests} shows that the $\varphi_{ca}/\phi_{\mathrm{proj}}$ 3D/2D correction factors for irregular and bullet rosette simulations are closer to unity compared to the other simulations, but have the largest variations. 

The irregular and linear tests shown in this section have demonstrated how plausible 3D ellipsoids or STL models with specific Fisher orientation distributions can explain the full suite of observed joint distributions between each ellipse parameter. However, how accurate can 3D morphological and geometric retrievals be with respect to using their 2D projections with their CGAL ellipse/circle measures? Furthermore, how much do the inferred 3D parameters deviate between the H and V channels and how do these uncertainties depend upon habit- and (potentially) instrument-based orientation discrepancies? Figure~\ref{fig:STL_MISC_tests} investigates the inverse (i.e., retrieval) problem by comparing various 3D measures with their 2D H/V projection counterparts using a suite of habit-specific STL (see Fig.~\ref{fig:stl_ex}) simulation tests from Table~\ref{table:ellip_orient_tests}. To better resolve the ambiguity in 3D and 2D measures and definitions, we utilize seven separate combinations of 3D and 2D measures: one for shape, three for size, two for density, and one for complexity. In this way, we show correction factors needed to convert each 2D measure into a 3D counterpart. For shape, we define a 3D/2D correction factor of $\varphi_{ca}/\phi_{\mathrm{proj}}$, which is the ratio of the 3D best-fit ellipsoid $\varphi_{ca}$ divided by the CGAL ellipse fit aspect ratio. For size, we define $D_{\max}/D_{\mathrm{circ}}$, which is the caliper (maximum) dimension divided by the CGAL circle diameter of minimum enclosure, $D_{\max}/(2 r_{a})$, which is the caliper (maximum) dimension divided by the CGAL ellipse fit major dimension, and $a/r_{a}$, which is the best-fit ellipsoid semi-major axis divided by the CGAL ellipse fit semi-major dimension. For density/area ratio, we define $\rho_{\mathrm{ellip}}/\mathrm{AR}_{\mathrm{ellip}}$ and $\rho_{\mathrm{sph}}/\mathrm{AR}_{\mathrm{circ}}$ which are the ratios of the best-fit ellipsoid density to the CGAL ellipse fit area ratio and the best-fit sphere density to the circle of minimum enclosure area ratio, respectively. Here, we assume no hollowing such that $\rho_{\mathrm{ellip}}$ and $\rho_{\mathrm{sph}}$ are nearly mathematically equivalent to the volume ratio of the STL volume and the circumscribing ellipsoid or sphere. Note that this is slightly different than typical ellipsoid effective densities which are often defined as the ratio of particle mass and enclosing ellipsoid volume (note that this is also different than how effective densities are defined in \citet{Heymsfield2004} and other studies). Finally, we show the ratio of sphericity [$\Psi\equiv \pi^{1/3} (6 V)^{2/3} S^{-1}$] to circularity ($C \equiv 4 \pi A P^{-2}$) where $V$ and $S$ are the STL volume and surface area, respectively, and $A$ and $P$ are the 2D projection area and perimeter, respectively. It is important to note that most of these correction factors can be both above and below unity due to the competing effects of orientation-induced foreshortening with discrepancies introduced from the geometric fits. Foreshortening is most apparent for the two linear simulations whereby the prolate nature of these columns serves to maximize the discrepancy between the true aspect ratio and the apparent aspect ratio through the shortening of the apparent 2D projected maximum sizes. Alternatively, geometric bounding of the ellipsoids, ellipses, and circles can artificially increase the sizes and aspect ratios needed to ensure volumetric and areal enclosure of each STL model and projection, respectively (this can be seen in Fig.~\ref{fig:P3}c, for example). Furthermore, binary closure and filling of, for example the SHPSG-fused and bullet rosette STL model projections, can lead to slightly larger CGAL ellipse aspect ratios than their best-fit ellipsoids when particles are small.

Figure~\ref{fig:STL_MISC_tests} shows that 2D-to-3D multiplicative correction factors between the horizontal and vertical channels are virtually indistinguishable from one another for most habits and metrics. This apparent channel insensitivity reiterates the H/V symmetry in Figs.~\ref{fig:boxplots},  \ref{fig:ellipse_fit_joint2_2020}, and~\ref{fig:ellipse_fit_joint2_2022_2023} where the marginal and joint distributions of the ellipse parameters, with the exception of orientation, are nearly identical between  channels. This is surprising as these simulations use three different orientation assumptions. As shown in Table~\ref{table:ellip_orient_tests}, irregular and bullet rosettes use the 2020 irregular-inferred orientations whereby $\overline{\phi}_{p}=45^{\circ}$, $\overline{\theta}_{p}\approx 120^{\circ}$, and $\kappa=4.5$ whereas the linear and capped column STL models use the 2022-2023 inferred orientations for linear particles where $\overline{\phi}_{p}=100^{\circ}$, $\overline{\theta}_{p}\approx 45^{\circ}$, and $\kappa=8$. While the assumed random orientations of dendrites and plates naturally leads to the theoretical same H and V projection distributions within stochastic variability, the same behavior for irregular and linear particles is instead caused by very different geometric reasons. For irregular and bullet rosette particles, the major dimensions are preferentially oriented at $\phi_{p} = 45^{\circ}$ in the vertical (X-Y) plane. This, coupled with a preferred approximate $30^{\circ}$ canting angle from the horizontal (i.e., $\overline{\theta}_{p} \approx 120^{\circ}$) leads to similar amounts of foreshortening of the major dimension in the X-Y and X-Z planes. The 2022-2023 linear particles, however, achieve the same sort of symmetry effect but in a completely different way. The 2022-2023 linear particles, unlike the irregular particles, appear to have their major dimensions oriented almost completely along the $\hat{y}$ direction where $\overline{\phi}_{p} = 80^{\circ}$ ($\overline{\phi}_{p} = 100^{\circ}$) for the 2020 (2022-2023) simulations.  Therefore, if $\phi_{p}=90^{\circ}$ and $\theta_{p}=90^{\circ}$, then the maximum dimension would be maximally foreshortened in the H channel but without any foreshortening in the V channel. However, the concentrated orientations of $\overline{\theta}_{p} \approx 45^{\circ}$ and $\kappa = 8$ ensures that the degree of foreshortening of the major axis in \textit{both} channels is controlled by the skewed $\theta_{p} \approx 45^{\circ}$ canting angle. This skewed canting therefore reduces the foreshortening in the H channel the same amount as the increase in  foreshortening in the V channel. Thus, the aspect ratio distributions remain virtually same between both channels and the orientations, coincidentally, also appear similar as demonstrated in Fig.~\ref{fig:EllipSim_Linear_2022_2023}. Table~\ref{table:channel_similarity_wide} shows the median relative difference percentage and the Kolmogorov-Smirnov (K-S) statistic between the H and V channels for each metric and STL simulation from Fig.~\ref{fig:STL_MISC_tests}. Besides the Linear 2020 and Capped Column 2020 simulations, the maximum deviation of the median values between H and V channels is only $1.1\%$ and the maximum deviation of cumulative distribution functions from each channel deviates by at most $5.08\%$. The Linear 2020 and Capped Column 2020  simulations deviate substantially from this behavior because of the shared $25^{\circ}$ canting angle rather than the near perfect $45^{\circ}$ canting angle from the 2022-2023 simulations. This difference alone leads to a nearly 15\% difference in aspect ratios and between 11.65\% and 17.35\% difference in the median sizes in the Linear 2020 simulations with K-S statistics larger than 0.34. The volume-based measures for the Linear 2020 simulations show more similarity between H and V channels where the median relative differences are less than 9.04\% with K-S statistics less than 0.3483. Interestingly enough, the Capped Columns 2020 simulations invert the size-volume biases such that the aspect ratios and sizes exhibit lower median relative differences than the volume-based measures. This H/V discrepancy maximizes for $\Psi/C$ where the median relative difference is 33.6\% with a K-S statistic over 0.45. These size-volume biases for linear and capped column particles can be explained by the fact that the capped column particles generally have larger aspect ratios due to the ellipsoid fits having to accommodate the top and bottom capped plates. Similarly, the blockier capped columns also exhibit less foreshortening due to the larger aspect ratios. However, the additional capped plates also lead to an increase in both 3D and 2D complexity which similarly increases the sensitivity of $C$ to orientations in the H and V projections.

%(note that $\phi_{p} = 90^{\circ}$ aligns the major dimension along the $\hat{y}$, directly into the projection direction of the H receiver; see Fig.~\ref{fig:P3}).

%\FloatBarrier
%\begin{figure}[ht]
%    \centering
%    \includegraphics[width=\textwidth]{Figures/STL_MISC_histogram_comp.pdf}
%    \caption{Histograms of ellipse fit aspect ratios and area ratios for STL models from \cref{fig:stl_ex} with parameters from \cref{table:sim_tests}.}
%    \label{fig:STL_MISC_tests}
%\end{figure}
%\FloatBarrier


\FloatBarrier
\begin{figure}[ht]
    \centering
    \includegraphics[width=\textwidth]{Figures/Boxplot_2D_3D_habits_7cats_8habits_cbf.pdf}
    \caption{Boxplots of the correction factors (3D/2D) needed to convert 2D H and V projection measures into their corresponding 3D STL and minimum-enclosing ellipsoid/sphere measures for STL simulations in Table~\ref{table:ellip_orient_tests}. Boxplot whiskers represent 1.5 times the interquartile range.}
    \label{fig:STL_MISC_tests}
\end{figure}
\FloatBarrier

The 3D/2D correction factor values themselves in Fig.~\ref{fig:STL_MISC_tests} illustrate the large variability due to both orientation effects and habit effects for each measure. The Linear 2020 simulations exhibit the largest variability in the H channel due to the $25^{\circ}$ canting. Therefore, while the median $D_{\max}/D_{\mathrm{circ}}$ in the V channel is approximately at unity with an IQR less than 0.2, the H has a median above unity and an IQR between unity and approximately 1.25. $D_{\max}/2r_{a}$ has an overall negative bias whereas $a/r_{a}$ has a large positive bias with a much larger IQR due to the minimum-enclosing volume ellipsoid elongation effect for columns (i.e., Fig.~\ref{fig:P3}c). Overall, $D_{\max}/D_{\mathrm{circ}}$ produces 3D/2D correction factors that are closer to unity for almost all habits in both channels. The volumetric density ratios $\rho_{\mathrm{ellip}}/\mathrm{AR}_{\mathrm{ellip}}$ and $\rho_{\mathrm{sph}}/\mathrm{AR}_{\mathrm{circ}}$ generally show IQR values that are less than 0.25 which are also generally smaller for the ellipsoid ratios than for the sphere/circle ratios. The complexity parameter ratio $\Psi/C$ by far exhibit the most variability in each channels and this variability depends upon habit. Simple habits such as plates exhibit little variability in correction values whereas more complex habits such as the Dendrite STL model exhibits a huge variability where the IQR is over 0.5 and the non-outlier range is over 1.5. The capped columns and bullet rosette simulations similarly exhibit a large range of 3D/2D correction values where the Capped Columns 2020 simulation exhibits median values closer to unity in the V channel than in the H channel. 

\FloatBarrier
\begin{table}[htbp]
\centering
\caption{Horizontal vs. vertical channel divergence metrics for the simulations in Fig.~\ref{fig:STL_MISC_tests}. Values are presented as the Median Relative Difference (\%) with the mean value from the both channels combined as the reference value. The Kolmogorov-Smirnov (KS) statistic is shown in parentheses. Bold values indicate significant divergence of H/V correction factors.}
\label{table:channel_similarity_wide}
\resizebox{\textwidth}{!}{%
\begin{tabular}{c cccccccc}
\toprule
\toprule
\textbf{Metric} & \makecell{\textbf{Irregular}\\\textbf{2020}} & \makecell{\textbf{Bullet}\\\textbf{Rosettes}\\\textbf{2020}} & \makecell{\textbf{Linear}\\\textbf{2020}} & \makecell{\textbf{Linear}\\\textbf{2022--2023}} & \makecell{\textbf{Capped}\\\textbf{Columns}\\\textbf{2020}} & \makecell{\textbf{Capped}\\\textbf{Columns}\\\textbf{2022--2023}} & \textbf{Dendrite} & \textbf{Plates} \\
\midrule
Aspect Ratio (3D/2D) & 0.47 (0.0140) & 0.36 (0.0109) & \textbf{14.96 (0.3408)} & 0.18 (0.0050) & \textbf{6.32 (0.2252)} & 0.03 (0.0033) & 0.43 (0.0046) & 0.06 (0.0022) \\
Max Dim / $D_{\mathrm{circ}}$ & 0.96 (0.0492) & 0.61 (0.0429) & \textbf{11.65 (0.4105)} & 0.07 (0.0050) & \textbf{6.87 (0.3009)} & 0.00 (0.0037) & 0.02 (0.0040) & 0.02 (0.0042) \\
Max Dim / $2r_a$ & 0.94 (0.0445) & 0.75 (0.0413) & \textbf{17.35 (0.4023)} & 0.07 (0.0049) & \textbf{13.89 (0.4300)} & 0.12 (0.0049) & 0.01 (0.0034) & 0.02 (0.0042) \\
Major Axis ($a/r_a$) & 1.10 (0.0508) & 0.65 (0.0346) & \textbf{17.35 (0.4023)} & 0.07 (0.0049) & \textbf{13.97 (0.4284)} & 0.12 (0.0051) & 0.01 (0.0034) & 0.02 (0.0042) \\
Ellipsoid Area-Density & 0.27 (0.0089) & 0.44 (0.0120) & \textbf{5.43 (0.3483)} & 0.01 (0.0044) & \textbf{32.23 (0.3516)} & 0.35 (0.0053) & 0.17 (0.0059) & 0.00 (0.0034) \\
Sphere Area-Density & 0.62 (0.0108) & 0.54 (0.0135) & \textbf{9.04 (0.2404)} & 0.16 (0.0058) & \textbf{22.92 (0.2041)} & 0.07 (0.0036) & 0.12 (0.0055) & 0.05 (0.0018) \\
Morph. Complexity & 0.23 (0.0079) & 0.80 (0.0215) & \textbf{6.51 (0.1410)} & 0.07 (0.0039) & \textbf{33.60 (0.4513)} & 0.18 (0.0088) & 0.52 (0.0045) & 0.01 (0.0026) \\
\bottomrule
\end{tabular}%
}
\end{table}
\FloatBarrier

%\FloatBarrier
%\begin{table}[htbp]
%\centering
%\caption{Horizontal vs. Vertical channel divergence metrics for the simulations in Fig.~\ref{fig:STL_MISC_tests}. Values are presented as the Median Relative Difference (\%) with the Kolmogorov-Smirnov (KS) statistic in parentheses.}
%\label{tab:channel_similarity_wide}
%\resizebox{\textwidth}{!}{%
%\begin{tabular}{c ccccccc}
%\toprule
%\toprule
%\textbf{Metric} & \makecell{\textbf{Irregular}\\\textbf{2020}} & \makecell{\textbf{Bullet}\\\textbf{Rosettes}\\\textbf{2020}} & \makecell{\textbf{Linear}\\\textbf{2020}} & \makecell{\textbf{Linear}\\\textbf{2022--2023}} & \makecell{\textbf{Capped}\\\textbf{Columns}\\\textbf{2022-2023}} & \textbf{Dendrite} & \textbf{Plates} \\
%\midrule
%Aspect Ratio (3D/2D) & 0.47 (0.0140) & 0.36 (0.0109) & \textbf{14.96 (0.3408)} & 0.18 (0.0050) & 0.03 (0.0033) & 0.43 (0.0046) & 0.06 (0.0022) \\
%Max Dim / $D_{\mathrm{circ}}$ & 0.96 (0.0492) & 0.61 (0.0429) & \textbf{11.65 (0.4105)} & 0.07 (0.0050) & 0.00 (0.0037) & 0.02 (0.0040) & 0.02 (0.0042) \\
%Max Dim / $2r_a$ & 0.94 (0.0445) & 0.75 (0.0413) & \textbf{17.35 (0.4023)} & 0.07 (0.0049) & 0.12 (0.0049) & 0.01 (0.0034) & 0.02 (0.0042) \\
%Major Axis ($a/r_a$) & 1.10 (0.0508) & 0.65 (0.0346) & \textbf{17.35 (0.4023)} & 0.07 (0.0049) & 0.12 (0.0051) & 0.01 (0.0034) & 0.02 (0.0042) \\
%Ellipsoid Area-Density & 0.27 (0.0089) & 0.44 (0.0120) & \textbf{5.43 (0.3483)} & 0.01 (0.0044) & 0.35 (0.0053) & 0.17 (0.0059) & 0.00 (0.0034) \\
%Sphere Area-Density & 0.62 (0.0108) & 0.54 (0.0135) & \textbf{9.04 (0.2404)} & 0.16 (0.0058) & 0.07 (0.0036) & 0.12 (0.0055) & 0.05 (0.0018) \\
%Morph. Complexity & 0.23 (0.0079) & 0.80 (0.0215) & \textbf{6.51 (0.1410)} & 0.07 (0.0039) & 0.18 (0.0088) & 0.52 (0.0045) & 0.01 (0.0026) \\
%\bottomrule
%\end{tabular}%
%}
%\end{table}
%\FloatBarrier


%\FloatBarrier
%\begin{figure}[ht]
%    \centering
%    \includegraphics[width=\textwidth]{Figures/IMPACTS_2DS_Ellipsoid_Simulation_NGL_Linear_2020.png}
%    \caption{As in \cref{fig:EllipSim_Irreg_2020} but for linear NGL particles using the Hexagonal Columns parameters from %\cref{table:ellip_dist_params}.}
%    \label{fig:EllipSim_Linear_2020}
%\end{figure}
%\FloatBarrier





%We use \cref{}

% 1.) 


%\subsubsection{Spheroid-like distributions}

% Oblate-like

% Prolate-like


%\subsubsection{Triaxial ellipsoid distributions}

%\subsection{Single ellipsoids}\label{sim:single}




% MAIN POINTS:
% 1.) b/a and c/b determine where the local peaks are in the projections. This is apparently true regardless of orientation (with some exceptions).
% 2.) Distributions of aspect ratio projections look highly unrealistic: they look too peaked and not smoothed enough. Distributions of orientation angle look more-or-less what the observations look like (i.e., governed by a von Mises distribution) and what the beta-distributed 3D aspect ratios look like in next section.
% 3.) Results from Jiang et al. 2016 and Dunnavan and Jiang 2019 are reiterated here for single spheroids. 


%\subsection{Distributions of ellipsoids}

% MAIN POINTS:
% 1.) If random psi, then we can relatively accurately retrieve the 3D orientation parameters through the following equations.
% 2.) Projected Orientation angles are roughly independent of the shapes of the ellipsoids and are approximately distributed as von Mises Distributions themselves with circular mean values that are consistent with the 3D orientation projections of the major axis direction.
% 3.) Any asymmetry in psi reflects into the X-Y and X-Z projected aspect ratio asymmetry. The fact that we don't see this in the in situ observations suggests that psi is probably symmetric (or uniform). (psi asymmetry reflects between H and V in narrowness/broadness of distribution) phiba/phicb asymmetry reflects in H and V distribution asymmetry.
% 4.) In order to have a symmetry in the H/V aspect ratio distributions for irregular particles, the orientations either have to be random or the ellipsoids themselves have to be symmetric in b/a and c/b (which seems to me to be much more likely).
% 5.)  


% x-z and y-z plane 3D retrieval estimates
%\begin{subequations}
%    \begin{equation}
%        \overline{\phi}_{p} \approx \mathrm{atan2}(\tan \overline{\theta}_{\mathrm{proj},xz},\tan \overline{\theta}_{\mathrm{proj},yz})
%    \end{equation}
%    \begin{equation}
%        \overline{\theta}_{p} \approx \mathrm{atan2}(\sec \overline{\phi}_{p},\tan \overline{\theta}_{\mathrm{proj},xz})
%    \end{equation}
%\end{subequations}


%\subsubsection{Random $\psi_p$ distribution}\label{sim:random_dist}


%\subsubsection{Retrieval of mean 3D angles from H-V projections}

%\Cref{sim:random_dist} illustrates that ellipsoids with random $\psi_p$ angles have projected orientation angles that are distributed similarly to a von Mises distribution. 



%    \overline{\phi}_{p} \approx	\begin{cases}
%       \overline{\theta}_{\mathrm{proj},v} & \overline{\theta}_{\mathrm{proj},v} \leq 90^{\circ} \\
%       \overline{\theta}_{\mathrm{proj},v} -180^{\circ} & (\overline{\theta}_{\mathrm{proj},v} > 90^{\circ}) \&  %(\overline{\theta}_{\mathrm{proj},v} \leq 270^{\circ})\\
%                          \overline{\theta}_{\mathrm{proj},v} -  360^{\circ} & (\overline{\theta}_{\mathrm{proj},v} > 270^{\circ}) \&  (\overline{\theta}%_{\mathrm{proj},v} \leq 360^{\circ})\\
%        \end{cases}





%\subsubsection{Von mises $\psi_p$ distribution}

%\FloatBarrier
%\begin{figure}[ht]
%    \centering
%    \includegraphics[width=\textwidth]{Figures/IMPACTS_EllipProj_Simulations_Obs_Irregular.png}
%    \caption{Joint distributions of ellipse fit angles and aspect %ratios for each 2020 IMPACTS case and for $T_{c}\leq -10\, ^{\circ}\mathrm{C}$. Theoretical ellipsoid simulations for $\overline{\varphi}_{ba}=\overline{\varphi}_{cb}=0.75$ and $\zeta=0.25$. }
%    \label{fig:IMPACTS_irregular_jointdist_2020}
%\end{figure}
%\FloatBarrier







%\FloatBarrier
%\begin{figure}[ht]
%\centering
%\includegraphics[width=\textwidth]{IMPACTS_2DS_SHPSG_Simulation_NGL_Linear_2020_MA94.png}
%\caption{As in \cref{fig:SHPSG_irreg_PCGL_2020} but for hexagonal columns.}
%\label{fig:SHPSG_linear_NGL_2020}
%\end{figure}
%\FloatBarrier



%% Generate marginal ellipse distribution histograms for suite of STL simulation tests.

% Figure 1: Ellipse parameters

% Figure 2: Difference between ellipse parameters and 3D ellipsoid parameters (e.g., a kind of error)

% Figure 3: Complexity parameters?









%\subsection{Implications for in situ observations}

%\FloatBarrier
%\begin{figure}[ht]
%    \centering
%    \includegraphics[width=\textwidth]%{Figures/P3_WFF_IMPACTS_0123.png}
%    \caption{3D model of the P3 aircraft with cloud probes}
%    \label{fig:P3}
%\end{figure}
%\FloatBarrier



%\subsection{Implications for microphysics models and in situ observations}

%Currently, both microphysics models and in situ observations commonly will assume single-aspect-ratio spheroidal shapes in order to perform various calculations. What are the implications for such calculations if instead  

%\FloatBarrier
%\begin{figure}[ht]
%\centering
%\includegraphics[width=\textwidth]{Figures/Irregular_2020_SHPSG_ellip_arcomp.png}
%\caption{Joint density comparison of SHPSG-generated ellipsoid aspect ratios for irregular particles (filled) with the bivariate beta distribution of ellipsoid aspect ratios used to generate ellipsoid results from \cref{fig:EllipSim_Irreg_2020} (black contours).}
%\label{fig:SHPSG_irreg_PCGL_2020_ellip_ar_comp}
%\end{figure}
%\FloatBarrier


\conclusions  %% \conclusions[modified heading if necessary]
%% Generate summary figure of joint ellipse parameters with boxes and example 2D-S images within boxes

%The results shown in this study, both from the 2D-S and from Monte Carlo simulations, suggests that many of the common assumptions regarding ice particle shapes and orientations in both models and inferred from observations are inconsistent with 3D considerations of ellipsoidal shapes and orientations. 

The results shown in this study, both from the 2D-S and from Monte Carlo simulations, suggest that many common spheroidal modeling assumptions of atmospheric ice particles are inconsistent with plausible 3D considerations of their best-fit ellipsoidal shapes and densities. NWP models and retrieval relations have traditionally treated ice particle shapes and densities for a variety of habits, such as irregular particles or snow aggregates, as oblate spheroids with fixed values or values that vary with size. The use of oblate spheroids to represent ice particles have historically stemmed from the need to introduce particle asymmetry into various microphysical or radiative calculations while assuming that particles have major axes that are randomly oriented in the horizontal. For example, dual polarization radar signatures in the DGL often show enhancements in differential reflectivity and specific differential phase. This polarimetric signature has been shown by \citet{Dunnavan2022} to result from local increases in irregular particle concentrations. However, irregular particles, like snow aggregates, apparently are not best represented as single aspect ratio oblate spheroids. Instead, there is a near-universal beta-like variability in both aspect ratios and area ratios for both H and V projections. This random variability about both the mean ellipsoid aspect ratios and mean area ratios strongly suggests that the 3D shape and densities similarly exhibit a form of random, symmetric variability. The minimum-volume enclosing ellipsoid $\varphi_{ba}$ and $\varphi_{ca}$ values are rather concentrated for irregular particles where $\overline{\varphi}_{ba}= \overline{\varphi}_{cb}=0.75$ and $\zeta = 0.6$ produce joint distributions for each combination of 2D ellipse parameter that are over 90\% similar to not only the composite distributions compiled from all eight of the 2020 cases, but also from each one individually. Perhaps most surprising is the lack of any size dependence needed to explain the in situ-derived ellipse aspect ratios or area ratios. This is in contrast to other studies such as \citet{Korolev2003} who, for example, found a slight size dependence in their aspect ratio data estimated from 2D particle images. However, as indicated in Fig.~\ref{fig:EllipSim_Irreg_2020}, this behavior can be a reflection of the in situ cloud probe processing itself, with larger particles more likely to not completely fit within the 1.28 mm 2D-S buffer width window. It is possible that other imagers, such as the CPI, could exhibit a similar sort of effect based upon each instrument's limitation. Overall, the apparent independence of aspect ratio and area ratio (i.e., ellipsoid effective densities) with size suggests that NWP and radiative models could possibly calculate irregular particle properties once and then scale these properties according to differing observed and modeled size distributions. Both numerical weather prediction models and observations that assume underlying 3D geometries in their derivations of microphysical properties \citep[e.g.,][]{Heymsfield2004} do not consider the size-independent variability of 3D irregular particle shapes and their cloud-probe-measured orientations. As such, these in situ studies and microphysics schemes rely upon an assumed unique 1-1 relationship between a chosen 2D particle length scale and 3D particle properties such as mass, volume, and density. Therefore, these studies do not take into account the observed shape and density beta-like variability as well as the 3D/2D correction factors needed based upon habit and orientation (Fig.~\ref{fig:STL_MISC_tests}). \citet{Heymsfield2004} in particular derives empirical power-law relations for effective densities of bullet rosettes using simplistic theoretical particle geometries. While these effective densities are based upon bulk microphysical quantities such as the ice water content, the inherent particle-by-particle variability of crystals such as bullet rosettes suggests that these semi-empirical studies are probably very limited. However, it appears that much of the variability that results from different habit types can potentially be captured analytically using beta distribution functions or other similar functions. 

Perhaps surprising in this study is the common symmetry and similarity in ellipse properties exhibited between the 2D-S H and V channels. While this observed symmetry for all ellipse parameters except for orientations would necessarily result from complete random particle orientations, as is sometimes assumed when using cloud probe instruments, it appears that this symmetry actually results from several other factors. For irregular particles, both ellipsoidal aspect ratio symmetry through $\overline{\varphi}_{ba} \approx \overline{\varphi}_{cb}$ as well as $\psi_{p}$ uniformity and either $\overline{\phi}_{p} \approx 45^{\circ}$ or $\overline{\phi}_{p} \approx 0^{\circ}$ and moderate canting angles of $\approx 30^{\circ}$ is necessary to explain the behavior of the joint ellipse distributions of aspect ratios and orientations. This observation confirms the generality of similar 2D-S ellipsoid simulation results from \citet{Dunnavan2022} which only considered one leg of the 20200207 IMPACTS case. Any deviation from $\varphi_{ba}=\varphi_{cb}$ leads to H/V channel asymmetries in the aspect ratio distributions and orientation angles that are not observed (Fig.~\ref{fig:EllipSim_Irreg_AR_2020_hists}). This is true when using both triaxial ellipsoids and more realistic and explicit randomly generated STL irregular-like particles. In fact, the spherical harmonic randomness of the SHPSG particles and the randomness associated with their fusions appears to explain why beta distributions are capable of producing similar ellipse distributions as deviations from the mean aspect ratio and density could themselves be considered somewhat random due to different habits and forms. It is unclear why there are vast differences in the irregular particle orientation distributions between the 2020 and 2022-2023 flights. Part of the discrepancy appears to be due to a potential mishap in the installation of the 2D-S in 2020 and the (corrected) re-installation of the 2D-S in 2021 right before the 2022-2023 flights (Appendix~\ref{sect:discuss_King}). The two main differences in each 2D-S configuration are the 2.8 degree corrected angle and the installation of a fairing onto the 2D-S cannister's extended pylon (see Fig.~\ref{fig:pylon_comp}). While these changes appear to be quite slight, explicit Computational Fluid Dynamic (CFD) simulations using a Computer-Aided Design (CAD) model of the P3 aircraft and the 2D-S instrument do show a shift in preferred equilibrium angles for the NGL linear particles which align better with the 2022-2023 IMPACTS data than the 2020 data. Appendix~\ref{sect:discuss_King} also provides evidence from the joint distributions of ellipse aspect ratios and ellipse orientations that the 2020 2D-S was rotated 90 degrees counter-clockwise compared to the correct 2022-2023 2D-S orientation (see Fig.~\ref{fig:rot_ex_discrep}). This swapping explains why the 2022-2023 H channel joint distributions are shockingly similar to those from the 2020 V channel. However, reconciling the 2020 H channel with the 2022-2023 V channel is harder due to the near $-10^{\circ}$ orientation mode of the primary irregular particles in the PCGL and the extended tail of negative orientation angles that seemingly represent blocky, prolate-like crystals (Fig.~\ref{fig:EllipSim_Irreg_2022_2023}) with orientation modes consistent with those of the NGL 2022-2023 simulations (Fig.~\ref{fig:EllipSim_Linear_2022_2023}). We hypothesize that the combination of a larger number of these blocky, prolate-like crystals in the 2022-2023 flights along with the fixed 2D-S orientation and correcting for the H/V channel swapping explains these peculiar discrepancies.  

%In particular, irregular particles as viewed by the 2D-S are not randomly oriented, as the mirrored, unimodal orientation distributions from Figs.~\ref{fig:ellipse_fit_joint2_2020} and~\ref{fig:ellipse_fit_joint2_2022_2023} show, but instead are preferentially oriented with respect to specific angles. For the 2020 IMPACTS cases, the mean orientation of the major axis appears to be $\overline{\phi}_{p} = 45^{\circ}$ and $\overline{\theta}_{p}\approx 120^{\circ}$ for the composite of all eight flights as well as for each individual flight in the PCGL. These mean orientations result primarily from irregular particles.   

It appears that many particles such as irregular particles, aggregates, hexagonal columns (needles), and liquid droplets can be identified by their 2D ellipse fit when guided by knowledge the habit-specific characteristic 2D projection ellipse parameter joint distributions. Thus, it is possible to use the results in this study to enhance 2D-S habit identification algorithms such as the \citet{Schima2024} one. While it is somewhat difficult to infer planar crystals (hexagonal plates and dendrites) from H and V ellipse parameter joint distributions, linear particles (e.g., NGL hexagonal columns and needles) have characteristic distributions that can be immediately inferred. As demonstrated in several other studies, the column/needle orientations in the vertical plane are quite specific and are anywhere from $60^{\circ}$ to $70^{\circ}$ depending on the mounting location of the probe below and in front of the wing tip. It also appears that slight changes to probe configurations and/or pylons can lead to significant changes in ellipse fit orientation angles and this can dramatically impact retrievals of columnar shapes, sizes, and densities. This can be seen in Fig.~\ref{fig:STL_MISC_tests} and Table~\ref{table:channel_similarity_wide} which show that the linear habit median maximum dimension varies between H and V channels by over 17\% in the ``Linear 2020'' simulation but only 0.07\% in the ``Linear 2022-2023'' simulation. The main difference between these simulations is with $\theta_{p}$ which results in an approximate $25^{\circ}$ canting (with respect to the labeled H channel) in 2020 compared to a $45^{\circ}$ canting in the 2022-2023 inferred orientations. It is the combination of columns, needles, and other prolate-like particles exhibiting a preferential vertical canting and flow-induced broadside orientation in the 2D-S that is the primary cause of the perceived H/V channel invariance in ellipse fit parameters. This invariance is much stronger in the 2022-2023 cases due to the near $45^{\circ}$ (i.e., $\overline{\theta}_{p} \approx 45^{\circ}$) canting angle compared to the lower $25^{\circ}$ (i.e., $\overline{\theta}_{p} \approx 115^{\circ}$) needed to explain the 2020 NGL linear observations. This $45^{\circ}$ vertical canting could possibly have resulted from the corrected 2D-S configuration contributing to better airflow within the sampling volume as suggested by explicit CFD simulations (Fig.~\ref{fig:CFD_tests}). Future in situ aircraft studies using OAPs will have to be mindful of such aircraft changes as these changes can apparently lead to different ice particle orientations as viewed from the 2D-S and potentially other instruments as well.  

% Despite this variability in size, the correction factors needed to retrieve density from area ratio are much more concentrated where mean values are between approximately 0.62 and 0.65 with standard deviations that are approximately 0.03 between both sets of experiments and channels. It is not clear why, for example, the 2022-2023 orientation distribution from the horizontal channel (labeled as the vertical channel during the 2022-2023 campaigns) deviates from the 2020 horizontal channel orientations. However, it appears that the P-3 modifications must have, in some way, altered the viewing direction and/or airflow around the 2D-S probe. Appendix B provides CFD simulations that explores the possibility of airflow differences as well as a hypothesis regarding different 2D-S mounting rotations. However, at this point, we still cannot explain the discrepancies between the 2020 and 2022-2023 ellipse fit orientations. Therefore, future in situ  studies will have to be mindful of such aircraft changes as these changes can apparently lead to different ice particle orientations as viewed from the 2D-S and potentially other instruments.   

%Overall, Fig.~\ref{fig:STL_MISC_tests} highlights that both habit type and 2D-S probe orientations lead to varying degrees of sensitivity when attempting to infer 3D properties from their 2D OAP projections and ellipse/circle fits.

%Similarly, the horizontal projections for the particles are roughly $-70^{\circ}$ or $-80^{\circ}$ but with more variability due to the increased gradients of torque exerted on the columns as a result of vortices generated by the aircraft wingtips in the vertical direction (\cref{fig:King_tests}).

%Planar crystals such as hexagonal plates and dendrites are difficult to identify both their aspect ratios and habit type. This is because planar crystals appear and will be classified as linear if viewed along their prism facets. However, proper identification requires a view parallel to planar crystal basal facets. If planar crystals flutter such that their basal facets are generally parallel to the vertical channel and their prism facets are generally parallel to the horizontal channel, then there will be asymmetric ellipse joint distributions between the H and V channels. However, the observed joint distributions are remarkably consistent between channels as demonstrated in
%\cref{fig:ellipse_fit_joint2_2020,fig:ellipse_fit_joint2_2022,fig:ellipse_fit_test,fig:boxplots}. It is possible that the lack of observed planar crystals  (\cref{fig:habit_heatmaps}) compared to irregular and columnar particles  explains why the ellipse distributions are consistent across the H and V channels. Future studies will have to consider these orientation effects when inferring planar crystal properties from OAP images.

%Bullet rosettes have highly variable ellipse fit properties and can therefore appear in many different places within the joint distributions diagrams for ellipse major axis lengths, aspect ratios, orientations, and area ratios. Much of this complexity comes from bullet rosettes with only a few branches as the best-fit ellipsoids depends crucially on the geometry of these branches. This suggests that complexity parameters, in conjunction with ellipse fit properties, can be used to infer the 3D morphologies of bullet rosettes and similarly complex irregular-like particles. Studies such as \citet{Schmitt2014} have similarly found success using such complexity parameters in describing irregular particles. The incorporation of ellipse fits however provides more information regarding the type and complexity of bullet rosettes since rosettes with more branches tend to appear more spherical in order to accommodate each additional branch. Therefore, rosettes with only a few branches can be quite elongated such that this information can be better identified using ellipse fit data. 



%% The following commands are for the statements about the availability of data sets and/or software code corresponding to the manuscript.
%% It is strongly recommended to make use of these sections in case data sets and/or software code have been part of your research the article is based on.

%\codeavailability{TEXT} %% use this section when having only software code available


%\dataavailability{TEXT} %% use this section when having only data sets available


%\sampleavailability{TEXT} %% use this section when having geoscientific samples available


%\videosupplement{TEXT} %% use this section when having video supplements available


\appendix
\section{Ellipsoid projection equations}\label{appen:A}   %% Appendix A


The canonical equation for a triaxial ellipsoid can be described by
\begin{equation}\label{d}
    \left(\frac{x}{c}\right)^{2} + \left(\frac{y}{b}\right)^{2} + \left( \frac{z}{a}\right)^{2} = 1. 
\end{equation}

The $Z-Y-Z$ Euler angle intrinsic rotation matrix using the convention in Fig.~\ref{fig:ellip_orient} gives 
\begin{equation}
\mathbf{A} = 
\begin{bmatrix}
\cos\phi_{p}\cos\theta_{p}\cos\psi_{p}-\sin\phi_p\sin\psi_p & -\cos\psi_p\sin\phi_p-\cos\phi_p\cos\theta_p\sin\psi_p & \cos\phi_p\sin\theta_p \\
\cos\phi_p\sin\psi_p+\cos\theta_p\cos\psi_p\sin\phi_p & \cos\phi_p\cos\psi_p-\cos\theta_p\sin\phi_p\sin\psi_p & \sin\phi_p\sin\theta_p \\
-\cos\psi_p\sin\theta_p & \sin\theta_p\sin\psi_p & \cos\theta_p
\end{bmatrix}.
\end{equation}
Ellipsoids projected onto a 2D $x^{\prime}$-$y^{\prime}$ plane will have rotated ellipses represented in the form
\begin{equation}
    C_{xx} {x^{\prime}}^{2} + C_{yy} {y^{\prime}}^{2} + C_{xy} x^{\prime} y^{\prime} = 1.
\end{equation}

Ellipsoids projected on the x-y plane (i.e., the V channel) will have corresponding ellipses with the parameters

\begin{subequations}\label{ellip_eq_vert}
    \begin{equation}
        C_{xx} = D_{22} - \frac{D_{23}^{2}}{D_{33}}
    \end{equation}
    \begin{equation}
        C_{yy} = D_{11} - \frac{D_{13}^{2}}{D_{33}}
    \end{equation}
    \begin{equation}
        C_{xy} = 2 \left( D_{12} - D_{13} \frac{D_{23}}{D_{33}}\right)
    \end{equation}
\end{subequations}
whereas ellipsoids projected onto the x-z plane (i.e., the H channel) will have corresponding ellipses with the parameters
\begin{subequations}\label{ellip_eq_horz}
    \begin{equation}
        C_{xx} = D_{33} - \frac{D_{23}^{2}}{D_{22}}
    \end{equation}
    \begin{equation}
        C_{yy} = D_{11} - \frac{D_{12}^{2}}{D_{22}}
    \end{equation}
    \begin{equation}
        C_{xy} = 2 \left( D_{13} - D_{12} \frac{D_{23}}{D_{22}}\right),
    \end{equation}
\end{subequations}
and ellipsoids projected onto the y-z plane will have the ellipse parameters
\begin{subequations}\label{ellip_eq_yz}
    \begin{equation}
        C_{xx} = D_{33} - \frac{D_{13}^{2}}{D_{11}}
    \end{equation}
    \begin{equation}
        C_{yy} = D_{22} - \frac{D_{12}^{2}}{D_{11}}
    \end{equation}
    \begin{equation}
        C_{xy} = 2 \left( D_{23} - D_{12} \frac{D_{13}}{D_{11}}\right),
    \end{equation}
\end{subequations}

where 
\begin{equation}
    D_{ij} = \left(\frac{A_{i1} A_{j1}}{c}\right)^{2} +  \left(\frac{A_{i2} A_{j2}}{b}\right)^{2} +  \left(\frac{A_{i3} A_{j3}}{a}\right)^{2}.
\end{equation}

The semi-major ($r_{a}$) and semi-minor ($r_{b}$) ellipse can then be expressed as
\begin{subequations}
\begin{equation}
r_{a} = \sqrt{\frac{2 \left( C_{yy} + C_{xx} + \sqrt{C_{xy}^{2} + (C_{yy}-C_{xx})^{2}}\right)}{4 C_{xx} C_{yy} - C_{xy}^{2}}}
\end{equation}
\begin{equation}
r_{b} = \sqrt{\frac{2 \left( C_{yy} + C_{xx} - \sqrt{C_{xy}^{2} + (C_{yy}-C_{xx})^{2}}\right)}{4 C_{xx} C_{yy} - C_{xy}^{2}}}.
\end{equation}
\end{subequations}
The projected ellipse aspect ratio $\varphi_{\mathrm{proj}}$ and angle deviation from horizontal $\theta_{\mathrm{proj}}$ can then be retrieved as

\begin{equation}
    \varphi_{\mathrm{proj}} = \frac{r_{b}}{r_{a}} = \sqrt{\frac{C_{yy} + C_{xx} - \sqrt{C_{xy}^{2} + (C_{yy}-C_{xx})^{2}}}{C_{yy} + C_{xx} + \sqrt{C_{xy}^{2} + (C_{yy}-C_{xx})^{2}}}},
\end{equation}

and
\begin{equation}
 \theta_{\mathrm{proj}} = 
\begin{cases}
   \arctan\left(\frac{C_{xy}}{C_{yy}-C_{xx} - \sqrt{(C_{xx}-C_{yy})^{2}+C_{xy}^{2}}}\right), & C_{xy} \neq 0\\
   90^{\circ}, & (C_{xy} =0) \land (C_{xx} < C_{yy})\\
   0^{\circ}, & (C_{xy} =0) \land (C_{xx} > C_{yy}).
\end{cases}
\end{equation}

\section{Plausible explanations for 2020 to 2022-2023 projection differences}\label{sect:discuss_King}

%\subsection{2D-S Simulations}

%\subsection{Plausible explanations for 2020 to 2022-2023 projection differences}\label{sect:discuss_King}

The stark difference in orientation angles between the 2020 and 2022-2023 campaigns and the striking similarity between the 2022 and 2023 campaigns suggest that there might have been differences in the 2D-S installation between the 2020 and 2022-2023 deployments. While it might seem possible that these differences are purely coincidental, due perhaps to different meteorological environments between the 2020 and 2022-2023 campaigns, there happens to be several pieces of strong evidence that these differences most likely resulted from the modifications to the aircraft and probe themselves. This is because in 2021, between the 2020 and 2022 IMPACTS campaign IOPs, the NASA GSFC Wallops Flight Facility (WFF) team found cracks in the P-3 starboard wing aileron (Mike Cropper WFF, personal communication). This prompted WFF engineers to add a fairing that was placed on the extended pylon from which the 2D-S was mounted. Figure~\ref{fig:pylon_comp} shows images of the P-3 wing before (panel c) and after (panel d) these modifications. As seen in Fig.~\ref{fig:pylon_comp}, this additional fairing (red dotted ellipse) covered the topmost portion of the extended pylon nearest to the wing. Computer-Aided Design (CAD) measurements (panels e and f in Fig.~\ref{fig:pylon_comp}) suggest that the total displacement of the 2D-S between the 2020 and 2022-2023 campaigns was only 3.4 inches as well as a 2.8$^{\circ}$ rotation downward. In this section, we explore plausible hypotheses for these discrepancies by utilizing these detailed CAD measurements of the P-3 aircraft and installed 2D-S locations for the 2020 and 2022-2023 configurations alongside Computational Fluid Dynamic (CFD) simulations of expected equilibrium orientation angles.

\FloatBarrier 
\begin{figure}[ht]
\centering
\includegraphics[width=\textwidth]{Pylon_faring_comp_2020_2022_w_coords.png}
\caption{P-3 aircraft coordinates in the X-Z (a) and X-Y (b) planes (note that this is a different basis than what is shown in Fig.~\ref{fig:P3}). Panels c and d show pictures of the P-3 station 18 pylons and probes for the IMPACTS 2020 (c) and IMPACTS 2022 (d) campaigns. The dotted red ellipse shows the additional fairing added in 2021 to the extended pylon. Panels e and f show the differences in 2D-S probe location between the 2020 campaign (blue) and the 2022-2023 campaigns (green) in the Y-Z (e) and X-Z (f) planes. The purple canister, probes and pylon shown in panel e represent the configuration of the original pylon before the IMPACTS 2020 campaign whereas the blue and green configurations represent the extended pylons used during the IMPACTS campaign for 2020 and 2022-2023 campaigns. Purple circles in panels a and b indicate the origin of the aircraft used in Fig.~\ref{fig:CFD_tests}. More information regarding the P-3 CAD measurements can be found in \citet{Hunter2025}.}
\label{fig:pylon_comp}
\end{figure}
\FloatBarrier

%Because of how slight these changes are, 

%It would seem that these slight changes shouldn't significantly alter the particle orientation behaviors. However, it is possible that these slight changes are meaningful due to the sensitivity of the equilibrium orientation angles with subtle displacements of the 2D-S. 

The NASA Langley TetrUSS computational fluid dynamics (CFD) system \citep{Frink2000} with the USM3D Mixed Element \citep[USM3D-ME;][]{Pandya2021} flow solver component was used to simulate the background fields of velocity, pressure, and density describing ambient airflow around the P-3 aircraft and the 2D-S volume. TetrUSS is a CFD software suite developed at NASA Langley over the last 35 years. It has become an aerodynamic tool of choice, with a wide user base inside and outside of NASA. Simulations were run as viscous, steady state, turbulent, Reynolds-averaged Navier-Stokes (RANS). The P-3’s four co-rotating propellers were simulated using a force-term source distribution in the Navier-Stokes equations. The computational domain stretched approximately 10 fuselage lengths from the P-3 aircraft to the far field in all directions. Preliminary grid resolution studies were conducted and showed acceptable convergence at a nominal grid size of 95 million cells. While numerous flight conditions were analyzed in support of the IMPACTS missions, the analysis here focused on high altitude (approximately 24,000 feet) and mid-altitude (approximately 10,000 feet) flight at cruise Mach numbers ranging from 0.44 to 0.53.  These conditions are representative of long-duration portions of the IMPACTS mission flights where 2D-S data were sampled. CFD solutions were run until iterative convergence was demonstrated over at least 1000 iterations, with steady forces and moments having under 2\% standard deviation. CFD solutions were run on the Pleiades supercomputer at NASA Ames Research Center, using 2000 cores and 1.5-2.0 hours of wall-clock time per case. Most solutions met convergence criteria in 4000 iterations using the HANIM acceleration methodology in USM3D-ME. Follow on computations with a separate 3-DOF solver were used to simulate two sets of 4,000 ice particle trajectories -- one for each 2D-S channel. Ice particles were modeled as point masses, using characteristics of a 1 mm rod with a drag coefficient of 0.8. With the assumption that the rod would always turn broadside to the local wind (i.e., the maximum drag direction), orientations were computed based on each particle’s three velocity components. As a check, initial 3-DOF orientation results were compared to the King theory \citep{King1985d}, and good agreement was seen in the far field and mid field away from the aircraft, before the effects of propeller swirl and local aerodynamic interference became strong (such effects are not included in the King theory).

%We simulated $4,000$ Lagrangian trajectories for both the H and V 2D-S channels using winds predicted by a 3 degree of freedom (3-DOF) CFD model to test the sensitivity of the slight 2D-S changes from 2020 to 2022-2023. Each trajectory was used to determine the hypothetical orientation of an assumed 1 mm rod with a drag coefficient assumed to be 0.8 and an angle of attack of 0.5 degrees. We assumed that each simulated rod continuously reorients itself to maximize its projected area with respect to the wind direction as in the column orientation model of \citet{King1985d}.

These Lagrangian trajectory results are shown in Fig.~\ref{fig:CFD_tests} alongside the ellipse fit angles from all 2020 and 2023 flights in the NGL. The 2020 configuration CFD H channel equilibrium mode occurs at roughly -70$^{\circ}$ whereas the V mode occurs at approximately -80$^{\circ}$. The IMPACTS H and V channel ellipse fit data for columns/needles however indicate modes at -70$^{\circ}$ and $+70^{\circ}$. Furthermore, CFD simulations strongly suggest that these equilibrium orientations would not substantially deviate from this preferred orientation. Similarly, the 2022-2023 CFD simulations in the H (V) channel peak at approximately $-80^{\circ}$ ($\pm 90^{\circ}$) compared to the $-70$ ($\pm 90^{\circ}$) from the IMPACTS ellipse data. Figure~\ref{fig:CFD_tests} also show the CFD trajectory depth positions within the H and V channels in order to determine orientation sensitivity to particle positions within the 2D-S H and V channel volumes. While there is some positional bias in the equilibrium orientations within each channel, this bias is overall insufficient to explain the observed ellipse fit orientation modes for either the 2020 or 2022-2023 configurations. Therefore, while there are deviations in the angle modes for both configurations, the 2022-2023 2D-S configuration appears to be much closer to the CFD simulations, particularly in the V channel. Pictures from the 2022 and 2023 campaigns (not shown) show physical ``Vertical'' and ``Horizontal'' labels that are on the correct 2D-S arms; such labels were not present during the 2020 campaign.  

%We set the origin of the P-3 aircraft approximately along the quarter-chord line of the P-3 aircraft as shown by the purple dots in \cref{fig:pylon_comp}a,b and we reorient the (x,y,z) coordinate system as described in \citet{King1985d}. \cref{fig:King_tests} shows the predicted $\theta_{\mathrm{proj},h}$ and $\theta_{\mathrm{proj},v}$ equilibrium angles for columns based on the 2D-S locations with respect to the P-3 origin. Dashed lines shown the 2D-S location in this coordinate system during the 2020 (black) and 2022-2023 (gray) campaigns. Overall, the predicted $\theta_{\mathrm{proj},v}\approx 70^{\circ}$ and $\theta_{\mathrm{proj},h}\approx -70^{\circ}$ for columns shown at warmer temperatures in \cref{fig:ar_angle_2020_temp} is well predicted by the \citet{King1985d} model. Furthermore, the approximate 3 in. displacement between the 2020 and 2022-2023 campaigns seems to have a minimal impact on the resulting equilibrium orientation angles. In particular, $\theta_{\mathrm{proj},h}$ would theoretically provide the biggest sensitivity to depth as indicated by the sharper gradients in panels a and b. $\theta_{\mathrm{proj},v}$ on the other hand is very insensitive to lateral displacements. The \citet{King1985d} model also explains why the observed $\theta_{\mathrm{proj},v}$ angles are much more concentrated at the predicted $70^{\circ}$ angle whereas the $\theta_{\mathrm{proj},h}$ has a much broader distribution. This can be seen in panels b and c where the tighter gradients in $\theta_{\mathrm{proj},h}$ suggests that there is a larger sensitivity to small changes in environmental conditions. The \citet{King1985d} model also explains why the observed $\theta_{\mathrm{proj},v}$ is slightly higher than previous studies such as \citet{Obrien2016} which found a slightly lower equilibrium angle of $60^{\circ}$. The IMPACTS campaign used an extended pylon which increased the distance the 2D-S was from the quarter-chord line. Therefore, while previous studies like \citet{Obrien2016} used a pylon like the purple one shown in panel e of \cref{fig:pylon_comp}, the extended pylon used during IMPACTS was farther separated from the aircraft wing which leads to a more negative orientation angle as shown in \cref{fig:King_tests}f. It is worth noting that the differences between the 2020 and 2022-2023 data can't be explained even if the true origin was different by a full meter in any direction. 

%\FloatBarrier 
%\begin{figure}[ht]
%\centering
%\includegraphics[width=\textwidth]{King_P3_orientation_tests_w_cbar.pdf}
%\caption{Output from the \citet{King1985d} orientation model. Panels b (c) shows the predicted equilibrium angles for $\theta_{h,\mathrm{proj}}$ ($\theta_{v,\mathrm{proj}}$) based on the $x$ distance  from the quarter-chord line and the $z$  depth ($y$ spanwise distance) from the aircraft origin at $x=600\, \mathrm{in.}$, $y=0\, \mathrm{in.}$ and $z=200\, \mathrm{in.}$ based on the coordinate system in panels a and b of \cref{fig:pylon_comp}. $\theta_{v,\mathrm{proj}}$ and $\theta_{h,\mathrm{proj}}$ represent angles in the 2D-S basis (see \cref{fig:P3}).}
%\label{fig:King_tests}
%\end{figure}
%\FloatBarrier

\FloatBarrier
\begin{figure}[ht]
\centering
\includegraphics[width=\textwidth]{IMPACTS_2DS_angle_CFD_comp_4k_freq_newcoords.pdf}
\caption{Frequency distribution comparison of equilibrium angles between CFD trajectory calculations and the IMPACTS ellipse fit data for all eight 2020 flights (top row) and all thirteen 2023 flights (middle row). (Bottom row) Scatterplot of the y position (left) and z position (right) from the center of the simulated 2D-S volume and corresponding equilibrium angle for all $4,000$ CFD trajectories for 2020 (blue) and 2023 (orange) configurations in the H and V channels. Only aspect ratios less than 0.4 in the NGL were included in order to better highlight the $\pm 70^{\circ}$ needle/column mode in the IMPACTS ellipse fit data.}
\label{fig:CFD_tests}
\end{figure}
\FloatBarrier

%It would seem that these modifications would not explain all the discrepancies shown in the particle orientation joint distributions from \cref{fig:ellipse_fit_joint2_2020} and \cref{fig:ellipse_fit_joint2_2022,fig:ellipse_fit_test}.

%An additional possibility is that the H and V channels were somehow flipped between the 2020 and 2022 deployments after the 2D-S was reinstalled. While this might seem improbable, it is noteworthy that the well-known $70^{\circ}$ preferred orientation of columns in the vertical plane is quite apparent in 2020 but not during the 2022 or 2023 campaigns (see \cref{fig:ellipse_fit_joint2_2020,fig:ellipse_fit_joint2_2022_2023}). 

%The fact that the $70^{\circ}$ columns appear in the vertical channel only for the 2020 cases strongly suggests that the 2020 2D-S was properly labeled. 

Interestingly, it appears that reversing the 2020 V channel angles produces an ellipse fit aspect ratio/orientation joint distribution that is shockingly similar to the 2023 H channel distribution. This is shown in Fig.~\ref{fig:rot_ex_discrep} for all 13 cases from the 2023 campaign and for all 8 cases from the 2020 campaign. This similarity can be explained if the 2D-S was incorrectly rotated by 90$^{\circ}$ clockwise in the probe coordinate system (see panel b in Fig.~\ref{fig:P3}). As shown in Fig.~\ref{fig:rot_ex_discrep} in this situation, the 2020 V channel plane would actual be in the H channel plane and the H channel plane would be in the V channel plane but with the V channel photodiode array rotated 180$^{\circ}$ with respect to the 2023 H channel (i.e., the 2020 labeled ``V'' photodiode array appearing on the port horizontal side rather than on the bottom 2D-S arm). This effectively would reverse the orientation angles if particle orientations were truly consistent among all 26 cases. While the above analysis fully explains the discrepancy in the labeled ``V'' channel for the 2020 flights, it is harder to reconcile the discrepancies between the 2020 labeled ``H'' channel (likely actually oriented in the 2022-2023 V channel location as shown in Fig.~\ref{fig:rot_ex_discrep}). 

%We hypothesize that the addition of the pylon fairing (as shown in \cref{fig:King_tests}) disturbed the airflow around the probe such that the Froude number went from supercritical to critical. 

%The Lagrangian trajectories in \cref{fig:CFD_tests} however suggests that 

%the free stream velocity would have to be quite low ($\sim 1.0$ m/s) in order for the equilibrium $\theta_{\mathrm{proj},h}$ angle for columns to be close to $0^{\circ}$. Even if the fairing blocked the free air velocity, it seems unlikely that the effect near the 2D-S probe would be disturbed to this degree. 

%Tests using the \citet{King1985d} model however suggests that the free stream velocity would have to be quite low ($\sim 1.0$ m/s) in order for the equilibrium $\theta_{\mathrm{proj},h}$ angle for columns to be close to $0^{\circ}$. Even if the fairing blocked the free air velocity, it seems unlikely that the effect near the 2D-S probe would be disturbed to this degree. 

%Instead, it is possible that the 2022-2023 IMPACTS cases had more instances of rimed particles which c

\FloatBarrier 
\begin{figure}[ht]
\centering
\includegraphics[width=\textwidth]{IMPACTS_2DS_AspectRatio_EllipseAngle_ROT_comp_newcoords.pdf}
\caption{(Top row) Original joint distributions of ellipse fit aspect ratio and angle for all 8 V-channel 2020 cases (left) and all 13 H-channel 2023 cases (right) where $T_{c}<0^{\circ}\mathrm{C}$. (Middle row) As in the top row but with the 2020 angles flipped. The bottom row shows how hypothetically this reversal could happen if the 2020 IMPACTS flights had the 2D-S probe rotated 90$^{\circ}$ counter-clockwise compared to the 2022 and 2023 flights.}
\label{fig:rot_ex_discrep}
\end{figure}
\FloatBarrier


\noappendix       %% use this to mark the end of the appendix section. Otherwise the figures might be numbered incorrectly (e.g. 10 instead of 1).

%% Regarding figures and tables in appendices, the following two options are possible depending on your general handling of figures and tables in the manuscript environment:

%% Option 1: If you sorted all figures and tables into the sections of the text, please also sort the appendix figures and appendix tables into the respective appendix sections.
%% They will be correctly named automatically.

%% Option 2: If you put all figures after the reference list, please insert appendix tables and figures after the normal tables and figures.
%% To rename them correctly to A1, A2, etc., please add the following commands in front of them:

\appendixfigures  %% needs to be added in front of appendix figures

\appendixtables   %% needs to be added in front of appendix tables

%% Please add \clearpage between each table and/or figure. Further guidelines on figures and tables can be found below.

\codedataavailability{The Matlab and Python code used to produce figures in the study are available on Github at \url{https://github.com/EdwinDunnavan/IMPACTS_AMT_Study} and is also available on Zenodo \citep{Dunnavan2026a}. The 2D-S ellipse fit NASA IMPACTS data is available from \citet{Dunnavan2026b} as a .mat file and the STL simulation data used in this paper is available from \citet{Dunnavan2026c}.} %% use this section when having data sets and software code available


\authorcontribution{ELD conceived of the study, developed and performed the simulations, calculated the statistics for the IMPACTS data and simulations, and wrote the manuscript. ELD processed the ellipse fit data using the UIOOPS software and JAF assisted in incorporating quality control measures for 2D-S particle images. DJD and CN collected the 2D-S probe data during the IMPACTS campaign and contributed to the analysis in section 5. AB provided SODA-2 image output using the simulations performed in this study. CH produced panels a and b in Fig.~\ref{fig:pylon_comp} as well as performed CFD-based trajectories used in Appendix B. All authors contributed to the discussion and improvement of the initial paper.} %% this section is mandatory

\competinginterests{The authors assert no conflicts of interest.} %% this section is mandatory even if you declare that no competing interests are present

%\disclaimer{TEXT} %% optional section

\begin{acknowledgements}\label{acknowledgements}
The authors would like to thank Mike Cropper (NASA Langley) and Joseph O'Brien (Argonne National Laboratory) for their help understanding the 2D-S probe and technical aspects of the 2D-S during the IMPACTS campaign, Julian Schima (Penn State) for providing the \citet{Schima2024} habit identification algorithm, Emma J{\"a}rvinen and Martin Schnaiter (University of Wuppertal) for their PHIPS images, Lee Thornhill (NASA Langley) for providing angle of attack data during IMPACTS, and Greg McFarquhar (CIWRO/OU) for his suggestions that improved this paper. The authors acknowledge the use of Google Gemini Pro as a tool for refactoring the Matlab and Python code used to generate the plots used in this paper. Funding for this work was provided by the NOAA/Office of Oceanic and Atmospheric Research under NOAA–University of Oklahoma Cooperative Agreement NA21OAR4320204, U.S. Department of Commerce.
\end{acknowledgements}




%% REFERENCES

%% The reference list is compiled as follows:

%\begin{thebibliography}{}

%\bibitem[AUTHOR(YEAR)]{LABEL1}
%REFERENCE 1

%\bibitem[AUTHOR(YEAR)]{LABEL2}
%REFERENCE 2

%\end{thebibliography}

%% Since the Copernicus LaTeX package includes the BibTeX style file copernicus.bst,
%% authors experienced with BibTeX only have to include the following two lines:
%%
\bibliographystyle{copernicus}
\bibliography{references.bib}
%%
%% URLs and DOIs can be entered in your BibTeX file as:
%%
%% URL = {http://www.xyz.org/~jones/idx_g.htm}
%% DOI = {10.5194/xyz}


%% LITERATURE CITATIONS
%%
%% command                        & example result
%% \citet{jones90}|               & Jones et al. (1990)
%% \citep{jones90}|               & (Jones et al., 1990)
%% \citep{jones90,jones93}|       & (Jones et al., 1990, 1993)
%% \citep[p.~32]{jones90}|        & (Jones et al., 1990, p.~32)
%% \citep[e.g.,][]{jones90}|      & (e.g., Jones et al., 1990)
%% \citep[e.g.,][p.~32]{jones90}| & (e.g., Jones et al., 1990, p.~32)
%% \citeauthor{jones90}|          & Jones et al.
%% \citeyear{jones90}|            & 1990



%% FIGURES

%% When figures and tables are placed at the end of the MS (article in one-column style), please add \clearpage
%% between bibliography and first table and/or figure as well as between each table and/or figure.

% The figure files should be labelled correctly with Arabic numerals (e.g. fig01.jpg, fig02.png).


%% ONE-COLUMN FIGURES

%%f
%\begin{figure}[t]
%\includegraphics[width=8.3cm]{FILE NAME}
%\caption{TEXT}
%\end{figure}
%
%%% TWO-COLUMN FIGURES
%
%%f
%\begin{figure*}[t]
%\includegraphics[width=12cm]{FILE NAME}
%\caption{TEXT}
%\end{figure*}
%
%
%%% TABLES
%%%
%%% The different columns must be seperated with a & command and should
%%% end with \\ to identify the column brake.
%
%%% ONE-COLUMN TABLE
%
%%t
%\begin{table}[t]
%\caption{TEXT}
%\begin{tabular}{column = lcr}
%\tophline
%
%\middlehline
%
%\bottomhline
%\end{tabular}
%\belowtable{} % Table Footnotes
%\end{table}
%
%%% TWO-COLUMN TABLE
%
%%t
%\begin{table*}[t]
%\caption{TEXT}
%\begin{tabular}{column = lcr}
%\tophline
%
%\middlehline
%
%\bottomhline
%\end{tabular}
%\belowtable{} % Table Footnotes
%\end{table*}
%
%%% LANDSCAPE TABLE
%
%%t
%\begin{sidewaystable*}[t]
%\caption{TEXT}
%\begin{tabular}{column = lcr}
%\tophline
%
%\middlehline
%
%\bottomhline
%\end{tabular}
%\belowtable{} % Table Footnotes
%\end{sidewaystable*}
%
%
%%% MATHEMATICAL EXPRESSIONS
%
%%% All papers typeset by Copernicus Publications follow the math typesetting regulations
%%% given by the IUPAC Green Book (IUPAC: Quantities, Units and Symbols in Physical Chemistry,
%%% 2nd Edn., Blackwell Science, available at: http://old.iupac.org/publications/books/gbook/green_book_2ed.pdf, 1993).
%%%
%%% Physical quantities/variables are typeset in italic font (t for time, T for Temperature)
%%% Indices which are not defined are typeset in italic font (x, y, z, a, b, c)
%%% Items/objects which are defined are typeset in roman font (Car A, Car B)
%%% Descriptions/specifications which are defined by itself are typeset in roman font (abs, rel, ref, tot, net, ice)
%%% Abbreviations from 2 letters are typeset in roman font (RH, LAI)
%%% Vectors are identified in bold italic font using \vec{x}
%%% Matrices are identified in bold roman font
%%% Multiplication signs are typeset using the LaTeX commands \times (for vector products, grids, and exponential notations) or \cdot
%%% The character * should not be applied as mutliplication sign
%
%
%%% EQUATIONS
%
%%% Single-row equation
%
%\begin{equation}
%
%\end{equation}
%
%%% Multiline equation
%
%\begin{align}
%& 3 + 5 = 8\\
%& 3 + 5 = 8\\
%& 3 + 5 = 8
%\end{align}
%
%
%%% MATRICES
%
%\begin{matrix}
%x & y & z\\
%x & y & z\\
%x & y & z\\
%\end{matrix}
%
%
%%% ALGORITHM
%
%\begin{algorithm}
%\caption{...}
%\label{a1}
%\begin{algorithmic}
%...
%\end{algorithmic}
%\end{algorithm}
%
%
%%% CHEMICAL FORMULAS AND REACTIONS
%
%%% For formulas embedded in the text, please use \chem{}
%
%%% The reaction environment creates labels including the letter R, i.e. (R1), (R2), etc.
%
%\begin{reaction}
%%% \rightarrow should be used for normal (one-way) chemical reactions
%%% \rightleftharpoons should be used for equilibria
%%% \leftrightarrow should be used for resonance structures
%\end{reaction}
%
%
%%% PHYSICAL UNITS
%%%
%%% Please use \unit{} and apply the exponential notation


\end{document}
